% Les amines — Chimie Tle D — Cours signé OnBuch+
% Programme officiel MINESEC, Terminale C, D, E. Compiler avec : tectonic L03-les-amines.tex
\newcommand{\DOCMATIERE}{Chimie}
\newcommand{\DOCNIVEAU}{Tle D}
\newcommand{\DOCMODULE}{Module 1 · Chimie organique}
\newcommand{\DOCLECON}{03}
\newcommand{\DOCTITRE}{Les amines}
% OnBuch+ — préambule « cours signés » ENRICHI (compilé avec tectonic / XeLaTeX)
% Système visuel « Pop » d'OnBuch+ : fond crème (jamais blanc), bordures encre
% épaisses, ombres « dures » décalées, titres Archivo Black en capitales.
% Version enrichie pour les cours de chimie : chimie (mhchem, chemfig),
% graphes (pgfplots), boîtes pédagogiques supplémentaires (définition,
% propriété, méthode, attention, le savais-tu, exercice résolu, exercice,
% TP, bilan), page de garde refaite et signature OnBuch+ ancrée en bas.
%
% Macros attendues dans main.tex AVANT \input{preamble} :
%   \DOCMATIERE  \DOCNIVEAU  \DOCMODULE  \DOCLECON (numéro)  \DOCTITRE
\documentclass[11pt]{article}

\usepackage{fontspec}
\usepackage[french]{babel}
\usepackage[a4paper,margin=2cm,top=3.4cm,bottom=2.8cm,headheight=1.4cm,footskip=1.3cm]{geometry}
\usepackage{amsmath,amssymb,amsfonts}
\usepackage[table]{xcolor} % [table] : fournit aussi \rowcolors (alternance de lignes)
\usepackage{pagecolor}
\usepackage{tikz}
\usetikzlibrary{arrows.meta,positioning,calc,shapes.geometric,shapes.misc,shapes.arrows,decorations.pathmorphing,decorations.markings,patterns,shadows,mindmap,trees,fit,backgrounds,matrix,intersections}
\usepackage{pgfplots}
\pgfplotsset{compat=1.18}
\usepackage[version=4]{mhchem}
\usepackage{chemfig}
\usepackage{enumitem}
\usepackage{fancyhdr}
% [most] charge toutes les bibliothèques tcolorbox (listings, minted, raster,
% documentation...) et gonfle inutilement la mémoire TeX sur un document long
% (~300 boîtes) : on ne charge que ce qui est réellement utilisé.
\usepackage[skins,breakable]{tcolorbox}
\usepackage{graphicx}
\usepackage{colortbl}
\usepackage{booktabs}
\usepackage{tabularx}
\usepackage{array}
\usepackage{multicol}
\usepackage{siunitx}
% parse-numbers=false : accepte aussi \SI{2,5 \times 10^{-3}}{...}
\sisetup{locale=FR,output-decimal-marker={,},group-minimum-digits=4,parse-numbers=false}
\usepackage{qrcode}
% \degree : commande fréquemment utilisée par les modèles pour un angle ou une
% température en dehors de \SI{}{} (non fournie par les paquets chargés).
\providecommand{\degree}{^{\circ}}
\usepackage{lastpage}
\usepackage{hyperref}
\hypersetup{hidelinks}

% ============================================================
% Palette « Pop » OnBuch+ — mêmes hex que la classe P (plus_ui.dart)
% ============================================================
\definecolor{poporange}{HTML}{F59321}
\definecolor{popdark}{HTML}{E07A0C}
\definecolor{popcream}{HTML}{FFF6E8}
\definecolor{popink}{HTML}{1C1714}
\definecolor{poppurple}{HTML}{7A5AE0}
\definecolor{popgreen}{HTML}{2E9E6A}
\definecolor{poppink}{HTML}{E0457B}
\definecolor{popblue}{HTML}{2F7DE1}
\definecolor{popgold}{HTML}{C98A12}
\definecolor{popmuted}{HTML}{8A7F76}
\definecolor{popline}{HTML}{EADFD2}
\definecolor{popgray}{HTML}{8A7F76}
\colorlet{popgrey}{popgray}
% Alias tolérants (le modèle nomme parfois les couleurs différemment)
\colorlet{popwhite}{white}
\colorlet{popblack}{popink}
\colorlet{popred}{poppink}
\colorlet{popyellow}{popgold}
% Teintes claires pour fonds de tableaux / zones
\colorlet{popblueL}{popblue!10!white}
\colorlet{poporangeL}{poporange!14!white}
\colorlet{popgreenL}{popgreen!12!white}
\colorlet{poppurpleL}{poppurple!12!white}
\colorlet{poppinkL}{poppink!12!white}
\colorlet{popgoldL}{popgold!14!white}

\pagecolor{popcream}

% ============================================================
% Typographie
% ============================================================
\defaultfontfeatures{Ligatures=TeX}
\setmainfont{PlusJakartaSans-Medium.ttf}[Path=../../../fonts/,BoldFont=PlusJakartaSans-Bold.ttf]
% Maths en sans-serif (Fira Math) avec chiffres et lettres droites de Plus
% Jakarta, pour que les formules se fondent dans le texte.
\usepackage[math-style=ISO,bold-style=ISO]{unicode-math}
\setmathfont{FiraMath-Regular.otf}[Path=../../../fonts/]
\setmathfont{PlusJakartaSans-Medium.ttf}[Path=../../../fonts/,range={up/num,up/{Latin,latin},\mathup}]
\usepackage{icomma} % virgule décimale sans espace en mode math
% Caractères absents de la police de texte (sinon carré vide) : rendus par
% leur équivalent mathématique, en texte comme en formule.
\usepackage{newunicodechar}
\newunicodechar{α}{\ensuremath{\alpha}}
\newunicodechar{Α}{\ensuremath{\alpha}}
\newunicodechar{β}{\ensuremath{\beta}}
\newunicodechar{δ}{\ensuremath{\delta}}
\newunicodechar{✓}{\popcheck{popgreen}}
\newfontfamily\popdisplay{ArchivoBlack-Regular.ttf}[Path=../../../fonts/]
\newfontfamily\poplogo{SpaceGrotesk-Bold.ttf}[Path=../../../fonts/]
\newfontfamily\popsemi{PlusJakartaSans-SemiBold.ttf}[Path=../../../fonts/]
\newfontfamily\popxbold{PlusJakartaSans-ExtraBold.ttf}[Path=../../../fonts/]
\newfontfamily\popxboldls{PlusJakartaSans-ExtraBold.ttf}[Path=../../../fonts/,LetterSpace=3.0]

\setlist[enumerate]{leftmargin=1.7em,itemsep=5pt,topsep=4pt}
\setlist[itemize]{leftmargin=1.7em,itemsep=4pt,topsep=4pt,label={\color{poporange}\textbullet}}
\setlength{\parindent}{0pt}
\setlength{\parskip}{5pt}
\renewcommand{\arraystretch}{1.25}

% chemfig : liaisons un peu plus épaisses, encre Pop
\setchemfig{atom sep=2em,bond style={line width=0.9pt,popink},cram width=3pt}

% pgfplots : style Pop par défaut
\pgfplotsset{
  every axis/.append style={axis line style={popink,line width=1pt},tick style={popink},
    grid style={popline,line width=0.5pt},label style={font=\small},tick label style={font=\footnotesize}},
  popcurve/.style={poporange,line width=1.8pt,no markers,smooth},
}

% ============================================================
% Pill / squiggle
% ============================================================
\newcommand{\poppill}[3]{%
  \tikz[baseline=-0.55ex]\node[draw=popink,line width=1.2pt,rounded corners=2.6pt,%
    fill=#2,inner xsep=6.5pt,inner ysep=3pt]{%
    \popxboldls\fontsize{7}{7}\selectfont\color{#3}\MakeUppercase{#1}};%
}
\newcommand{\popsquiggle}[1]{%
  \tikz[baseline=-0.4ex]{
    \draw[#1,line width=2.6pt,line cap=round]
      plot [smooth] coordinates {(0,0.09) (0.34,0.01) (0.68,0.21) (1.02,0.01) (1.36,0.10)};
  }%
}
% Mot-clé surligné (marqueur orange sous le texte)
\newcommand{\cle}[1]{{\popxbold\color{popdark}#1}}

% ============================================================
% En-tête / pied
% ============================================================
\pagestyle{fancy}
\fancyhf{}
\renewcommand{\headrulewidth}{0pt}
\renewcommand{\footrulewidth}{0pt}
\lhead{\raisebox{-0.15ex}{\poplogo\fontsize{13}{13}\selectfont\color{popink}OnBuch\color{poporange}+}}
\rhead{\poppill{\DOCMATIERE\ \textbullet\ \DOCNIVEAU\ \textbullet\ Leçon \DOCLECON}{white}{popink}}
\lfoot{\popxbold\fontsize{7.6}{7.6}\selectfont\color{popmuted}\MakeUppercase{Cours signé OnBuch+}\ \textbullet\ {\popsemi\DOCTITRE}}
\rfoot{\tikz[baseline=-0.6ex]\node[rounded corners=8.5pt,fill=popink,minimum height=17pt,minimum width=17pt,inner xsep=5pt,inner sep=0]{\popxbold\fontsize{8}{8}\selectfont\color{popcream}\thepage\,/\,\pageref*{LastPage}};}

% ============================================================
% Titres
% ============================================================
\newcommand{\chaptitre}[2]{%
  \begin{tcolorbox}[enhanced,colback=poporange,colframe=popink,boxrule=2.6pt,arc=11pt,
      drop shadow=popink,left=15pt,right=15pt,top=12pt,bottom=11pt,before skip=3pt,after skip=18pt]
    \poppill{#2}{popink}{popcream}\par\vspace{9pt}
    {\raggedright\hyphenpenalty=10000\exhyphenpenalty=10000\popdisplay\fontsize{22}{24}\selectfont\color{popcream}\MakeUppercase{#1}\par}
  \end{tcolorbox}
}
% Section numérotée : pastille orange avec le numéro + titre Archivo
\newcounter{popsec}
\newcommand{\coursec}[1]{%
  \stepcounter{popsec}\setcounter{popsub}{0}%
  \par\vspace{16pt}\noindent
  \begin{minipage}{\linewidth}
  \tikz[baseline=-0.7ex]\node[draw=popink,line width=1.6pt,fill=poporange,rounded corners=4pt,minimum size=22pt,inner sep=0]{\popdisplay\fontsize{12}{12}\selectfont\color{popcream}\thepopsec};%
  \hspace{8pt}{\popdisplay\fontsize{15}{17}\selectfont\color{popink}\MakeUppercase{#1}}\par
  \vspace{4pt}\noindent{\color{poporange}\rule{36pt}{3.6pt}}
  \end{minipage}\par\nopagebreak\vspace{8pt}
}
\newcounter{popsub}
\newcommand{\courssub}[1]{%
  \stepcounter{popsub}%
  \par\vspace{9pt}\noindent{\popxbold\fontsize{11.5}{13}\selectfont\color{popink}{\color{poporange}\thepopsec.\thepopsub}\hspace{6pt}#1}\par\nopagebreak\vspace{4pt}
}

% ============================================================
% Boîtes pédagogiques — même recette : fond blanc, bordure épaisse colorée,
% ombre dure encre, badge plein. Titre optionnel [#1].
% ============================================================
\tcbset{popbase/.style={breakable,enhanced,colback=white,boxrule=2.3pt,arc=9pt,
  shadow={3pt}{-3pt}{0pt}{popink},left=14pt,right=14pt,top=11pt,bottom=11pt,before skip=11pt,after skip=15pt,
  fontupper=\popsemi\fontsize{10.6}{14.5}\selectfont\color{popink},
  fontlower=\popsemi\fontsize{10.6}{14.5}\selectfont\color{popink}}}
% Coche dessinée (le glyphe ✓ n'existe pas dans les polices du document)
\newcommand{\popcheck}[1]{\tikz[baseline=-0.5ex]\draw[#1,line width=1.6pt,line cap=round,line join=round](-0.13,0.0)--(-0.04,-0.09)--(0.13,0.1);}
\newcommand{\popboxhead}[3]{\poppill{#1}{#2}{white}\ifx\relax#3\relax\else\hspace{6pt}{\popxbold\fontsize{10.6}{12}\selectfont\color{popink}#3}\fi\par\vspace{7pt}}

\newtcolorbox{aretenir}[1][]{popbase,colframe=popgreen,before upper={\popboxhead{À retenir}{popgreen}{#1}}}
\newtcolorbox{exemplebox}[1][]{popbase,colframe=poporange,before upper={\popboxhead{Exemple}{poporange}{#1}}}
\newtcolorbox{definition}[1][]{popbase,colframe=popblue,before upper={\popboxhead{Définition}{popblue}{#1}}}
\newtcolorbox{propriete}[1][]{popbase,colframe=poppurple,before upper={\popboxhead{Propriété}{poppurple}{#1}}}
\newtcolorbox{methode}[1][]{popbase,colframe=popgold,before upper={\popboxhead{Méthode}{popgold}{#1}}}
\newtcolorbox{attention}[1][]{popbase,colframe=poppink,before upper={\popboxhead{Attention}{poppink}{#1}}}
\newtcolorbox{experience}[1][]{popbase,colframe=popblue,colback=popblueL,before upper={\popboxhead{Expérience}{popblue}{#1}}}
% « Le savais-tu ? » : carte violette pleine (contexte, culture, Cameroun)
\newtcolorbox{savaistu}[1][]{popbase,colback=poppurple,colframe=popink,
  fontupper=\popsemi\fontsize{10.4}{14.2}\selectfont\color{white},
  before upper={\poppill{Le savais-tu\,?}{white}{poppurple}\ifx\relax#1\relax\else\hspace{6pt}{\popxbold\color{white}#1}\fi\par\vspace{7pt}}}
% Exercice résolu : énoncé en haut, \tcblower puis la solution sur fond crème
\newcounter{popexo}\newcounter{popexr}
\newtcolorbox{exoresolu}[1][]{popbase,colframe=poporange,colbacklower=popcream,
  segmentation style={popink,line width=1.2pt,dashed},
  before upper={\stepcounter{popexr}\popboxhead{Exercice résolu \thepopexr}{poporange}{#1}},
  before lower={\poppill{Solution}{popink}{popcream}\par\vspace{6pt}}}
% Exercice (non résolu en place) — la correction est dans la partie Corrigés
\newtcolorbox{exercice}[1][]{popbase,colframe=popink,
  before upper={\stepcounter{popexo}\popboxhead{Exercice \thepopexo}{popink}{#1}}}
% Corrigé
\newtcolorbox{corrige}[1][]{popbase,colframe=popgreen,colback=popgreenL,
  before upper={\popboxhead{Corrigé}{popgreen}{#1}}}
% Objectifs / prérequis / bilan
\newtcolorbox{objectifs}{popbase,colframe=popink,colback=poporangeL,
  before upper={\popboxhead{Objectifs de la leçon}{popink}{}}}
\newtcolorbox{prerequis}{popbase,colframe=popink,colback=popblueL,
  before upper={\popboxhead{Prérequis}{popblue}{}}}
\newtcolorbox{bilan}{popbase,colframe=popink,colback=white,boxrule=2.8pt,
  before upper={\popboxhead{Fiche bilan}{poporange}{L'essentiel à connaître pour l'examen}}}
% Figure encadrée avec légende
\newtcolorbox{popfigure}[1][]{enhanced,breakable=false,colback=white,colframe=popline,boxrule=1.4pt,arc=8pt,
  left=10pt,right=10pt,top=10pt,bottom=8pt,before skip=10pt,after skip=12pt,halign=center,
  lower separated=false,halign lower=center,
  fontlower=\popsemi\fontsize{9}{11.5}\selectfont\color{popmuted},
  #1}
\newcommand{\legende}[1]{\par\vspace{6pt}{\popsemi\fontsize{9}{11.5}\selectfont\color{popmuted}#1}}

% ============================================================
% Signature OnBuch+
% ============================================================
\newcommand{\poplogoinline}[1]{{\poplogo\fontsize{#1}{#1}\selectfont\color{popink}OnBuch\color{poporange}+}}
\newcommand{\signatureonbuch}{%
  \begin{tcolorbox}[enhanced,colback=popink,colframe=popink,boxrule=2.2pt,arc=10pt,
      drop shadow=poporange,left=14pt,right=14pt,top=10pt,bottom=10pt,before skip=0pt,after skip=0pt]
    \tikz[baseline=-0.7ex]\node[circle,fill=popgreen,draw=popcream,line width=1.3pt,minimum size=22pt,inner sep=0]{\popcheck{white}};%
    \hspace{9pt}\begin{minipage}[c]{0.62\linewidth}
      {\popxbold\fontsize{12}{13}\selectfont\color{popcream}Signé\ \textbullet\ {\poplogo OnBuch\color{poporange}+}}\par\vspace{2pt}
      {\raggedright\popsemi\fontsize{8}{10.5}\selectfont\color{popline}\MakeUppercase{Équipe pédagogique OnBuch+}\\\MakeUppercase{Conforme au programme officiel MINESEC}\par}
    \end{minipage}\hfill
    {\poplogo\fontsize{22}{22}\selectfont\color{popcream}OnBuch\color{poporange}+}
  \end{tcolorbox}%
}

% ============================================================
% Page de garde
% #1 titre · #2 matière · #3 niveau · #4 module · #5 n° leçon · #6 durée ·
% #7 description / objectifs (texte ou itemize)
% ============================================================
\newcommand{\pagegarde}[7]{%
  \thispagestyle{empty}
  \newgeometry{a4paper,margin=1.7cm,top=1.6cm,bottom=1.4cm}%
  \begin{tikzpicture}[remember picture,overlay]
    \fill[poporange!18] ([xshift=-3.2cm,yshift=-3.6cm]current page.north east) circle (4.6cm);
    \fill[poppurple!14] ([xshift=2.4cm,yshift=7.6cm]current page.south west) circle (3.8cm);
  \end{tikzpicture}%
  {\poplogo\fontsize{30}{30}\selectfont\color{popink}OnBuch\color{poporange}+}\hfill
  \raisebox{0.6ex}{\poppill{Cours signé \textbullet\ Premium}{popink}{popcream}}\par
  \vspace{1pt}\popsquiggle{poppurple}\par
  \vspace{8pt}
  \begin{tcolorbox}[enhanced,colback=poporange,colframe=popink,boxrule=2.8pt,arc=13pt,
      drop shadow=popink,left=18pt,right=18pt,top=12pt,bottom=13pt,after skip=10pt]
    \poppill{#2 \textbullet\ #3}{popink}{popcream}\hspace{5pt}\poppill{#4}{white}{popink}\par\vspace{8pt}
    {\popdisplay\fontsize{14}{14}\selectfont\color{popink}LEÇON #5}\par\vspace{4pt}
    {\raggedright\hyphenpenalty=10000\exhyphenpenalty=10000\popdisplay\fontsize{25}{27}\selectfont\color{popcream}\MakeUppercase{#1}\par}
    \vspace{8pt}
    {\popxbold\fontsize{9.5}{9.5}\selectfont\color{popink}\MakeUppercase{Durée conseillée : #6}}
  \end{tcolorbox}
  \begin{tcolorbox}[enhanced,colback=white,colframe=popink,boxrule=2.2pt,arc=10pt,
      drop shadow=popink,left=16pt,right=16pt,top=10pt,bottom=10pt,after skip=8pt]
    \poppill{Dans cette leçon}{poppurple}{white}\par\vspace{5pt}
    \popsemi\fontsize{9.3}{12.6}\selectfont\color{popink}#7
  \end{tcolorbox}
  \noindent\begin{minipage}[t]{0.487\linewidth}
    \begin{tcolorbox}[enhanced,colback=popgreen,colframe=popink,boxrule=2.2pt,arc=10pt,
        drop shadow=popink,left=11pt,right=11pt,top=10pt,bottom=10pt,height=3.1cm,valign=center]
      \tcbox[colback=white,colframe=popink,boxrule=1.3pt,arc=5pt,boxsep=0pt,nobeforeafter,
        left=3pt,right=3pt,top=3pt,bottom=3pt,valign=center]{\qrcode[height=1.9cm]{https://play.google.com/store/apps/details?id=com.luvvix.onbuch&hl=fr}}%
      \hspace{8pt}\begin{minipage}[c]{0.5\linewidth}
        {\popdisplay\fontsize{11}{12.5}\selectfont\color{white}\MakeUppercase{Télécharge OnBuch}}\par\vspace{4pt}
        {\raggedright\popsemi\fontsize{8.4}{11}\selectfont\color{white}Gratuit \textbullet\ Google Play\\Tuteur IA, annales, concours, résultats\par}
      \end{minipage}
    \end{tcolorbox}
  \end{minipage}\hfill
  \begin{minipage}[t]{0.487\linewidth}
    \begin{tcolorbox}[enhanced,colback=poppurple,colframe=popink,boxrule=2.2pt,arc=10pt,
        drop shadow=popink,left=11pt,right=11pt,top=10pt,bottom=10pt,height=3.1cm,valign=center]
      \tikz[baseline=-0.7ex]\node[circle,fill=white,draw=popink,line width=1.3pt,minimum size=1.6cm,inner sep=0]{\popdisplay\fontsize{24}{24}\selectfont\color{poppurple}+};%
      \hspace{8pt}\begin{minipage}[c]{0.6\linewidth}
        {\popdisplay\fontsize{11}{12.5}\selectfont\color{white}\MakeUppercase{Passe à OnBuch+}}\par\vspace{4pt}
        {\raggedright\popsemi\fontsize{8.4}{11}\selectfont\color{white}Tous les cours signés, Tuteur IA illimité, fiches PDF — onglet OnBuch+ de l'app\par}
      \end{minipage}
    \end{tcolorbox}
  \end{minipage}
  \vspace{8pt}
  \popcontenu
  \vfill
  \signatureonbuch
  \restoregeometry
  \newpage
}

% Rangée « Ce cours contient » (page de garde)
\newcommand{\poptile}[3]{% #1 couleur #2 gros texte #3 légende
  \begin{tcolorbox}[enhanced,colback=#1,colframe=popink,boxrule=2pt,arc=9pt,shadow={2.5pt}{-2.5pt}{0pt}{popink},
      left=6pt,right=6pt,top=9pt,bottom=9pt,halign=center,valign=center,height=1.75cm,width=\linewidth,nobeforeafter]
    {\popdisplay\fontsize{12.5}{13}\selectfont\color{white}\MakeUppercase{#2}}\par\vspace{3pt}
    {\centering\popsemi\fontsize{7.8}{9.5}\selectfont\color{white}#3\par}
  \end{tcolorbox}}
\newcommand{\popcontenu}{%
  \par\noindent{\popxboldls\fontsize{7.5}{7.5}\selectfont\color{popmuted}CE COURS SIGNÉ CONTIENT}\par\vspace{7pt}
  \noindent\begin{minipage}[t]{0.235\linewidth}\poptile{popblue}{Cours}{complet, illustré, pas à pas}\end{minipage}\hfill
  \begin{minipage}[t]{0.235\linewidth}\poptile{poporange}{Méthodes}{et exercices résolus}\end{minipage}\hfill
  \begin{minipage}[t]{0.235\linewidth}\poptile{poppink}{Exercices}{type examen + corrigés}\end{minipage}\hfill
  \begin{minipage}[t]{0.235\linewidth}\poptile{popgreen}{Fiche bilan}{l'essentiel à retenir}\end{minipage}\par}

% Sommaire
\newtcolorbox{sommaire}{enhanced,colback=white,colframe=popink,boxrule=2.2pt,arc=10pt,
  drop shadow=popink,left=18pt,right=18pt,top=13pt,bottom=13pt,before skip=6pt,after skip=20pt,
  before upper={\poppill{Sommaire}{poppurple}{white}\par\vspace{9pt}\popsemi\fontsize{10.6}{14.5}\selectfont\color{popink}}}

% Signature de fin de cours — poussée tout en bas de la dernière page
\newcommand{\findecours}{%
  \par\vfill\nopagebreak
  \begin{center}\popsquiggle{poporange}\end{center}\vspace{4pt}
  \signatureonbuch
}

%% FIGFIX-BEGIN v5 — correctifs globaux des figures (docs/onbuch-generation/tools/figfix)
\usepackage{adjustbox}\usepackage{etoolbox}\tcbuselibrary{raster}
% toute figure TikZ/pgfplots plus large que la ligne est réduite à la largeur disponible
\BeforeBeginEnvironment{tikzpicture}{\begin{adjustbox}{max width=\linewidth}}
\AfterEndEnvironment{tikzpicture}{\end{adjustbox}}
% légendes pgfplots lisibles par-dessus les courbes
\pgfplotsset{every axis/.append style={legend style={fill=white,fill opacity=0.92,text opacity=1,draw=black!15,inner sep=3pt}}}
% étiquettes d'arêtes (syntaxe edge["..."]) avec fond blanc
\tikzset{every edge quotes/.append style={fill=white,inner sep=1.2pt}}
%% FIGFIX-END
\begin{document}

\pagegarde{Les amines}{Chimie}{Tle D}{Module 1 · Chimie organique}{03}{3 h}{Cette leçon introduit les amines, dérivés azotés de l'ammoniac : structure, classes, nomenclature, propriétés physiques et chimiques (basicité, nucléophilie). Elle montre comment les amines permettent de préparer les ammoniums quaternaires (réaction d'Hofmann) et les amides, briques essentielles des médicaments.
\vspace{4pt}
\begin{multicols}{2}\small
\begin{itemize}
\item À la fin de cette leçon, je sais décrire la structure pyramidale de l'ammoniac et justifier le passage aux amines par substitution des atomes d'hydrogène
\item À la fin de cette leçon, je sais écrire la formule générale des amines et identifier les trois classes (primaire, secondaire, tertiaire) ainsi que l'ion ammonium quaternaire
\item À la fin de cette leçon, je sais nommer une amine à partir de sa formule semi-développée et inversement (y compris l'aniline)
\item À la fin de cette leçon, je sais expliquer les propriétés physiques des amines (état, solubilité, odeur) par les liaisons hydrogène
\item À la fin de cette leçon, je sais écrire l'équation-bilan de l'action de l'eau sur une amine et comparer sa basicité à celle de l'ammoniac
\item À la fin de cette leçon, je sais mettre en évidence expérimentalement le caractère basique des amines (indicateurs colorés, précipitation des ions Cu2+)
\end{itemize}\end{multicols}}

\begin{sommaire}
\begin{multicols}{2}
\begin{enumerate}
\item De l'ammoniac aux amines : structure, définition et classes
\item Nomenclature des amines
\item Propriétés physiques des amines
\item Caractère basique des amines
\item Caractère nucléophile : alkylation et réactions d'Hofmann
\item Action sur les chlorures d'acyle : préparation des amides
\item Travaux pratiques
\item Méthodes et exercices résolus
\item Exercices
\item Corrigés des exercices
\item Fiche bilan
\end{enumerate}
\end{multicols}
\end{sommaire}

\begin{objectifs}
\begin{itemize}
\item À la fin de cette leçon, je sais décrire la structure pyramidale de l'ammoniac et justifier le passage aux amines par substitution des atomes d'hydrogène
\item À la fin de cette leçon, je sais écrire la formule générale des amines et identifier les trois classes (primaire, secondaire, tertiaire) ainsi que l'ion ammonium quaternaire
\item À la fin de cette leçon, je sais nommer une amine à partir de sa formule semi-développée et inversement (y compris l'aniline)
\item À la fin de cette leçon, je sais expliquer les propriétés physiques des amines (état, solubilité, odeur) par les liaisons hydrogène
\item À la fin de cette leçon, je sais écrire l'équation-bilan de l'action de l'eau sur une amine et comparer sa basicité à celle de l'ammoniac
\item À la fin de cette leçon, je sais mettre en évidence expérimentalement le caractère basique des amines (indicateurs colorés, précipitation des ions Cu2+)
\item À la fin de cette leçon, je sais écrire les équations-bilans des réactions d'Hofmann successives d'une amine sur un dérivé halogéné jusqu'à l'ion ammonium quaternaire
\item À la fin de cette leçon, je sais écrire l'équation-bilan de l'action d'une amine sur un chlorure d'acyle (ou un anhydride) conduisant à un amide
\end{itemize}
\end{objectifs}

\begin{prerequis}
\begin{itemize}
\item Structure de Lewis et géométrie des molécules (méthode VSEPR, notion de doublet non liant)
\item Nomenclature des alcanes et des dérivés halogénés (Première)
\item Notion d'acide et de base de Brønsted, couples acide/base, pH (Première)
\item Dérivés halogénés R-X : polarité de la liaison C-X et caractère électrophile du carbone
\item Dérivés acides : acides carboxyliques, chlorures d'acyle et anhydrides (leçon précédente)
\end{itemize}
\end{prerequis}

\newpage

\coursec{De l'ammoniac aux amines : structure, définition et classes}

\textbf{Situation d'accroche.} À Douala, une pharmacienne prépare un sirop antipaludéen à base de chloroquine : cette molécule contient des atomes d'azote issus de la famille des \emph{amines}. Quelques kilomètres plus loin, au marché de Mokolo, l'odeur âcre du poisson qui se gâte est elle aussi due à des amines volatiles. Comment passe-t-on de la molécule d'ammoniac, simple et pyramidale, à ces composés azotés si présents dans les médicaments, les colorants et les odeurs de la vie quotidienne ? Et comment les chimistes fabriquent-ils, à partir d'amines, les amides qui constituent les protéines et de nombreux médicaments ? C'est tout l'objet de cette leçon.

\textbf{Problématique.} Quelle est la structure des amines, comment les classe-t-on, et pourquoi le petit doublet d'électrons porté par l'azote commande-t-il toute leur chimie ?

\courssub{Rappel : la structure de la molécule d'ammoniac \ce{NH3}}

\textbf{Intuition.} Avant de comprendre les amines, il faut revenir à leur « mère » : la molécule d'\cle{ammoniac} \ce{NH3}, bien connue des agriculteurs (engrais azotés) et des produits d'entretien. Tout ce que feront les amines dans cette leçon — se comporter en base, attaquer des molécules — est déjà écrit dans la structure de \ce{NH3}.

L'atome d'azote possède \num{5} électrons sur sa couche externe (configuration $1s^2\,2s^2\,2p^3$). Pour satisfaire la règle de l'octet, il forme \num{3} liaisons covalentes avec \num{3} atomes d'hydrogène, et il lui reste \num{2} électrons non engagés : un \cle{doublet non liant} (ou doublet libre).

\begin{popfigure}
\begin{center}
\chemfig{N(-[2]H)(-[6]H)(-[0]H)}
\hspace{2cm}
\begin{tikzpicture}[scale=1.1]
  % Azote au sommet
  \shade[ball color=popblue] (0,1.4) circle (0.32) node[white,font=\small\bfseries]{N};
  % Doublet non liant
  \fill[poporange] (-0.10,1.95) circle (0.05);
  \fill[poporange] (0.10,1.95) circle (0.05);
  \node[poporange,font=\footnotesize] at (0,2.35) {doublet non liant};
  % Hydrogènes à la base
  \shade[ball color=popmuted!60] (-1.1,-0.4) circle (0.22) node[white,font=\scriptsize\bfseries]{H};
  \shade[ball color=popmuted!60] (1.1,-0.4) circle (0.22) node[white,font=\scriptsize\bfseries]{H};
  \shade[ball color=popmuted!60] (0,-0.75) circle (0.22) node[white,font=\scriptsize\bfseries]{H};
  % Liaisons
  \draw[line width=2pt,popdark] (0,1.15) -- (-1.0,-0.22);
  \draw[line width=2pt,popdark] (0,1.15) -- (1.0,-0.22);
  \draw[line width=2pt,popdark] (0,1.12) -- (0,-0.52);
  % Angle
  \draw[popgreen,thick,->] (-0.55,0.35) arc[start angle=200,end angle=340,radius=0.65];
  \node[popgreen,font=\footnotesize] at (0,0.05) {$\approx \SI{107}{\degree}$};
  % Base en pointillés
  \draw[dashed,popmuted] (-1.1,-0.4) -- (1.1,-0.4) -- (0,-0.75) -- cycle;
\end{tikzpicture}
\end{center}
\legende{À gauche : représentation de Lewis simplifiée de \ce{NH3}. À droite : géométrie pyramidale à base triangulaire ; l'angle $\widehat{\mathrm{HNH}}$ vaut environ \SI{107}{\degree}. Le doublet non liant (points orange) occupe le sommet du tétraèdre virtuel.}
\end{popfigure}

\begin{propriete}[Géométrie de l'ammoniac]
La molécule d'ammoniac \ce{NH3} a une géométrie \cle{pyramidale} à base triangulaire : l'atome d'azote occupe le sommet, les trois atomes d'hydrogène la base. Les quatre doublets d'électrons autour de l'azote (trois liants, un non liant) se disposent approximativement en tétraèdre ; le doublet non liant, plus « répulsif », comprime les liaisons \ce{N-H} : l'angle $\widehat{\mathrm{HNH}}$ vaut environ \SI{107}{\degree} au lieu des \SI{109,5}{\degree} du tétraèdre parfait.
\end{propriete}

\textbf{Le doublet non liant : le fil conducteur de toute la leçon.} Ce doublet d'électrons « disponible » est la clé de voûte de la chimie des amines :

\begin{itemize}
  \item il peut \textbf{capter un proton} \ce{H+} : c'est le \cle{caractère basique} (section 4) ;
  \item il peut \textbf{attaquer un atome de carbone porteur d'une charge partielle positive} $\delta^+$ : c'est le \cle{caractère nucléophile} (sections 5 et 6).
\end{itemize}

\begin{aretenir}[L'idée-force]
Toute la chimie des amines s'explique par une seule phrase : \emph{l'atome d'azote porte un doublet non liant}. Retiens cette image : le doublet est une « main » tendue, prête à saisir un proton (basicité) ou un carbone électrophile (nucléophilie).
\end{aretenir}

\courssub{Définition des amines}

\textbf{Intuition.} Imagine la molécule d'ammoniac comme une « tête » d'azote coiffée de trois « chapeaux » hydrogène. Remplace un, deux ou trois de ces chapeaux par des groupes carbonés (alkyles ou aryles) : tu obtiens les amines. L'azote garde toujours son doublet non liant — donc son caractère basique et nucléophile.

\begin{definition}[Amine]
Une \cle{amine} est un composé organique \emph{dérivé de l'ammoniac} \ce{NH3} par remplacement de un, deux ou trois atomes d'hydrogène par des groupes carbonés : groupes \cle{alkyles} $\mathrm{R}$ (chaînes carbonées saturées, ex. \ce{CH3-}, \ce{C2H5-}) ou groupes \cle{aryles} $\mathrm{Ar}$ (cycles aromatiques, ex. phényle \ce{C6H5-}).
\end{definition}

Le groupe caractéristique des amines est le groupe \cle{amino} \ce{-NH2} (lorsque l'azote ne porte qu'un seul groupe carboné). L'atome d'azote y conserve toujours son doublet non liant, et la géométrie autour de l'azote reste pyramidale, comme dans \ce{NH3}.

\begin{exemplebox}[Des amines de la vie quotidienne]
\begin{itemize}
  \item La \textbf{méthylamine} \chemfig{CH_3-NH_2} : un seul hydrogène de \ce{NH3} remplacé par un groupe méthyle ;
  \item l'\textbf{éthylamine} \chemfig{CH_3-CH_2-NH_2} : présente dans certaines décompositions de protéines ;
  \item la \textbf{triméthylamine} \chemfig{N(-[2]CH_3)(-[6]CH_3)-CH_3} : responsable de l'odeur du poisson avarié ;
  \item l'\textbf{aniline} (phénylamine) \chemfig{C_6H_5-NH_2} : un groupe aryle ; matière première historique des colorants.
\end{itemize}
\end{exemplebox}

\courssub{Les trois classes d'amines}

\textbf{Intuition.} Combien de « chapeaux » hydrogène ont été remplacés ? C'est la seule question à poser pour classer une amine. On compte le nombre de groupes carbonés directement liés à l'atome d'azote.

\begin{definition}[Classe d'une amine]
La \cle{classe} d'une amine est le \textbf{nombre de groupes carbonés liés à l'atome d'azote}, c'est-à-dire le nombre de liaisons \ce{N-C} de la molécule :
\begin{itemize}
  \item \num{1} groupe carboné : amine \cle{primaire}, formule générale $R-\ce{NH2}$ ;
  \item \num{2} groupes carbonés : amine \cle{secondaire}, formule générale $R^1-\ce{NH}-R^2$ ;
  \item \num{3} groupes carbonés : amine \cle{tertiaire}, formule générale $R^1-\ce{N}(R^2)-R^3$.
\end{itemize}
Les groupes $\mathrm{R}$, $\mathrm{R'}$, $\mathrm{R''}$ peuvent être identiques ou différents, alkyles ou aryles.
\end{definition}

\begin{center}
\begin{tabularx}{\linewidth}{|l|X|X|X|}
\hline
\rowcolor{popblue!40} \textbf{Classe} & \textbf{Formule générale} & \textbf{Liaisons $N-C$} & \textbf{Exemple} \\
\hline
Primaire & $R-\ce{NH2}$ & \num{1} & Éthylamine \ce{C2H5NH2} \\
\hline
Secondaire & $R-\ce{NH}-R'$ & \num{2} & Diéthylamine \ce{(C2H5)2NH} \\
\hline
Tertiaire & $R-\ce{N}(R')-R''$ & \num{3} & Triméthylamine \ce{(CH3)3N} \\
\hline
\end{tabularx}
\end{center}

\begin{methode}[Classer une amine en trois étapes]
\begin{enumerate}
  \item Repérer l'atome d'\textbf{azote} dans la formule semi-développée ;
  \item \textbf{Compter les liaisons \ce{N-C}} partant de cet azote (ignorer les liaisons \ce{N-H}) ;
  \item Conclure : \num{1} liaison \ce{N-C} $\Rightarrow$ primaire ; \num{2} $\Rightarrow$ secondaire ; \num{3} $\Rightarrow$ tertiaire.
\end{enumerate}
\end{methode}

\begin{attention}[Le piège classique du Baccalauréat]
La classe d'une amine se détermine sur l'atome d'\textbf{azote}, \emph{jamais} sur le carbone ! Ne confonds pas avec la classe des alcools, qui porte sur le carbone fonctionnel.

Exemple piège : la \emph{tert}-butylamine \chemfig{CH_3-C(-[2]CH_3)(-[6]CH_3)-NH_2}, c'est-à-dire \chemfig{(CH_3)_3C-NH_2}. Le carbone lié à \ce{NH2} est bien un carbone \emph{tertiaire} (lié à trois autres carbones)... mais l'azote n'est lié qu'à \textbf{un seul} groupe carboné : c'est une amine \textbf{primaire}. Compter les liaisons \ce{N-C}, rien d'autre !
\end{attention}

\begin{exoresolu}[Classer trois amines]
On considère les trois composés suivants :
\[
\text{(a) } \chemfig{CH_3-CH_2-NH-CH_3} \qquad
\text{(b) } \chemfig{(CH_3)_3C-NH_2} \qquad
\text{(c) } \chemfig{C_2H_5-N(CH_3)-C_2H_5}
\]
Déterminer la classe de chacun.
\tcblower
\textbf{Solution détaillée.} On applique la méthode : compter les liaisons \ce{N-C}.

\textbf{(a)} \chemfig{CH_3-CH_2-NH-CH_3} : l'azote est lié à un groupe éthyle et à un groupe méthyle, soit \num{2} liaisons \ce{N-C} (et une liaison \ce{N-H}). C'est une amine \textbf{secondaire}.

\textbf{(b)} \chemfig{(CH_3)_3C-NH_2} : l'azote n'est lié qu'à \num{1} groupe carboné (le groupe \emph{tert}-butyle), même si ce groupe est ramifié. C'est une amine \textbf{primaire}. Le carbone tertiaire ne change rien à l'affaire.

\textbf{(c)} \chemfig{C_2H_5-N(CH_3)-C_2H_5} : l'azote est lié à deux groupes éthyles et un groupe méthyle, soit \num{3} liaisons \ce{N-C} (aucune liaison \ce{N-H}). C'est une amine \textbf{tertiaire}.
\end{exoresolu}

\courssub{De \ce{NH3} à l'ion ammonium quaternaire : l'arbre de substitution}

On peut visualiser le passage de l'ammoniac aux amines comme un arbre de substitutions successives des atomes d'hydrogène par des groupes alkyles $\mathrm{R}$ :

\begin{popfigure}
\begin{center}
\begin{tikzpicture}[
  node distance=1.1cm and 1.6cm,
  boite/.style={draw=popblue, thick, rounded corners=4pt, fill=popblueL, inner sep=6pt, align=center},
  fleche/.style={-{Stealth[length=3mm]}, thick, popdark}]
  \node[boite] (nh3) {\chemfig{N(-[2]H)(-[6]H)(-[0]H)}\\[2pt]\footnotesize ammoniac};
  \node[boite, right=of nh3] (prim) {\chemfig{N(-[2]H)(-[6]H)(-[0]R)}\\[2pt]\footnotesize amine primaire};
  \node[boite, right=of prim] (sec) {\chemfig{N(-[2]H)(-[6]R)(-[0]R')}\\[2pt]\footnotesize amine secondaire};
  \node[boite, right=of sec] (tert) {\chemfig{N(-[2]R'')(-[6]R)(-[0]R')}\\[2pt]\footnotesize amine tertiaire};
  \node[boite, below=of tert, fill=poporangeL, draw=poporange] (quat) {\footnotesize ion ammonium quaternaire\\ $\mathrm{R_4N^+}$ \; associé à un anion \ce{X-}};
  \draw[fleche] (nh3) -- node[above,font=\scriptsize]{$+\mathrm{R}$,\; $-\mathrm{H}$} (prim);
  \draw[fleche] (prim) -- node[above,font=\scriptsize]{$+\mathrm{R'}$,\; $-\mathrm{H}$} (sec);
  \draw[fleche] (sec) -- node[above,font=\scriptsize]{$+\mathrm{R''}$,\; $-\mathrm{H}$} (tert);
  \draw[fleche] (tert) -- node[right,font=\scriptsize]{$+\mathrm{R'''}$} (quat);
\end{tikzpicture}
\end{center}
\legende{Arbre de substitution : chaque remplacement d'un hydrogène de \ce{NH3} par un groupe carboné fait monter d'une classe. Une quatrième substitution conduit à l'ion ammonium quaternaire.}
\end{popfigure}

\begin{definition}[Ion ammonium quaternaire]
Si l'atome d'azote est lié à \textbf{quatre} groupes carbonés, il ne peut plus porter de doublet non liant : on obtient un cation, l'\cle{ion ammonium quaternaire}, de formule générale $\mathrm{R_4N^+}$. Sa charge positive est compensée par un anion (par exemple \ce{Cl-}, \ce{Br-}, \ce{OH-}) : on parle alors de \emph{sel d'ammonium quaternaire}.

Exemple : le \textbf{chlorure de tétraméthylammonium} \ce{(CH3)4N+ Cl-}, solide ionique, soluble dans l'eau.
\end{definition}

\begin{attention}[Ammonium quaternaire $\neq$ amine]
L'ion $\mathrm{R_4N^+}$ n'est \textbf{pas} une amine : il n'a plus de doublet non liant, donc plus de caractère basique ni nucléophile. C'est un cation stable, analogue à l'ion ammonium \ce{NH4+} où les quatre hydrogènes ont été remplacés par des groupes alkyles. Ne l'appelle jamais « amine quaternaire » dans une copie : dis « ion ammonium quaternaire » ou « sel d'ammonium quaternaire ».
\end{attention}

\begin{savaistu}[L'odeur du poisson avarié]
L'odeur caractéristique du poisson qui se gâte — celle qui monte parfois du marché de Mokolo en fin de journée — est due à la \textbf{triméthylamine} \chemfig{N(CH_3)_3}, une amine tertiaire volatile (elle bout vers \SI{3}{\celsius} seulement !). Elle est produite par la décomposition bactérienne des protéines et de l'oxyde de triméthylamine présent dans les tissus du poisson. Astuce de grand-mère validée par la chimie : arroser le poisson de jus de citron. L'acide citrique transforme la triméthylamine en ion triméthylammonium \chemfig{(CH_3)_3NH+}, non volatil... et l'odeur disparaît ! Tu retrouveras ce principe (base + acide $\rightarrow$ sel) dans la section 4.
\end{savaistu}

\begin{aretenir}[L'essentiel de la section]
\begin{itemize}
  \item \ce{NH3} est \textbf{pyramidale} (angle $\approx \SI{107}{\degree}$) et porte un \textbf{doublet non liant} : source de la basicité et de la nucléophilie des amines.
  \item Une \textbf{amine} dérive de \ce{NH3} par remplacement de 1, 2 ou 3 hydrogènes par des groupes carbonés : $R-\ce{NH2}$ (primaire), $R-\ce{NH}-R'$ (secondaire), $R-\ce{N}(R')-R''$ (tertiaire).
  \item La \textbf{classe} se détermine en comptant les liaisons \ce{N-C} sur l'\textbf{azote} — jamais sur le carbone.
  \item Quatre groupes carbonés sur l'azote donnent l'\textbf{ion ammonium quaternaire} $\mathrm{R_4N^+}$, associé à un anion.
\end{itemize}
\end{aretenir}

\begin{exercice}[À toi de jouer]
\begin{enumerate}
  \item Écrire les formules semi-développées de toutes les amines de formule brute \ce{C3H9N} et préciser la classe de chacune.
  \item La phénylamine (aniline) \chemfig{C_6H_5-NH_2} est-elle une amine primaire, secondaire ou tertiaire ? Justifier.
  \item On fait réagir de la triméthylamine avec de l'iodométhane \ce{CH3I}. Quelle espèce ionique obtient-on ? Écrire sa formule.
\end{enumerate}
\end{exercice}

\begin{corrige}[Exercice « À toi de jouer »]
\begin{enumerate}
  \item Il existe \textbf{quatre} isomères de constitution pour \ce{C3H9N} :
  \begin{itemize}
    \item \chemfig{CH_3-CH_2-CH_2-NH_2} : propan-1-amine (propylamine), \textbf{primaire} (1 liaison \ce{N-C}) ;
    \item \chemfig{CH_3-CH(NH_2)-CH_3} : propan-2-amine (isopropylamine), \textbf{primaire} (1 liaison \ce{N-C}) ;
    \item \chemfig{CH_3-CH_2-NH-CH_3} : N-méthyléthylamine (éthylméthylamine), \textbf{secondaire} (2 liaisons \ce{N-C}) ;
    \item \chemfig{N(CH_3)_3} : triméthylamine, \textbf{tertiaire} (3 liaisons \ce{N-C}).
  \end{itemize}
  \item L'azote de \chemfig{C_6H_5-NH_2} n'est lié qu'à \textbf{un} groupe carboné (le phényle) : c'est une amine \textbf{primaire}, même si le groupe est aromatique. Son nom systématique est \emph{phénylamine}, \emph{aniline} étant le nom retenu.
  \item La triméthylamine (tertiaire) fixe un quatrième groupe méthyle : on obtient l'\textbf{ion tétraméthylammonium} \ce{(CH3)4N+}, associé à l'ion iodure \ce{I-} : iodure de tétraméthylammonium \ce{(CH3)4N+ I-}. C'est la réaction d'Hofmann, étudiée en détail dans la section 5.
\end{enumerate}
\end{corrige}

\coursec{Nomenclature des amines}

Dans la section précédente, nous avons vu comment on passe de la molécule d'ammoniac \ce{NH3} aux amines en remplaçant un, deux ou trois atomes d'hydrogène par des groupes carbonés, ce qui définit les trois classes : amines \cle{primaires}, \cle{secondaires} et \cle{tertiaires}. Mais à l'écrit du Baccalauréat, savoir reconnaître une amine ne suffit pas : il faut aussi savoir la \textbf{nommer correctement}. Or il existe \emph{deux} systèmes de nomenclature qui cohabitent : la nomenclature systématique (officielle, type IUPAC) et la nomenclature courante (très utilisée au Bac). Un bon élève de Terminale doit maîtriser les deux, et surtout savoir passer de l'une à l'autre sans faute. C'est l'objet de cette section.

\courssub{Principe général : deux façons de nommer une amine}

Regardons la molécule \ce{CH3-CH2-NH2}. Elle est constituée d'un groupe éthyle fixé sur un atome d'azote portant deux hydrogènes. Deux lectures sont possibles :

\begin{itemize}
\item On considère que le groupe fonctionnel est \ce{-NH2} greffé sur une chaîne carbonée : c'est le point de vue de la \cle{nomenclature systématique}, qui traite l'amine comme un alcane « aminé ». On obtient alors le nom \textbf{éthanamine}.
\item On considère que la molécule est de l'ammoniac dont un hydrogène a été remplacé par un groupe éthyle : c'est le point de vue de la \cle{nomenclature courante}. On nomme alors le(s) groupe(s) alkyle(s) suivi(s) du mot « amine » : \textbf{éthylamine}.
\end{itemize}

Les deux noms désignent la même molécule. Retiens cette idée directrice : \emph{la nomenclature systématique raisonne sur la chaîne carbonée, la nomenclature courante raisonne sur les groupes fixés à l'azote}.

\courssub{La nomenclature systématique : les alcanamines}

\begin{definition}[Règles de la nomenclature systématique]
Pour nommer une amine de manière systématique :
\begin{enumerate}
\item On choisit la \cle{chaîne carbonée la plus longue} contenant le carbone lié à l'azote ; elle donne le nom de l'alcane correspondant.
\item On remplace le « e » final de l'alcane par le suffixe \cle{-amine}.
\item On fait précéder le suffixe d'un \cle{indice de position} indiquant le numéro du carbone porteur du groupe azoté, en numérotant la chaîne de façon à donner à cet indice la valeur \textbf{la plus petite possible}.
\item Les éventuels substituants sur la chaîne sont cités en préfixe, par ordre alphabétique, comme pour les alcènes et les alcools.
\end{enumerate}
\end{definition}

\begin{exemplebox}[Nommer \ce{CH3-CH(NH2)-CH2-CH3}]
\textbf{Étape 1.} La chaîne la plus longue contient 4 atomes de carbone : c'est un dérivé du \textbf{butane}.

\textbf{Étape 2.} Le groupe \ce{-NH2} est porté par un carbone. Numérotons la chaîne dans les deux sens :
\begin{itemize}
\item de gauche à droite : \ce{NH2} est sur le carbone 2 ;
\item de droite à gauche : \ce{NH2} serait sur le carbone 3.
\end{itemize}

\textbf{Étape 3.} On retient l'indice le plus petit : 2. Le nom est donc \textbf{butan-2-amine} (et \emph{surtout pas} butan-3-amine).
\end{exemplebox}

\begin{attention}[Erreur fréquente : le sens de numérotation]
Beaucoup d'élèves numérotent la chaîne « dans le sens où elle est écrite ». C'est faux : on numérote toujours de manière à donner au groupe fonctionnel \ce{-NH2} l'\textbf{indice le plus petit}. Ainsi \ce{CH3-CH2-CH(NH2)-CH3} s'écrit de gauche à droite, mais se nomme \textbf{butan-2-amine} car on numérote depuis la droite. Vérifie systématiquement les deux sens de numérotation avant de conclure.
\end{attention}

\courssub{La nomenclature courante : les alkylamines}

\begin{definition}[Règles de la nomenclature courante]
On nomme les \cle{groupes alkyles} fixés sur l'atome d'azote \textbf{par ordre alphabétique}, puis on ajoute le mot \cle{amine} (collé, sans espace). Si plusieurs groupes identiques sont présents, on utilise les préfixes multiplicatifs « di- » et « tri- ».
\end{definition}

\begin{exemplebox}[Exemples de noms courants]
\begin{itemize}
\item \ce{CH3-NH2} : méthylamine ;
\item \ce{(CH3)2NH} : diméthylamine ;
\item \ce{(CH3)3N} : triméthylamine ;
\item \ce{CH3-CH2-NH-CH3} : éthylméthylamine (éthyle avant méthyle, ordre alphabétique) ;
\item \ce{CH3-CH2-N(CH3)-CH2-CH3} : diéthylméthylamine.
\end{itemize}
\end{exemplebox}

La nomenclature courante est très prisée au Baccalauréat car elle est rapide. Mais attention à l'\textbf{ordre alphabétique} : dans « éthylméthylamine », le « é » de éthyle précède le « m » de méthyle. On écrirait « méthylpropylamine » et non « propylméthylamine ».

\courssub{Amines secondaires et tertiaires : les préfixes N- et N,N-}

Voici le point le plus délicat de la nomenclature systématique des amines. Considérons l'amine secondaire \ce{CH3-CH2-NH-CH3}. Deux groupes carbonés sont fixés sur l'azote : un éthyle et un méthyle. Lequel choisir comme « chaîne principale » ?

\begin{propriete}[Règle des préfixes N-]
Pour une amine \textbf{secondaire} ou \textbf{tertiaire} dont les groupes fixés sur l'azote sont différents :
\begin{enumerate}
\item on choisit comme chaîne principale le groupe carboné \textbf{le plus long} (ou le plus substitué), qui donne le nom de l'alcanamine ;
\item les autres groupes, directement liés à l'azote, sont cités en préfixe avec le \cle{locant N-} (ou \cle{N,N-} s'il y en a deux), qui précise qu'ils sont portés par l'\textbf{atome d'azote} et non par un carbone de la chaîne.
\end{enumerate}
\end{propriete}

\begin{exemplebox}[Application de la règle des préfixes N-]
\begin{itemize}
\item \ce{CH3-CH2-NH-CH3} : la chaîne principale est l'éthyle (2 C) ; le méthyle est sur l'azote : \textbf{N-méthyléthanamine}.
\item \ce{CH3-CH2-CH2-N(CH3)2} : chaîne principale à 3 C, deux méthyles sur l'azote : \textbf{N,N-diméthylpropanamine} (ou N,N-diméthylpropan-1-amine).
\item \ce{(CH3)2CH-NH-CH2-CH3} : chaîne principale à 2 C : \textbf{N-éthylpropan-2-amine}.
\end{itemize}
\end{exemplebox}

\begin{attention}[Piège classique : oublier le préfixe N-]
L'erreur la plus fréquente au Bac est d'écrire « méthyléthanamine » sans le « N- ». Or « 2-méthyléthanamine » ou « méthyléthanamine » suggère que le méthyle est fixé sur un \emph{carbone} de la chaîne, ce qui correspond à une molécule différente (un isomère !). Le locant \textbf{N-} est \emph{obligatoire} dès qu'un substituant est porté par l'azote d'une amine secondaire ou tertiaire. Sans lui, le nom est faux et la copie perd des points.
\end{attention}

\begin{methode}[Pour nommer une amine sans se tromper]
\begin{enumerate}
\item \textbf{Repérer la classe} : compter le nombre de groupes carbonés liés à l'azote (1 = primaire, 2 = secondaire, 3 = tertiaire).
\item \textbf{Identifier la chaîne principale} (la plus longue contenant le C lié à N) et les groupes fixés directement sur l'azote.
\item \textbf{Nomenclature systématique} : nom de l'alcane + indice de position (le plus petit possible) + suffixe « -amine » ; ajouter les préfixes N- ou N,N- pour les groupes sur l'azote.
\item \textbf{Nomenclature courante} : citer les groupes alkyles par ordre alphabétique + « amine ».
\item \textbf{Vérifier} : le nom obtenu doit permettre de réécrire la formule sans ambiguïté.
\end{enumerate}
\end{methode}

\courssub{Cas particuliers à connaître impérativement}

\textbf{L'aniline.}\par
Lorsque le groupe \ce{-NH2} est directement fixé sur un noyau benzénique, on obtient \ce{C6H5-NH2}. Son nom systématique est \cle{phénylamine}, mais son nom usuel, universellement employé (y compris au Bac), est \cle{aniline}. C'est une amine primaire aromatique, liquide à température ordinaire, qui s'oxyde à l'air en brunissant.

\begin{savaistu}[L'aniline, mère des colorants et du paracétamol]
L'aniline est historiquement le point de départ de l'industrie des \textbf{colorants} : c'est à partir de dérivés de l'aniline que furent synthétisés la mauvéine (premier colorant de synthèse, 1856) puis l'\textbf{indigo}, le bleu des pagnes et des jeans. L'aniline sert aussi de matière première à la fabrication du \textbf{paracétamol}, l'antalgique le plus consommé au Cameroun et dans le monde. Une simple molécule, \ce{C6H5-NH2}, se cache ainsi derrière des pans entiers de l'industrie chimique et pharmaceutique !
\end{savaistu}

\textbf{Les diamines usuelles.}\par
Certaines molécules portent \emph{deux} groupes amine. On les nomme avec le suffixe « -diamine » et deux indices de position :

\begin{itemize}
\item \ce{H2N-CH2-CH2-NH2} : \textbf{éthane-1,2-diamine} (nom courant : éthylènediamine), très utilisée comme agent complexant ;
\item \ce{H2N-(CH2)6-NH2} : \textbf{hexane-1,6-diamine} (ou hexaméthylènediamine). C'est l'un des deux monomères du \textbf{nylon-6,6}, obtenu par polycondensation avec l'acide hexanedioïque (acide adipique) : les deux « 6 » du nom du nylon rappellent les 6 atomes de carbone de chacun des deux monomères.
\end{itemize}

\courssub{Formules et noms : le tableau de correspondance}

\begin{popfigure}
\centering
\begin{tikzpicture}[scale=0.92, transform shape]
\node at (0,0) {\chemfig{CH_3-CH_2-CH_2-CH_2-NH_2}};
\node[below, popblue, font=\small\bfseries] at (0,-0.7) {butan-1-amine (butylamine)};

\node at (6.2,0) {\chemfig{CH_3-CH(-[2]NH_2)-CH_2-CH_3}};
\node[below, poporange, font=\small\bfseries] at (6.2,-0.7) {butan-2-amine};

\node at (0,-3.2) {\chemfig{CH_3-CH_2-CH_2-NH-CH_3}};
\node[below, popgreen, font=\small\bfseries] at (0,-3.9) {N-méthylpropanamine (méthylpropylamine)};

\node at (6.2,-3.2) {\chemfig{*6(=-=-(-NH_2)=-)}};
\node[below, poppurple, font=\small\bfseries] at (6.2,-3.9) {phénylamine (aniline)};
\end{tikzpicture}
\legende{Quatre amines de référence avec leurs noms systématiques (et courants). Remarque le préfixe N- pour l'amine secondaire et le cas particulier de l'aniline.}
\end{popfigure}

\begin{center}
\begin{tabularx}{\linewidth}{|l|X|X|}
\hline
\rowcolor{popblueL} \textbf{Formule semi-développée} & \textbf{Nom systématique} & \textbf{Nom courant} \\
\hline
\ce{CH3-NH2} & méthanamine & méthylamine \\
\hline
\ce{CH3-CH2-NH2} & éthanamine & éthylamine \\
\hline
\ce{CH3-CH2-CH2-NH2} & propan-1-amine & propylamine \\
\hline
\ce{CH3-CH(NH2)-CH3} & propan-2-amine & isopropylamine \\
\hline
\ce{CH3-CH(NH2)-CH2-CH3} & butan-2-amine & sec-butylamine \\
\hline
\ce{CH3-NH-CH2-CH3} & N-méthyléthanamine & éthylméthylamine \\
\hline
\ce{(CH3)3N} & N,N-diméthylméthanamine & triméthylamine \\
\hline
\ce{C6H5-NH2} & phénylamine & aniline \\
\hline
\ce{H2N-(CH2)6-NH2} & hexane-1,6-diamine & hexaméthylènediamine \\
\hline
\end{tabularx}
\end{center}

\begin{aretenir}[L'essentiel de la nomenclature des amines]
\begin{itemize}
\item \textbf{Systématique} : chaîne la plus longue + suffixe « -amine » + indice de position le plus petit ; préfixes \textbf{N-} ou \textbf{N,N-} pour les groupes portés par l'azote des amines secondaires et tertiaires.
\item \textbf{Courante} : groupes alkyles par \textbf{ordre alphabétique} + « amine ».
\item \ce{C6H5-NH2} = \textbf{aniline} (phénylamine) ; \ce{H2N-(CH2)6-NH2} = hexane-1,6-diamine, monomère du nylon-6,6.
\end{itemize}
\end{aretenir}

\begin{exoresolu}[Nommer, puis retrouver une formule]
\textbf{Énoncé.} 1) Donner le nom systématique et le nom courant de l'amine de formule \ce{CH3-CH2-N(CH3)-CH2-CH2-CH3}. Préciser sa classe.

2) Écrire la formule semi-développée de la N-éthylbutan-2-amine.

\tcblower
\textbf{Solution détaillée.}

\textbf{1)} L'azote porte trois groupes carbonés (éthyle, méthyle, propyle) : c'est une amine \textbf{tertiaire}.

\emph{Nom systématique} : la chaîne la plus longue liée à l'azote compte 3 carbones (propyle) ; les groupes éthyle et méthyle sont portés par l'azote, cités par ordre alphabétique avec le locant N- : \textbf{N-éthyl-N-méthylpropanamine} (ou N-éthyl-N-méthylpropan-1-amine).

\emph{Nom courant} : groupes par ordre alphabétique (éthyle, méthyle, propyle) : \textbf{éthylméthylpropylamine}.

\textbf{2)} « Butan-2-amine » impose une chaîne de 4 carbones avec \ce{-NH-} sur le carbone 2 ; « N-éthyl » indique un groupe éthyle fixé sur l'azote. D'où la formule :
\[ \ce{CH3-CH(NH-CH2-CH3)-CH2-CH3} \]
c'est-à-dire \chemfig{CH_3-CH(-[2]NH-CH_2-CH_3)-CH_2-CH_3}. Vérification : l'azote porte deux groupes carbonés, c'est bien une amine secondaire, cohérent avec le préfixe N- unique.
\end{exoresolu}

\begin{exercice}[À toi de jouer]
1) Nommer (nomenclatures systématique et courante) les amines suivantes et préciser leur classe :
\begin{itemize}
\item \ce{CH3-CH2-CH2-CH(NH2)-CH3} ;
\item \ce{(CH3-CH2)2NH} ;
\item \ce{CH3-N(CH2-CH3)-CH2-CH2-CH2-CH3}.
\end{itemize}
2) Écrire les formules semi-développées de : a) la diéthylamine ; b) la N,N-diméthyléthanamine ; c) l'éthane-1,2-diamine.
\end{exercice}

\begin{corrige}[Exercice n°1]
\textbf{1)} a) Chaîne de 5 C, \ce{-NH2} sur le carbone 2 (numérotation depuis la droite) : \textbf{pentan-2-amine} (nom courant : 1-méthylbutylamine, rarement exigé) ; amine primaire.

b) \ce{(C2H5)2NH} : deux éthyles sur l'azote : \textbf{N-éthyléthanamine} (courant : \textbf{diéthylamine}) ; amine secondaire.

c) Azote portant méthyle, éthyle et butyle : \textbf{N-éthyl-N-méthylbutanamine} (courant : \textbf{butyléthylméthylamine}) ; amine tertiaire.

\textbf{2)} a) \ce{CH3-CH2-NH-CH2-CH3} ; b) \ce{CH3-CH2-N(CH3)2} ; c) \ce{H2N-CH2-CH2-NH2}.
\end{corrige}

Tu sais désormais nommer n'importe quelle amine et retrouver sa formule à partir de son nom. Cette compétence, apparemment anodine, conditionne la réussite de toute la suite du chapitre : impossible d'écrire correctement l'équation d'une réaction d'Hofmann ou la préparation d'un amide si l'on confond N-méthyléthanamine et propanamine ! Dans la section suivante, nous verrons comment la structure des amines — et en particulier leurs liaisons hydrogène — explique leurs propriétés physiques remarquables.

\coursec{Propriétés physiques des amines}

Dans la section précédente, tu as appris à nommer les amines et à reconnaître leur classe. Mais une question naturelle se pose maintenant : à quoi \emph{ressemblent} ces composés ? Sont-ils gazeux, liquides, solides ? Se dissolvent-ils dans l'eau ? Pourquoi certains poissons un peu trop attendus dégagent-ils cette odeur âcre si caractéristique ? Toutes ces propriétés physiques s'expliquent par un seul fait structural que tu connais déjà : la présence, sur l'atome d'azote, d'un \cle{doublet non liant} et, pour les amines primaires et secondaires, de liaisons \ce{N-H} polarisées. C'est toute la clé de cette section.

\courssub{État physique des amines à température ambiante}

\begin{propriete}[État physique]
À température ambiante (\SI{25}{\celsius}) :
\begin{itemize}
\item les amines de très faible masse molaire (\ce{CH3NH2}, \ce{C2H5NH2}, \ce{(CH3)3N}) sont \cle{gazeuses} ;
\item les amines de masse molaire intermédiaire sont \cle{liquides} ;
\item les amines à longue chaîne carbonée ou aromatiques (comme l'aniline \ce{C6H5NH2}) sont liquides visqueux ou \cle{solides}.
\end{itemize}
\end{propriete}

Prenons trois exemples concrets. La méthylamine \ce{CH3NH2} ($M = \SI{31}{g.mol^{-1}}$) bout à \SI{-6}{\celsius} : c'est un gaz à température ambiante. L'éthylamine \ce{C2H5NH2} ($M = \SI{45}{g.mol^{-1}}$) bout à \SI{17}{\celsius} : gaz ou liquide très volatil selon la saison. L'aniline \ce{C6H5NH2} ($M = \SI{93}{g.mol^{-1}}$) bout à \SI{184}{\celsius} : c'est un liquide huileux. On retrouve la tendance générale de la chimie organique : \textbf{plus la masse molaire augmente, plus la température d'ébullition augmente}. Mais ce n'est pas toute l'histoire, car à masse comparable, les amines bouillent beaucoup plus haut que les alcanes. Pourquoi ?

\courssub{Les liaisons hydrogène : la clé de tout}

\textbf{Polarité de la liaison N-H et doublet non liant}\par

L'atome d'azote est plus électronégatif que l'hydrogène ($\chi_N = 3{,}0$ contre $\chi_H = 2{,}2$). La liaison \ce{N-H} est donc \cle{polaire} : l'azote porte une charge partielle $\delta^-$ et l'hydrogène une charge partielle $\delta^+$. De plus, l'azote possède un \cle{doublet non liant} disponible. Les conditions sont réunies pour l'établissement de \cle{liaisons hydrogène} : un hydrogène porté par un azote (donneur) est attiré par le doublet non liant d'un autre azote (accepteur).

\begin{definition}[Liaison hydrogène dans les amines]
Une \cle{liaison hydrogène} \ce{N-H\cdots N} est une interaction attractive entre un atome d'hydrogène lié à un azote (pôle $\delta^+$) et le doublet non liant d'un autre atome d'azote (pôle $\delta^-$). Seules les amines \textbf{primaires} (\ce{R-NH2}) et \textbf{secondaires} (\ce{R2NH}) peuvent en former \emph{entre leurs propres molécules}, car elles possèdent au moins un H sur l'azote.
\end{definition}

\begin{popfigure}
\begin{center}
\begin{tikzpicture}[scale=1.0]
% Ethylamine - ethylamine H bond
\node at (-4.2,2.6) {\small\textbf{a) Entre molécules d'éthylamine}};
\node at (-5.6,0.6) {\chemfig{CH_3-CH_2-N(-[2]H)(-[6]H)}};
\node at (-2.6,0.6) {\chemfig{CH_3-CH_2-N(-[2]H)(-[6]H)}};
\draw[popblue,thick,dashed] (-4.35,1.35) -- (-3.15,1.35);
\node[popblue] at (-3.75,1.65) {\scriptsize liaison H};
\node at (-4.35,1.6) {\scriptsize $\delta^+$};
\node at (-3.1,1.6) {\scriptsize $\delta^-$};
% Ethylamine - water H bond
\node at (3.4,2.6) {\small\textbf{b) Entre éthylamine et eau}};
\node at (2.0,0.6) {\chemfig{CH_3-CH_2-N(-[2]H)(-[6]H)}};
\node at (5.2,0.6) {\chemfig{H-O(-[2]H)}};
\draw[poporange,thick,dashed] (3.25,1.35) -- (4.55,1.35);
\node[poporange] at (3.9,1.65) {\scriptsize liaison H};
\node at (3.3,1.6) {\scriptsize $\delta^+$};
\node at (4.6,1.6) {\scriptsize $\delta^-$};
\draw[popgreen,thick,dashed] (3.3,-0.15) -- (4.35,-0.15);
\node[popgreen] at (3.8,-0.45) {\scriptsize liaison H};
\end{tikzpicture}
\end{center}
\legende{Liaisons hydrogène (pointillés) : a) association entre molécules d'éthylamine \ce{N-H\cdots N} ; b) liaisons hydrogène avec l'eau — l'amine peut être donneuse (\ce{N-H\cdots O}) ou acceptrice (\ce{O-H\cdots N}) grâce à son doublet non liant.}
\end{popfigure}

\begin{attention}[Amines tertiaires et liaisons hydrogène]
Erreur très fréquente au Baccalauréat : affirmer que les amines \textbf{tertiaires} \ce{R3N} forment des liaisons hydrogène \emph{entre elles}. C'est \textbf{faux} : elles ne portent \textbf{aucun hydrogène sur l'azote}, elles ne peuvent donc pas être \emph{donneuses}. En revanche, leur doublet non liant leur permet d'\emph{accepter} des liaisons hydrogène avec l'eau (\ce{O-H\cdots N}), ce qui explique leur solubilité. Retiens : \textbf{pas de N-H, pas d'association entre molécules d'amines tertiaires}.
\end{attention}

\courssub{Conséquence sur les températures d'ébullition}

Pour faire bouillir un liquide, il faut séparer les molécules les unes des autres, donc rompre les liaisons hydrogène intermoléculaires : cela coûte de l'énergie. D'où la hiérarchie suivante, à masse molaire voisine :

\begin{aretenir}[Hiérarchie des températures d'ébullition]
À masse molaire comparable :
\[ \text{alcane} < \text{amine tertiaire} < \text{amine primaire} < \text{alcool} \]
\begin{itemize}
\item \textbf{Alcane} : molécule apolaire, seules de faibles interactions de Van der Waals existent.
\item \textbf{Amine tertiaire} : polaire, mais pas de liaison hydrogène intermoléculaire.
\item \textbf{Amine primaire} : liaisons hydrogène \ce{N-H\cdots N}.
\item \textbf{Alcool} : liaisons hydrogène \ce{O-H\cdots O}, \emph{plus fortes} car l'oxygène est plus électronégatif que l'azote ($\chi_O = 3{,}4 > \chi_N = 3{,}0$) : la liaison \ce{O-H} est plus polarisée que \ce{N-H}.
\end{itemize}
\end{aretenir}

\begin{exemplebox}[Comparaison chiffrée à masse voisine]
Comparons trois composés de masse molaire proche ($M \approx \SIrange{44}{46}{g.mol^{-1}}$) :
\begin{center}
\begin{tabularx}{\linewidth}{|l|c|c|X|}
\hline
\rowcolor{popblueL} \textbf{Composé} & $M$ (\si{g.mol^{-1}}) & $T_{\text{éb}}$ (\si{\celsius}) & \textbf{Interaction} \\
\hline
Propane \ce{C3H8} & 44 & $-42$ & Van der Waals seules \\
\hline
Triméthylamine \ce{(CH3)3N} & 59 & 3 & polaire, sans liaison H intermoléculaire \\
\hline
Éthylamine \ce{C2H5NH2} & 45 & 17 & liaisons H \ce{N-H\cdots N} \\
\hline
Éthanol \ce{C2H5OH} & 46 & 78 & liaisons H \ce{O-H\cdots O} (fortes) \\
\hline
\end{tabularx}
\end{center}
Commentaire : l'éthylamine bout \SI{59}{\celsius} plus haut que le propane grâce à ses liaisons hydrogène, mais \SI{61}{\celsius} plus bas que l'éthanol car la liaison \ce{N-H\cdots N} est plus faible que \ce{O-H\cdots O}. La triméthylamine, bien que plus lourde que l'éthylamine, bout plus bas : c'est une amine tertiaire, sans \ce{N-H}.
\end{exemplebox}

\begin{exemplebox}[Isomères : primaire contre tertiaire]
La propan-1-amine \ce{C3H7NH2} (primaire) bout à \SI{49}{\celsius} ; son isomère la triméthylamine \ce{(CH3)3N} (tertiaire) bout à \SI{3}{\celsius}, et le propan-1-ol \ce{C3H7OH} à \SI{97}{\celsius}. L'écart de \SI{46}{\celsius} entre les deux amines isomères illustre parfaitement le rôle des liaisons hydrogène intermoléculaires, présentes chez la primaire, absentes chez la tertiaire.
\end{exemplebox}

\begin{popfigure}
\begin{center}
\begin{tikzpicture}
\begin{axis}[
width=0.95\linewidth, height=7cm,
xlabel={Masse molaire $M$ (\si{g.mol^{-1}})},
ylabel={$T_{\text{éb}}$ (\si{\celsius})},
xmin=25, xmax=105, ymin=-60, ymax=120,
grid=both, grid style={popline!40},
legend style={at={(0.03,0.97)}, anchor=north west, font=\small},
tick label style={font=\small}, label style={font=\small}
]
% Alcanes : ethane 30/-89, propane 44/-42, butane 58/-0.5, pentane 72/36, hexane 86/69
\addplot[mark=*, popdark, thick] coordinates {(30,-89) (44,-42) (58,-0.5) (72,36) (86,69)};
\addlegendentry{Alcanes}
% Amines primaires : methylamine 31/-6, ethylamine 45/17, propylamine 59/49, butylamine 73/78, pentylamine 87/104
\addplot[mark=square*, popblue, thick] coordinates {(31,-6) (45,17) (59,49) (73,78) (87,104)};
\addlegendentry{Amines primaires}
% Alcools : methanol 32/65, ethanol 46/78, propanol 60/97, butanol 74/118
\addplot[mark=triangle*, poporange, thick] coordinates {(32,65) (46,78) (60,97) (74,118)};
\addlegendentry{Alcools}
\end{axis}
\end{tikzpicture}
\end{center}
\legende{Température d'ébullition en fonction de la masse molaire. À masse comparable : alcool $>$ amine primaire $>$ alcane. L'écart entre alcools et amines traduit la plus grande force des liaisons \ce{O-H\cdots O} par rapport à \ce{N-H\cdots N}.}
\end{popfigure}

\courssub{Solubilité des amines dans l'eau}

\begin{propriete}[Solubilité]
Les amines à \textbf{courte chaîne carbonée} (moins de 6 atomes de carbone) sont \cle{solubles dans l'eau}, quelle que soit leur classe (primaire, secondaire \emph{ou tertiaire}), car elles établissent des liaisons hydrogène avec les molécules d'eau :
\begin{itemize}
\item comme \emph{donneur} si elles portent un \ce{N-H} : \ce{N-H\cdots O-H} ;
\item comme \emph{accepteur} grâce au doublet de l'azote : \ce{R3N\cdots H-O} (possible même pour les tertiaires).
\end{itemize}
La solubilité \textbf{diminue lorsque la chaîne carbonée s'allonge}, car la partie hydrocarbonée, hydrophobe, l'emporte sur la tête polaire \ce{-NH2}. Les amines à plus de 6 carbones sont pratiquement insolubles.
\end{propriete}

C'est exactement le même raisonnement que pour les alcools : la solubilité résulte d'un compromis entre la partie \emph{hydrophile} (le groupe \ce{-NH2}, qui « aime » l'eau) et la partie \emph{hydrophobe} (la chaîne carbonée, qui la « fuit »). La méthylamine et l'éthylamine se mélangent à l'eau en toutes proportions ; l'aniline n'y est que très peu soluble (\SI{3,6}{g} par litre environ).

\begin{methode}[Expliquer une propriété physique d'une amine]
Face à toute question sur une propriété physique d'une amine, adopte ce réflexe en trois étapes :
\begin{enumerate}
\item \textbf{Repérer} la classe de l'amine et la présence ou non de liaisons \ce{N-H} (polarisées, car $\chi_N > \chi_H$).
\item \textbf{Invoquer les liaisons hydrogène adaptées} : entre molécules d'amines (\ce{N-H\cdots N}) pour expliquer l'ébullition ; avec l'eau (\ce{N-H\cdots O} et \ce{O-H\cdots N}) pour expliquer la solubilité.
\item \textbf{Comparer} si besoin avec l'alcool de masse voisine : la liaison \ce{O-H\cdots O} étant plus forte que \ce{N-H\cdots N}, l'alcool bout plus haut.
\end{enumerate}
\end{methode}

\courssub{Odeur des amines}

Les amines volatiles possèdent une odeur \cle{âcre}, \cle{ammoniacale} ou de \cle{poisson} très caractéristique, facilement perceptible même à faible concentration. C'est la triméthylamine \ce{(CH3)3N} qui est en grande partie responsable de l'odeur du poisson qui se gâte : elle provient de la dégradation bactérienne de certains composés azotés de la chair. L'odeur ammoniacale rappelle, elle, la parenté structurale avec \ce{NH3}.

\begin{savaistu}[Putrescine et cadavérine : des noms qui sentent mauvais]
Lors de la décomposition de la viande, les bactéries dégradent les acides aminés et produisent deux diamines aux noms terriblement évocateurs : la \textbf{putrescine} \ce{H2N-(CH2)4-NH2} (butane-1,4-diamine) et la \textbf{cadavérine} \ce{H2N-(CH2)5-NH2} (pentane-1,5-diamine). Leur odeur nauséabonde est un signal d'alerte biologique : elle nous détourne des aliments avariés. À l'inverse, certaines amines sont précieuses : l'adrénaline, la dopamine et la sérotonine, messagers chimiques de notre système nerveux, sont des amines !
\end{savaistu}

\begin{exoresolu}[Justifier des températures d'ébullition]
\textbf{Énoncé.} On considère les trois composés suivants : le butane \ce{C4H10} ($M = \SI{58}{g.mol^{-1}}$, $T_{\text{éb}} = \SI{-0,5}{\celsius}$), la propan-1-amine \ce{C3H7NH2} ($M = \SI{59}{g.mol^{-1}}$, $T_{\text{éb}} = \SI{49}{\celsius}$) et le butan-1-ol \ce{C4H9OH} ($M = \SI{74}{g.mol^{-1}}$, $T_{\text{éb}} = \SI{118}{\celsius}$). Expliquer l'ordre des températures d'ébullition.
\tcblower
\textbf{Solution détaillée.}
\begin{enumerate}
\item \textbf{Butane} : molécule apolaire, seules les faibles interactions de Van der Waals relient les molécules ; elles se séparent facilement, d'où une ébullition très basse (\SI{-0,5}{\celsius}).
\item \textbf{Propan-1-amine} : amine primaire, elle possède deux liaisons \ce{N-H} polarisées et un doublet non liant sur l'azote. Ses molécules s'associent par liaisons hydrogène \ce{N-H\cdots N}, qu'il faut rompre pour vaporiser le liquide : $T_{\text{éb}}$ nettement plus élevée (\SI{49}{\celsius}) malgré une masse quasi identique à celle du butane.
\item \textbf{Butan-1-ol} : alcool, il forme des liaisons hydrogène \ce{O-H\cdots O}, plus fortes que \ce{N-H\cdots N} car l'oxygène est plus électronégatif que l'azote ; de plus sa masse est supérieure. Il bout donc le plus haut (\SI{118}{\celsius}).
\end{enumerate}
Conclusion : à masse voisine, alcane $<$ amine primaire $<$ alcool, conformément à la hiérarchie des liaisons hydrogène.
\end{exoresolu}

\begin{exercice}[Pour t'entraîner]
On donne trois isomères de formule brute \ce{C3H9N} : la propan-1-amine \ce{C3H7NH2} ($T_{\text{éb}} = \SI{49}{\celsius}$), la N-méthyléthylamine \ce{C2H5-NH-CH3} ($T_{\text{éb}} = \SI{37}{\celsius}$) et la triméthylamine \ce{(CH3)3N} ($T_{\text{éb}} = \SI{3}{\celsius}$). Attribuer à chacune sa classe et justifier l'ordre des températures d'ébullition.
\end{exercice}

\begin{corrige}[Exercice 1]
Propan-1-amine : primaire (deux \ce{N-H}) ; N-méthyléthylamine : secondaire (un seul \ce{N-H}) ; triméthylamine : tertiaire (aucun \ce{N-H}). La primaire forme le plus de liaisons hydrogène intermoléculaires (deux donneurs H par molécule), la secondaire un peu moins (un seul donneur), la tertiaire aucune. On observe donc l'ordre : primaire (\SI{49}{\celsius}) $>$ secondaire (\SI{37}{\celsius}) $>$ tertiaire (\SI{3}{\celsius}). L'écart entre la primaire et la secondaire s'explique par le nombre de donneurs H ; la triméthylamine, sans \ce{N-H}, ne s'associe pas par liaisons hydrogène intermoléculaires et reste la plus volatile.
\end{corrige}

En résumé, une seule idée directrice gouverne toutes les propriétés physiques des amines : la polarité de la liaison \ce{N-H} et le doublet non liant de l'azote permettent des liaisons hydrogène, qui élèvent les températures d'ébullition (sauf pour les tertiaires) et assurent la solubilité dans l'eau des amines courtes. Ce même doublet non liant a une autre conséquence capitale, cette fois \emph{chimique} : il confère aux amines un caractère basique, que nous étudions dans la section suivante.

\coursec{Caractère basique des amines}

Dans la section précédente, nous avons vu que les amines sont des composés volatils, souvent odorants, et que les plus légères d'entre elles se dissolvent facilement dans l'eau grâce aux liaisons hydrogène. Mais que se passe-t-il \emph{chimiquement} lorsqu'une amine se dissout dans l'eau ? Si tu plonges un papier pH dans une solution d'éthylamine, il vire au bleu foncé : la solution est \cle{basique}. Cette propriété n'est pas un détail : c'est le comportement le plus important des amines, celui qui explique leur rôle dans les médicaments, les colorants et de nombreuses synthèses industrielles. Toute cette section repose sur un unique atout de la molécule : le \cle{doublet non liant} porté par l'atome d'azote.

\courssub{Origine de la basicité : le doublet non liant de l'azote}

Reprenons l'intuition avant la définition. Un proton \ce{H+} est un atome d'hydrogène qui a perdu son électron : c'est une \og case vide \fg{} électronique, chargée positivement. L'atome d'azote d'une amine, lui, possède un doublet d'électrons \og disponibles \fg{}, non engagés dans une liaison. L'attraction est immédiate : le doublet de l'azote vient combler la case vide du proton, formant une liaison covalente \ce{N-H} supplémentaire. L'amine joue le rôle de \emph{donneur de doublet}, le proton celui d'\emph{accepteur}.

\begin{definition}[Base de Brønsted et basicité des amines]
Une \cle{base de Brønsted} est une espèce chimique capable de \textbf{capter un proton} \ce{H+} grâce à un doublet non liant. Les amines sont des bases de Brønsted : le doublet non liant de l'atome d'azote capte un proton pour former un \cle{ion alkylammonium} :
\[ \ce{R-NH2 + H+ -> R-NH3+} \]
L'amine (base) et l'ion alkylammonium (acide conjugué) forment le couple \ce{R-NH3+}/\ce{R-NH2}.
\end{definition}

\begin{popfigure}
\begin{center}
\begin{tikzpicture}[>=Stealth, scale=1.05]
  % Amine avec doublet
  \node at (0,0) {\chemfig{R-[:30]N(-[2]H)(-[6]H)}};
  \node[popblue] at (0.55,0.55) {\large :};
  \node[popmuted, font=\small] at (0.9,0.95) {doublet non liant};
  % Flèche courbe du doublet vers H+
  \node[popink] at (3.6,0.35) {\ce{H+}};
  \draw[->, poporange, very thick, bend left=35] (0.75,0.6) to node[midway, above, font=\small, text=poporange]{capture du proton} (3.3,0.45);
  % Flèche de réaction
  \draw[->, popdark, ultra thick] (4.6,0.1) -- (5.8,0.1);
  % Produit
  \node at (7.9,0) {\chemfig{R-[:30]N^{+}(-[2]H)(-[6]H)-[,-0.6]H}};
  \node[popmuted, font=\small] at (7.9,-1.1) {ion alkylammonium};
  \node[popmuted, font=\small] at (-0.2,-1.1) {amine (base)};
\end{tikzpicture}
\end{center}
\legende{Capture d'un proton par le doublet non liant de l'azote : la flèche courbe part du doublet (donneur) vers le proton (accepteur).}
\end{popfigure}

\begin{attention}[La flèche part toujours du doublet]
Dans un mécanisme, la flèche courbe décrit un déplacement d'\textbf{électrons} : elle part du doublet non liant de l'azote et arrive sur le proton \ce{H+}. Ne la fais jamais partir de \ce{H+} : un proton n'a aucun électron à déplacer !
\end{attention}

\courssub{Action de l'eau sur une amine : un équilibre}

L'eau est un solvant \emph{amphotère} : elle peut céder un proton (elle se comporte alors en acide). Face à une amine, l'eau joue ce rôle d'acide : l'amine arrache un proton à une molécule d'eau, laissant derrière elle un ion hydroxyde \ce{HO-}. C'est la formation de ces ions \ce{HO-} qui rend la solution basique.

\begin{propriete}[Réaction d'une amine avec l'eau]
Toute amine réagit avec l'eau selon un \textbf{équilibre} (réaction limitée, d'où la double flèche) :
\[ \ce{R-NH2 + H2O <=> R-NH3+ + HO-} \]
Pour une amine secondaire ou tertiaire, on écrit de même :
\[ \ce{R2NH + H2O <=> R2NH2+ + HO-} \qquad \ce{R3N + H2O <=> R3NH+ + HO-} \]
La solution contient des ions hydroxyde \ce{HO-} : son pH est \textbf{supérieur à 7}.
\end{propriete}

\begin{methode}[Écrire l'équation de l'action de l'eau sur une amine]
\begin{enumerate}
  \item Écrire les réactifs : l'amine \ce{R-NH2} et l'eau \ce{H2O}.
  \item Identifier le proton transféré : l'eau cède \emph{un} de ses atomes d'hydrogène sous forme de \ce{H+} ; il reste \ce{HO-}.
  \item Greffer ce proton sur l'azote de l'amine : \ce{R-NH2} devient \ce{R-NH3+} (l'azote porte désormais 4 liaisons et une charge $+$).
  \item Vérifier la conservation : mêmes atomes de chaque côté, et charge totale nulle des deux côtés ($0 = (+1) + (-1)$).
  \item Placer une \textbf{double flèche} \ce{<=>} : la réaction est un équilibre, jamais une réaction totale.
\end{enumerate}
\end{methode}

\begin{exemplebox}[Avec l'éthylamine]
\[ \ce{CH3-CH2-NH2 + H2O <=> CH3-CH2-NH3+ + HO-} \]
L'ion formé est l'ion \textbf{éthylammonium}. Vérification : à gauche \ce{C2H9NO}, à droite \ce{C2H8N+ + HO-} soit \ce{C2H9NO} au total ; charges : $0$ à gauche, $(+1)+(-1)=0$ à droite. L'équation est correcte.
\end{exemplebox}

\begin{attention}[Piège classique du Baccalauréat]
L'erreur la plus fréquente est d'écrire une flèche simple \ce{->}. L'action de l'eau sur une amine est un \textbf{équilibre} : seule une faible fraction des molécules d'amine est transformée en ions. C'est pourquoi on dit que les amines sont des \cle{bases faibles} et que leurs solutions sont \emph{faiblement} basiques. Toute équation \ce{R-NH2 + H2O -> ...} avec flèche simple est fausse et coûte des points.
\end{attention}

\begin{exoresolu}[Le pH d'une solution d'éthylamine]
\textbf{Énoncé.} On dissout de l'éthylamine \ce{C2H5NH2} dans l'eau distillée de façon à obtenir une solution de concentration $C = \SI{0,10}{mol.L^{-1}}$. La mesure au pH-mètre donne $\mathrm{pH} \approx 11,9$. (a) Écrire l'équation de la réaction de l'éthylamine avec l'eau. (b) Calculer la concentration en ions \ce{HO-} et la comparer à celle d'une solution d'hydroxyde de sodium (soude) de même concentration, dont le pH vaut 13. Conclure.
\tcblower
\textbf{Solution.} (a) L'éthylamine est une base faible :
\[ \ce{C2H5NH2 + H2O <=> C2H5NH3+ + HO-} \]
(b) À \SI{25}{\celsius}, $\mathrm{pH} + \mathrm{pOH} = 14$, donc $\mathrm{pOH} = 14 - 11,9 = 2,1$ et :
\[ [\ce{HO-}] = 10^{-2,1} \approx \SI{8,1e-3}{mol.L^{-1}} \]
Pour la soude, base forte totalement dissociée : $[\ce{HO-}] = C = \SI{0,10}{mol.L^{-1}}$ ($\mathrm{pOH}=1$, $\mathrm{pH}=13$). Le rapport des concentrations est $\dfrac{0,10}{8,1e-3} \approx 12$ : la solution d'éthylamine contient environ \textbf{12 fois moins} d'ions \ce{HO-} que la soude de même concentration. Conclusion : l'éthylamine est bien une base, mais une \textbf{base faible} — seule une petite partie des molécules a capté un proton. C'est la signature expérimentale de l'équilibre.
\end{exoresolu}

\courssub{Comparaison de basicité : aniline, ammoniac, amines aliphatiques}

Toutes les amines ne se valent pas. La force d'une base dépend de la \textbf{disponibilité du doublet} de l'azote : plus le doublet est \og riche \fg{} et accessible, plus la base est forte.

\begin{itemize}
  \item \textbf{Amines aliphatiques} (\ce{R-NH2}, \ce{R2NH}, \ce{R3N}) : les groupes alkyles \ce{R} sont \emph{donneurs d'électrons} (effet inductif donneur). Ils \og poussent \fg{} de la densité électronique vers l'azote : le doublet est enrichi, plus attractif pour \ce{H+}. Les amines aliphatiques sont donc des bases \textbf{plus fortes que l'ammoniac}.
  \item \textbf{Ammoniac} \ce{NH3} : l'azote ne porte que des atomes d'hydrogène, sans effet donneur. Basicité intermédiaire.
  \item \textbf{Aniline} (phénylamine) \ce{C6H5NH2} : le doublet de l'azote est \emph{délocalisé} dans le cycle benzénique par conjugaison. Moins disponible sur l'azote, il capte difficilement un proton : l'aniline est une base \textbf{plus faible que l'ammoniac}.
\end{itemize}

\begin{popfigure}
\begin{center}
\begin{tikzpicture}[>=Stealth]
  \draw[->, ultra thick, popdark] (0,0) -- (12.5,0) node[below left=2pt and -30pt, font=\small]{basicité croissante};
  % Aniline
  \fill[poppink] (2,0) circle (3pt);
  \node[above, font=\small\bfseries, text=poppink] at (2,0.1) {aniline};
  \node[below, font=\footnotesize, text=popmuted, align=center] at (2,-0.15) {\ce{C6H5NH2}\\ doublet délocalisé};
  % Ammoniac
  \fill[poporange] (6,0) circle (3pt);
  \node[above, font=\small\bfseries, text=poporange] at (6,0.1) {ammoniac};
  \node[below, font=\footnotesize, text=popmuted, align=center] at (6,-0.15) {\ce{NH3}\\ référence};
  % Amines aliphatiques
  \fill[popblue] (10,0) circle (3pt);
  \node[above, font=\small\bfseries, text=popblue] at (10,0.1) {amines aliphatiques};
  \node[below, font=\footnotesize, text=popmuted, align=center] at (10,-0.15) {\ce{R-NH2}, \ce{R2NH}, \ce{R3N}\\ effet donneur des alkyles};
\end{tikzpicture}
\end{center}
\legende{Échelle de basicité : aniline $<$ ammoniac $<$ amines aliphatiques. Tout se joue sur la disponibilité du doublet de l'azote.}
\end{popfigure}

\begin{aretenir}[L'essentiel sur la force des amines]
La basicité d'une amine croît avec la disponibilité du doublet non liant de l'azote :
\[ \text{aniline} \;<\; \text{ammoniac} \;<\; \text{amines aliphatiques} \]
Effet donneur des groupes alkyles $\Rightarrow$ doublet enrichi $\Rightarrow$ base plus forte. Délocalisation du doublet dans le noyau benzénique $\Rightarrow$ doublet moins disponible $\Rightarrow$ base plus faible.
\end{aretenir}

\courssub{Mise en évidence expérimentale du caractère basique}

\textbf{Action sur les indicateurs colorés.}\par
Les ions \ce{HO-} libérés par l'équilibre avec l'eau modifient la teinte des indicateurs colorés usuels :
\begin{itemize}
  \item le \cle{bleu de bromothymol} (BBT) vire au \textbf{bleu} ;
  \item la \cle{phénolphtaléine} vire au \textbf{rose} ;
  \item le \cle{papier pH} indique une valeur \textbf{supérieure à 7} (typiquement entre 10 et 12 pour une solution diluée).
\end{itemize}
Ce sont exactement les mêmes observations qu'avec une solution d'ammoniac : le test ne distingue pas l'amine de \ce{NH3}, il prouve seulement la présence d'une base.

\begin{experience}[Action d'une amine sur une solution de sulfate de cuivre(II)]
\textbf{Protocole.} Dans un tube à essais contenant environ \SI{2}{mL} de solution bleue de sulfate de cuivre(II) \ce{CuSO4}, ajouter quelques gouttes d'une solution d'éthylamine. Agiter.

\textbf{Observation.} Il apparaît un \textbf{précipité bleu} gélatineux.

\textbf{Interprétation.} L'amine réagit d'abord avec l'eau : \ce{C2H5NH2 + H2O <=> C2H5NH3+ + HO-}. Les ions hydroxyde formés précipitent alors les ions cuivre(II) en hydroxyde de cuivre(II) :
\[ \ce{Cu^{2+} + 2 HO- -> Cu(OH)2} \]
Le précipité bleu de \ce{Cu(OH)2} met en évidence les ions \ce{HO-}, donc le caractère basique de l'amine — exactement comme le fait l'ammoniac.
\end{experience}

\begin{popfigure}
\begin{center}
\begin{tikzpicture}[scale=0.95]
  % Tube 1 : sulfate de cuivre
  \draw[thick, popdark] (0,0) -- (0,3.4) (1.4,0) -- (1.4,3.4);
  \draw[thick, popdark] (0,0) arc (180:360:0.7);
  \fill[popblue!35] (0.03,0.15) -- (0.03,1.8) -- (1.37,1.8) -- (1.37,0.15) arc (360:180:0.67);
  \draw[popblue!60] (0.03,1.8) -- (1.37,1.8);
  \node[font=\footnotesize, text=popblue!80!black] at (0.7,1.0) {\ce{Cu^{2+}}};
  \node[below, font=\small, align=center] at (0.7,-0.75) {solution de sulfate\\ de cuivre(II) (bleue)};
  % Pipette
  \draw[thick, popdark] (0.45,4.6) -- (0.45,3.1) (0.75,4.6) -- (0.75,3.1);
  \draw[thick, popdark] (0.45,4.6) arc (180:360:0.15);
  \fill[popgreen!50] (0.48,3.6) rectangle (0.72,4.5);
  \node[right, font=\small, text=popgreen!60!black] at (0.9,4.1) {éthylamine \ce{C2H5NH2}};
  \draw[->, thick, popgreen!60!black] (2.6,3.6) -- (2.6,2.6);
  % Flèche
  \draw[->, ultra thick, popdark] (2.2,1.4) -- (3.6,1.4);
  % Tube 2 : précipité
  \draw[thick, popdark] (4.4,0) -- (4.4,3.4) (5.8,0) -- (5.8,3.4);
  \draw[thick, popdark] (4.4,0) arc (180:360:0.7);
  \fill[popblue!20] (4.43,0.9) -- (4.43,1.8) -- (5.77,1.8) -- (5.77,0.9);
  \fill[popblue!70] (4.43,0.15) -- (4.43,0.9) -- (5.77,0.9) -- (5.77,0.15) arc (360:180:0.67);
  \draw[popblue!60] (4.43,1.8) -- (5.77,1.8);
  \foreach \x/\y in {4.7/0.5, 5.1/0.35, 5.5/0.6, 4.9/0.75, 5.3/0.2, 5.6/0.4}
    \fill[popblue!90!black] (\x,\y) circle (1.2pt);
  \node[font=\footnotesize, text=popblue!90!black] at (5.1,0.55) {};
  \node[below, font=\small, align=center] at (5.1,-0.75) {précipité bleu\\ de \ce{Cu(OH)2}};
  \node[right, font=\footnotesize, text=popmuted, align=left] at (6.1,0.5) {\ce{Cu^{2+} + 2 HO- -> Cu(OH)2}};
\end{tikzpicture}
\end{center}
\legende{Test caractéristique : l'ajout d'éthylamine à une solution de sulfate de cuivre(II) provoque un précipité bleu d'hydroxyde de cuivre(II), preuve de la formation d'ions \ce{HO-}.}
\end{popfigure}

\courssub{Réaction avec les acides : formation de sels d'ammonium}

Contrairement à l'eau, un acide fort cède \emph{quantitativement} ses protons. La réaction d'une amine avec un acide est alors \textbf{totale} (flèche simple) :
\[ \ce{R-NH2 + H3O+ -> R-NH3+ + H2O} \]
Avec l'acide chlorhydrique, par exemple, on obtient un \cle{sel d'alkylammonium}, le chlorure d'alkylammonium, ionique et très soluble dans l'eau :
\[ \ce{CH3-CH2-NH2 + HCl -> CH3-CH2-NH3+ + Cl-} \]
Le produit s'appelle \textbf{chlorure d'éthylammonium} (ou \og chlorhydrate d'éthylamine \fg{} dans le langage pharmaceutique). Cette réaction est aussi utilisée en sens inverse : ajouter une base forte (soude) à un sel d'ammonium \emph{régénère} l'amine, ce qui permet de la séparer d'un mélange.

\begin{savaistu}[Des sels d'ammonium dans ton armoire à pharmacie]
La plupart des principes actifs azotés (antipaludéens, antihistaminiques, anesthésiques locaux) sont des amines... souvent peu solubles dans l'eau et instables à l'air. L'astuce des pharmaciens : les transformer en \textbf{sels d'ammonium} — chlorhydrates, phosphates, sulfates — qui sont ioniques, donc \textbf{très solubles dans l'eau} et bien plus stables. C'est ainsi que la chloroquine, longtemps utilisée contre le paludisme au Cameroun, était administrée sous forme de \emph{chlorhydrate} ou de \emph{phosphate de chloroquine}, en sirop ou en injectable. Sans le caractère basique des amines, pas de médicaments buvables !
\end{savaistu}

\begin{exercice}[Pour t'entraîner]
On considère la méthylamine \ce{CH3NH2} et l'aniline \ce{C6H5NH2}.
\begin{enumerate}
  \item Écrire l'équation-bilan de la réaction de la méthylamine avec l'eau, puis celle de sa réaction avec une solution d'acide chlorhydrique.
  \item On verse quelques gouttes de BBT dans une solution aqueuse de méthylamine. Quelle teinte observe-t-on ? Justifier.
  \item Laquelle de ces deux amines est la base la plus forte ? Justifier par la structure.
  \item On ajoute de la méthylamine à une solution de sulfate de cuivre(II). Décrire l'observation et écrire l'équation de la réaction responsable du précipité.
\end{enumerate}
\end{exercice}

\begin{corrige}[Exercice 1]
\begin{enumerate}
  \item Avec l'eau (équilibre) : \ce{CH3NH2 + H2O <=> CH3NH3+ + HO-}. Avec l'acide chlorhydrique (réaction totale) : \ce{CH3NH2 + H3O+ -> CH3NH3+ + H2O} (soit, avec \ce{HCl} : \ce{CH3NH2 + HCl -> CH3NH3+ + Cl-}, chlorure de méthylammonium).
  \item Le BBT vire au \textbf{bleu} : la réaction avec l'eau produit des ions \ce{HO-}, la solution est basique ($\mathrm{pH} > 7$).
  \item La \textbf{méthylamine} est la base la plus forte : le groupe méthyle, donneur d'électrons, enrichit le doublet de l'azote. Dans l'aniline, le doublet est délocalisé dans le cycle benzénique et donc moins disponible pour capter \ce{H+}.
  \item On observe un \textbf{précipité bleu} d'hydroxyde de cuivre(II) : \ce{Cu^{2+} + 2 HO- -> Cu(OH)2}, les ions \ce{HO-} provenant de l'équilibre de la méthylamine avec l'eau.
\end{enumerate}
\end{corrige}

\begin{aretenir}[Bilan de la section]
\begin{itemize}
  \item L'amine est une \textbf{base de Brønsted} grâce au doublet non liant de l'azote, qui capte un proton \ce{H+}.
  \item Avec l'eau : \textbf{équilibre} \ce{R-NH2 + H2O <=> R-NH3+ + HO-} ; la solution est faiblement basique (pH $> 7$). Couple \ce{R-NH3+}/\ce{R-NH2}.
  \item Force des bases : aniline $<$ ammoniac $<$ amines aliphatiques (disponibilité du doublet).
  \item Tests : BBT bleu, phénolphtaléine rose, précipité bleu de \ce{Cu(OH)2} avec les ions \ce{Cu^{2+}}.
  \item Avec un acide fort : réaction \textbf{totale} \ce{R-NH2 + H3O+ -> R-NH3+ + H2O}, formant des sels d'ammonium solubles — la clé de nombreuses formes pharmaceutiques.
\end{itemize}
\end{aretenir}

Le doublet non liant de l'azote ne fait pas que capter des protons : il peut aussi attaquer des atomes de carbone porteurs de charges partielles positives. C'est le \emph{caractère nucléophile} des amines, que nous exploiterons dans la section suivante avec les réactions d'Hofmann.

\coursec{Caractère nucléophile : alkylation et réactions d'Hofmann}

Dans la section précédente, nous avons vu que le \cle{doublet non liant} de l'azote confère aux amines un \cle{caractère basique} : l'amine capte un proton \ce{H+}, c'est-à-dire un \emph{atome d'hydrogène privé d'électron}. Mais ce doublet peut faire bien plus ! Au lieu d'attaquer un simple proton, il peut attaquer un \emph{atome de carbone} porteur d'une charge partielle positive. L'amine cesse alors d'être simplement une base : elle devient un \cle{nucléophile}. C'est cette propriété qui permet de construire de nouvelles liaisons \ce{C-N}, donc de fabriquer des amines de plus en plus substituées, jusqu'aux sels d'\cle{ammonium quaternaire} que l'on retrouve dans les désinfectants de nos hôpitaux. C'est toute la richesse des \cle{réactions d'Hofmann}, que nous allons maintenant étudier pas à pas.

\courssub{Nucléophile et électrophile : deux acteurs complémentaires}

\begin{definition}[Nucléophile et électrophile]
Un \cle{nucléophile} (du grec \emph{nucleus}, noyau, et \emph{philos}, qui aime) est une espèce chimique \textbf{riche en électrons} — elle possède un \textbf{doublet non liant} ou porte une \textbf{charge négative} — qui attaque un site pauvre en électrons pour former une liaison covalente.

Un \cle{électrophile} (qui aime les électrons) est une espèce ou un site \textbf{pauvre en électrons} : cation, atome porteur d'une charge partielle positive $\delta^+$, ou atome possédant une lacune électronique.
\end{definition}

L'intuition est simple : les charges opposées s'attirent. Le nucléophile offre son doublet d'électrons ; l'électrophile l'accepte. La nouvelle liaison covalente naît de ce « don » de doublet.

Dans une amine \ce{R-NH2}, l'azote possède un doublet non liant : l'amine est donc un \textbf{nucléophile}, exactement comme l'ammoniac \ce{NH3}. Reste à trouver le partenaire électrophile.

\courssub{Le dérivé halogéné : un carbone électrophile}

Dans un dérivé halogéné \ce{R-X} (où \ce{X} = \ce{Cl}, \ce{Br}, \ce{I}), l'halogène est \textbf{plus électronégatif que le carbone}. Il attire vers lui les électrons de la liaison \ce{C-X}, qui est donc \textbf{polarisée} :

\[ \ce{R-\overset{\delta^+}{C}H2-\overset{\delta^-}{X}} \]

Le carbone porteur de l'halogène, appauvri en électrons ($\delta^+$), devient un \cle{site électrophile}. De plus, l'ion halogénure \ce{X-} est un bon \emph{groupe partant} : il peut quitter la molécule en emportant le doublet de la liaison \ce{C-X}. Toutes les conditions sont réunies pour une \cle{substitution nucléophile} : le nucléophile (l'amine) se substitue à l'halogène sur le carbone.

\begin{popfigure}
\begin{center}
\chemfig{H_3C-\chembelow{N}{\text{\tiny\textbullet\textbullet}}H_2}
\hspace{1.5cm}
\chemfig{H_3C-[@{b1}]@{X}Br}
\chemmove{
  \draw[-{Stealth},popink,thick] ([yshift=2mm]b1.north) .. controls +(0.3,0.9) and +(0.3,0.9) .. node[above,popink]{\scriptsize attaque du doublet de N sur C$^{\delta+}$} ++(-3.6,0);
  \draw[-{Stealth},poporange,thick] ([yshift=-1mm]b1.south) .. controls +(0.2,-0.5) and +(0,-0.6) .. node[below,poporange]{\scriptsize départ de \ce{Br-}} (X.south);
}
\end{center}
\legende{Substitution nucléophile : le doublet non liant de l'azote attaque le carbone $\delta^+$ du dérivé halogéné ; la liaison \ce{C-Br} se rompt et l'ion bromure part.}
\end{popfigure}

\courssub{Première alkylation : d'une amine primaire à une amine secondaire}

Faisons réagir une amine primaire sur un dérivé halogéné. Le doublet de l'azote attaque le carbone électrophile, l'halogène s'en va : une nouvelle liaison \ce{N-C} est créée, et le groupe alkyle \ce{R$'$} s'est \emph{greffé} sur l'azote. On parle d'\cle{alkylation} de l'amine.

\begin{propriete}[Alkylation d'une amine primaire]
Une amine primaire réagit avec un dérivé halogéné par \textbf{substitution nucléophile} pour donner une \textbf{amine secondaire} et un hydracide \ce{HX} :
\[ \ce{R-NH2 + R$'$-X -> R-NH-R$'$ + HX} \]
Cette réaction est \textbf{lente} ; on la réalise en \textbf{chauffant}, classiquement en \textbf{tube scellé} (les réactifs sont souvent volatils).
\end{propriete}

\begin{exemplebox}[La méthanamine attaque le bromoéthane]
\[ \ce{CH3-NH2 + CH3-CH2-Br -> CH3-NH-CH2-CH3 + HBr} \]
La méthanamine (amine primaire) est transformée en \textbf{N-méthyléthanamine} (amine secondaire). Le groupe éthyle \ce{C2H5} a remplacé l'un des hydrogènes de l'azote.
\end{exemplebox}

\begin{attention}[Le rôle de l'excès d'amine]
L'hydracide \ce{HX} formé est un acide fort : il ne reste pas libre en présence d'amine, qui est une base. Il est \textbf{neutralisé} par une seconde molécule d'amine :
\[ \ce{R-NH2 + HX -> R-NH3+ X-} \]
On obtient un \textbf{sel d'ammonium}. C'est pourquoi, si l'on veut récupérer l'amine secondaire, on travaille avec un \textbf{excès d'amine} : une partie sert de nucléophile, l'autre sert de base pour piéger \ce{HX}. L'équation globale s'écrit alors :
\[ \ce{2 R-NH2 + R$'$-X -> R-NH-R$'$ + R-NH3+ X-} \]
\end{attention}

\courssub{Les réactions d'Hofmann successives : jusqu'à l'ammonium quaternaire}

Voici le point capital de la leçon. L'amine secondaire formée possède \emph{encore} un doublet non liant sur l'azote : elle est \textbf{elle-même nucléophile} ! Si du dérivé halogéné reste disponible, la réaction \textbf{ne s'arrête pas} : elle se poursuit.

\begin{propriete}[Réactions d'Hofmann successives]
Tant que l'azote porte encore au moins un atome d'hydrogène, l'amine peut être alkylée à nouveau :
\begin{align*}
\ce{R-NH2 + R$'$-X} &\ce{-> R-NH-R$'$ + HX} \quad \text{(amine secondaire)}\\
\ce{R-NH-R$'$ + R$'$-X} &\ce{-> R-N(R$'$)2 + HX} \quad \text{(amine tertiaire)}
\end{align*}
L'amine tertiaire n'a plus d'hydrogène sur l'azote, mais elle a \textbf{toujours son doublet} : elle attaque une dernière fois pour donner un \cle{ion ammonium quaternaire}, sans dégagement de \ce{HX} (il n'y a plus de H à éliminer) :
\[ \ce{R-N(R$'$)2 + R$'$-X -> R-N(R$'$)3+ X-} \]
L'ion ammonium quaternaire \ce{R4N+} est associé à l'ion halogénure \ce{X-} : c'est un \textbf{sel d'ammonium quaternaire}, produit \textbf{final} de l'alkylation.
\end{propriete}

\begin{methode}[Écrire une réaction d'Hofmann]
\begin{enumerate}
\item \textbf{Repérer} le doublet non liant de l'azote (le nucléophile) et le carbone $\delta^+$ du dérivé halogéné (l'électrophile).
\item \textbf{Former} la nouvelle liaison \ce{N-C} : le groupe alkyle se fixe sur l'azote, l'halogène part sous forme \ce{X-}.
\item \textbf{Éliminer} \ce{HX} si l'azote portait un hydrogène (l'ion \ce{X-} arrache \ce{H+} à l'azote) ; s'il n'y a plus d'hydrogène (amine tertiaire), on obtient directement l'ion ammonium quaternaire.
\item \textbf{Recommencer} tant qu'il reste un hydrogène sur l'azote et du dérivé halogéné disponible : primaire $\to$ secondaire $\to$ tertiaire $\to$ ammonium quaternaire.
\end{enumerate}
\end{methode}

\courssub{Exemple fil rouge : de l'ammoniac à l'ion tétraméthylammonium}

Suivons la cascade complète à partir de l'ammoniac \ce{NH3} et de l'iodométhane \ce{CH3-I}, en chauffant en tube scellé. Chaque étape est une substitution nucléophile.

\textbf{Étape 1 — formation de l'amine primaire :}
\[ \ce{NH3 + CH3-I -> CH3-NH2 + HI} \]
L'ammoniac donne la \textbf{méthanamine}.

\textbf{Étape 2 — formation de l'amine secondaire :}
\[ \ce{CH3-NH2 + CH3-I -> (CH3)2NH + HI} \]
La méthanamine donne la \textbf{diméthylamine}.

\textbf{Étape 3 — formation de l'amine tertiaire :}
\[ \ce{(CH3)2NH + CH3-I -> (CH3)3N + HI} \]
La diméthylamine donne la \textbf{triméthylamine}.

\textbf{Étape 4 — formation de l'ion ammonium quaternaire :}
\[ \ce{(CH3)3N + CH3-I -> (CH3)4N+ I-} \]
La triméthylamine donne l'\textbf{iodure de tétraméthylammonium}. C'est le terme de la cascade : l'azote, lié à quatre carbones, n'a plus de doublet disponible.

\begin{popfigure}
\begin{center}
\begin{tikzpicture}[node distance=0.35cm and 0.9cm,
  mol/.style={draw=popblue, thick, rounded corners=3pt, fill=popblueL, inner sep=5pt},
  cls/.style={font=\scriptsize\itshape, text=popmuted},
  fl/.style={-{Stealth}, thick, poporange}]
\node[mol] (n1) {\chemfig{NH_3}};
\node[cls, below=of n1] {ammoniac};
\node[mol, right=of n1] (n2) {\chemfig{CH_3-NH_2}};
\node[cls, below=of n2] {amine primaire};
\node[mol, right=of n2] (n3) {\chemfig{(CH_3)_2NH}};
\node[cls, below=of n3] {amine secondaire};
\node[mol, right=of n3] (n4) {\chemfig{(CH_3)_3N}};
\node[cls, below=of n4] {amine tertiaire};
\node[mol, right=of n4, draw=poppurple, fill=poppurpleL] (n5) {\chemfig{(CH_3)_4N^{+}\;I^{-}}};
\node[cls, below=of n5] {ammonium quaternaire};
\draw[fl] (n1) -- node[above, font=\scriptsize]{+\ce{CH3I}} node[below, font=\scriptsize]{-\ce{HI}} (n2);
\draw[fl] (n2) -- node[above, font=\scriptsize]{+\ce{CH3I}} node[below, font=\scriptsize]{-\ce{HI}} (n3);
\draw[fl] (n3) -- node[above, font=\scriptsize]{+\ce{CH3I}} node[below, font=\scriptsize]{-\ce{HI}} (n4);
\draw[fl] (n4) -- node[above, font=\scriptsize]{+\ce{CH3I}} (n5);
\end{tikzpicture}
\end{center}
\legende{Cascade des réactions d'Hofmann : alkylation successive de l'ammoniac par l'iodométhane, en chauffant en tube scellé.}
\end{popfigure}

\begin{attention}[Erreur fréquente : arrêter la réaction à la première étape]
Beaucoup d'élèves écrivent \ce{NH3 + CH3-I -> CH3-NH2 + HI} et s'arrêtent là. C'est \textbf{faux} si le dérivé halogéné est en excès ! L'amine primaire formée est encore nucléophile : l'alkylation \textbf{se poursuit} jusqu'à l'ion ammonium quaternaire. En réalité, on obtient toujours un \textbf{mélange} des quatre produits ; c'est le \textbf{réactif mis en excès} qui détermine le produit \emph{majoritaire} :
\begin{itemize}
\item excès d'\textbf{ammoniac} (ou d'amine) $\Rightarrow$ l'amine la moins substituée est majoritaire ;
\item excès de \textbf{dérivé halogéné} $\Rightarrow$ le sel d'ammonium quaternaire est majoritaire.
\end{itemize}
\end{attention}

\begin{center}
\begin{tabularx}{\linewidth}{|l|X|X|}
\hline
\rowcolor{popblueL}
\textbf{Réactif en excès} & \textbf{Produit majoritaire} & \textbf{Exemple avec \ce{NH3} + \ce{CH3-I}} \\
\hline
Ammoniac \ce{NH3} & Amine primaire & Méthanamine \ce{CH3-NH2} \\
\hline
Proportions voisines & Amines secondaire et tertiaire (mélange) & \ce{(CH3)2NH}, \ce{(CH3)3N} \\
\hline
Dérivé halogéné \ce{R-X} & Sel d'ammonium quaternaire & Iodure de tétraméthylammonium \ce{(CH3)4N+ I-} \\
\hline
\end{tabularx}
\end{center}

\courssub{Exemple chiffré détaillé : alkylation de la méthanamine par le bromoéthane}

\begin{exoresolu}[De la méthanamine au bromure d'éthyltriméthylammonium]
On chauffe en tube scellé la méthanamine \ce{CH3-NH2} avec un \textbf{excès} de bromoéthane \ce{CH3-CH2-Br}. Écrire les équations-bilan des réactions successives et nommer le produit final.
\tcblower
\textbf{Étape 1} : la méthanamine, nucléophile, attaque le carbone $\delta^+$ du bromoéthane :
\[ \ce{CH3-NH2 + CH3-CH2-Br -> CH3-NH-CH2-CH3 + HBr} \]
Produit : la \textbf{N-méthyléthanamine} (amine secondaire). Le HBr est neutralisé par l'excès d'amine sous forme de bromure de méthylammonium \ce{CH3-NH3+ Br-}.

\textbf{Étape 2} : l'amine secondaire est encore nucléophile :
\[ \ce{CH3-NH-CH2-CH3 + CH3-CH2-Br -> CH3-N(CH2-CH3)2 + HBr} \]
Produit : la \textbf{N-éthyl-N-méthyléthanamine}, c'est-à-dire \ce{CH3-N(C2H5)2} (amine tertiaire).

\textbf{Étape 3} : l'amine tertiaire fixe un dernier groupe éthyle, sans dégagement de HBr :
\[ \ce{CH3-N(C2H5)2 + CH3-CH2-Br -> [CH3-N(C2H5)3]+ Br-} \]
Produit final : le \textbf{bromure de triéthylméthylammonium}, un sel d'ammonium quaternaire. Comme le bromoéthane est en excès, c'est lui le produit \textbf{majoritaire}.
\end{exoresolu}

\begin{exercice}[À toi de jouer]
On fait réagir l'éthanamine \ce{C2H5-NH2} avec un excès d'iodométhane \ce{CH3-I}, à chaud, en tube scellé.
\begin{enumerate}
\item Écrire les équations-bilan des trois étapes d'alkylation.
\item Nommer l'amine secondaire et l'amine tertiaire obtenues.
\item Quel est le produit majoritaire final ? Justifier.
\end{enumerate}
\end{exercice}

\begin{corrige}[Exercice 1]
\textbf{1.} Les trois étapes successives :
\begin{align*}
\ce{C2H5-NH2 + CH3-I} &\ce{-> C2H5-NH-CH3 + HI}\\
\ce{C2H5-NH-CH3 + CH3-I} &\ce{-> C2H5-N(CH3)2 + HI}\\
\ce{C2H5-N(CH3)2 + CH3-I} &\ce{-> [C2H5-N(CH3)3]+ I-}
\end{align*}
\textbf{2.} Amine secondaire : la \textbf{N-méthyléthanamine}. Amine tertiaire : la \textbf{N,N-diméthyléthanamine}.

\textbf{3.} L'iodométhane étant en excès, l'alkylation va à son terme : le produit majoritaire est l'\textbf{iodure d'éthyltriméthylammonium} \ce{[C2H5-N(CH3)3]+ I-}, sel d'ammonium quaternaire.
\end{corrige}

\begin{savaistu}[Hofmann et les ammoniums quaternaires, des laboratoires aux hôpitaux]
C'est le chimiste allemand \textbf{August Wilhelm von Hofmann} qui découvrit en \textbf{1850} cette série d'alkylations successives, en chauffant de l'ammoniac avec des iodures d'alkyle dans des tubes scellés — d'où le nom de \emph{réactions d'Hofmann}. Cette méthode fut l'une des premières voies de synthèse des amines.

Aujourd'hui, les sels d'ammonium quaternaire sont partout : le \textbf{chlorure de benzalkonium} est le principe actif de nombreux \textbf{désinfectants} et antiseptiques utilisés dans les hôpitaux et les pharmacies du Cameroun (on le retrouve dans certaines solutions hydro-alcooliques et lingettes). D'autres ammoniums quaternaires servent d'\textbf{adoucissants} dans les lessives : leurs longues chaînes carbonées se fixent sur les fibres textiles et les rendent souples. Leur charge positive permanente, qui attire les membranes bactériennes chargées négativement, explique leur pouvoir bactéricide.
\end{savaistu}

\begin{aretenir}[L'essentiel de la section]
\begin{itemize}
\item Le doublet non liant de l'azote fait de l'amine un \cle{nucléophile} : elle attaque le carbone $\delta^+$ d'un dérivé halogéné \ce{R-X} (liaison \ce{C-X} polarisée).
\item La réaction est une \textbf{substitution nucléophile} (alkylation), lente, réalisée à chaud en tube scellé : \ce{R-NH2 + R$'$-X -> R-NH-R$'$ + HX}.
\item L'alkylation est \textbf{successive} (réactions d'Hofmann) : amine primaire $\to$ secondaire $\to$ tertiaire $\to$ \textbf{ion ammonium quaternaire} \ce{R4N+ X-}.
\item Le \textbf{réactif en excès} oriente le produit majoritaire : excès d'amine $\Rightarrow$ amine peu substituée ; excès de dérivé halogéné $\Rightarrow$ ammonium quaternaire.
\item Avec une amine tertiaire, il n'y a plus de \ce{HX} dégagé : \ce{R3N + R$'$-X -> R3R$'$N+ X-}.
\end{itemize}
\end{aretenir}

Le doublet de l'azote ne s'attaque pas qu'aux dérivés halogénés : il peut aussi attaquer le carbone, encore plus électrophile, d'un \textbf{chlorure d'acyle}. Cette réaction, plus rapide et non successive, ouvre la voie à la synthèse des \textbf{amides} — et donc de nombreux médicaments. C'est l'objet de la section suivante.

\coursec{Action sur les chlorures d'acyle : préparation des amides}

Dans la section précédente, nous avons vu que le doublet non liant de l'azote confère aux amines un \cle{caractère nucléophile} : l'amine attaque le carbone $\delta^+$ d'un dérivé halogéné \ce{R-X} et s'y fixe par \emph{alkylation} (réactions d'Hofmann). Mais le carbone porteur d'un halogène n'est pas le seul site électrophile que l'amine peut attaquer. Il existe une famille de dérivés d'acides carboxyliques — les \cle{chlorures d'acyle} \ce{R-COCl} et les \cle{anhydrides d'acide} \ce{R-CO-O-CO-R} — dont le carbone du groupe carbonyle est encore plus déficitaire en électrons. Lorsqu'une amine rencontre un tel composé, elle y fixe le groupe \emph{acyle} \ce{R-CO} : c'est une \cle{acylation}. Le produit obtenu, l'\cle{amide}, est l'une des fonctions les plus importantes de la chimie du vivant et du médicament : c'est elle qui relie les acides aminés dans les protéines. C'est d'ailleurs ainsi que l'industrie pharmaceutique fabrique le paracétamol que tu trouves dans toutes les pharmacies de Yaoundé ou de Douala.

\courssub{Pourquoi le carbone du chlorure d'acyle est-il électrophile ?}

Reprenons l'intuition avant la théorie. Dans un dérivé halogéné \ce{R-X}, le carbone est rendu $\delta^+$ par \emph{un seul} atome attracteur (l'halogène). Dans un chlorure d'acyle \ce{R-COCl}, le même carbone est lié à \emph{deux} atomes très électronégatifs : l'oxygène (double liaison \ce{C=O}) et le chlore (liaison \ce{C-Cl}). Ce carbone, dit \cle{carbonyle}, est donc fortement appauvri en électrons : c'est un \cle{site électrophile} particulièrement réactif.

\begin{propriete}[Électrophilie du carbone du chlorure d'acyle]
Dans un chlorure d'acyle \ce{R-COCl}, le carbone du groupe \ce{C=O} est lié à deux atomes attracteurs d'électrons (O et Cl). Il porte une charge partielle $\delta^+$ marquée : c'est un \cle{électrophile} puissant, facilement attaqué par le doublet non liant de l'azote d'une amine (nucléophile).
\end{propriete}

\begin{popfigure}
\begin{center}
\begin{tikzpicture}[baseline=(amine.base)]
\node (chlorure) {\chemfig{[:-30]R-C(=[2]O)(-[6]Cl)}};
\node[right=1.5cm of chlorure] (amine) {\chemfig{[:-30]R'-N(-[2]H)(-[6]H)}};
\draw[popblue, thick, ->] (chlorure.east) -- ++(0.5,0) node[midway, above] {\small Attaque};
\draw[poporange, thick, <-] (amine.west) -- ++(-0.5,0);
\node[popblue, font=\small\bfseries, below=2mm of chlorure] {Site électrophile : C $\delta^+$};
\node[poporange, font=\small\bfseries, below=2mm of amine] {Site nucléophile : doublet sur N};
\end{tikzpicture}
\end{center}
\legende{Les deux partenaires de la réaction : le chlorure d'acyle (électrophile) et l'amine (nucléophile). Le doublet non liant de l'azote attaque le carbone $\delta^+$ du carbonyle.}
\end{popfigure}

\courssub{La réaction : formation d'un amide}

\begin{definition}[Amide et liaison amide]
Un \cle{amide} est un composé organique de formule générale \ce{R-CO-NR$'$R$''$}, où R', R'' sont des atomes d'hydrogène ou des groupes alkyles (ou aryles). La liaison \ce{C-N} entre le groupe carbonyle et l'azote est appelée \cle{liaison amide} ; lorsqu'elle relie des acides aminés dans les protéines, on l'appelle \cle{liaison peptidique}. Selon le nombre de groupes carbonés portés par l'azote, on distingue les amides \emph{non substitués} \ce{R-CO-NH2}, \emph{N-substitués} \ce{R-CO-NH-R$'$} et \emph{N,N-disubstitués} \ce{R-CO-NR$'$2}.
\end{definition}

\textbf{Avec une amine primaire.} L'amine primaire \ce{R$'$-NH2} attaque le carbone du chlorure d'acyle ; la liaison \ce{C-Cl} se rompt et un atome d'hydrogène de l'azote s'élimine avec le chlore sous forme de \ce{HCl}. On obtient un amide \emph{N-substitué} :

\[ \ce{R-COCl + R$'$-NH2 -> R-CO-NH-R$'$ + HCl} \]

\textbf{Avec une amine secondaire.} Le mécanisme est identique ; l'azote portant deux groupes alkyles, on obtient un amide \emph{N,N-disubstitué} :

\[ \ce{R-COCl + R$'$2NH -> R-CO-NR$'$2 + HCl} \]

\textbf{Avec l'ammoniac.} L'ammoniac \ce{NH3} possède trois hydrogènes sur l'azote : il réagit de la même façon et donne un amide \emph{non substitué} :

\[ \ce{R-COCl + NH3 -> R-CO-NH2 + HCl} \]

\textbf{Et les amines tertiaires ?} Une amine tertiaire \ce{R3N} ne possède \emph{aucun} hydrogène sur l'azote : l'élimination de \ce{HCl} est impossible, et la réaction ne conduit pas à un amide.

\begin{attention}[Piège classique du Bac : l'amine tertiaire]
Ne fais \textbf{jamais} réagir une amine tertiaire \ce{R3N} avec un chlorure d'acyle pour « préparer un amide » : c'est \textbf{impossible}. La formation de l'amide exige le départ de \ce{HCl}, donc la présence d'\emph{au moins un atome d'hydrogène sur l'azote}. Avant d'écrire l'équation, compte les H portés par N : amine primaire (2 H, réagit), secondaire (1 H, réagit), tertiaire (0 H, ne réagit pas). Cette question est un classique des sujets de Baccalauréat C et D.
\end{attention}

\begin{methode}[Écrire l'équation amine + chlorure d'acyle]
Pour écrire sans erreur l'équation-bilan de la réaction :
\begin{enumerate}
\item \textbf{Identifier} le chlorure d'acyle \ce{R-COCl} et l'amine (primaire ou secondaire) ; vérifier qu'il reste \emph{au moins un H} sur l'azote.
\item \textbf{Rompre} mentalement la liaison \ce{C-Cl} du chlorure d'acyle.
\item \textbf{Fixer} le groupe acyle \ce{R-CO} sur l'azote de l'amine, à la place d'un de ses hydrogènes.
\item \textbf{Éliminer} \ce{HCl} (le H de l'azote + le Cl du chlorure).
\item \textbf{Vérifier} la conservation des atomes et nommer l'amide obtenu : chaîne du groupe \ce{R-CO} (suffixe \emph{-amide}) précédée de \emph{N-} ou \emph{N,N-} suivi du nom du ou des groupes fixés sur l'azote.
\end{enumerate}
\end{methode}

\begin{exemplebox}[Préparation du N-éthyléthanamide]
On fait réagir le chlorure d'éthanoyle sur l'éthanamine (amine primaire) :
\[ \ce{CH3-COCl + C2H5-NH2 -> CH3-CO-NH-C2H5 + HCl} \]
\begin{center}
\chemfig{CH_3-C(=[1]O)-[7]Cl} \; \ce{+} \; \chemfig{CH_3-CH_2-NH_2} \; \ce{->} \; \chemfig{CH_3-C(=[1]O)-[7]NH-CH_2-CH_3} \; \ce{+} \; \ce{HCl}
\end{center}
Le produit est le \textbf{N-éthyléthanamide} : le préfixe \emph{N-éthyl} indique que le groupe éthyle est fixé sur l'azote, et \emph{éthanamide} désigne la chaîne \ce{CH3-CO-NH-} à deux atomes de carbone.
\end{exemplebox}

\courssub{Caractéristiques de la réaction et variante avec un anhydride}

\begin{propriete}[Caractéristiques de l'acylation d'une amine]
La réaction d'une amine avec un chlorure d'acyle est :
\begin{itemize}
\item \textbf{rapide} et \textbf{totale} (ce n'est pas un équilibre, contrairement à l'estérification) ;
\item \textbf{exothermique} : on opère souvent à froid, en ajoutant le chlorure d'acyle goutte à goutte ;
\item productrice de \ce{HCl} : on utilise fréquemment un \textbf{excès d'amine} (ou on ajoute une base comme \ce{NaOH}) pour neutraliser l'acide chlorhydrique formé, qui sinon transformerait l'amine en ion alkylammonium \ce{R-NH3+}, privée de son doublet nucléophile.
\end{itemize}
\end{propriete}

On peut remplacer le chlorure d'acyle par un \cle{anhydride d'acide} \ce{R-CO-O-CO-R}. La réaction est analogue mais \textbf{plus douce} (moins violente, plus facile à contrôler au laboratoire) ; elle produit un acide carboxylique au lieu de \ce{HCl} :

\[ \ce{R-CO-O-CO-R + R$'$-NH2 -> R-CO-NH-R$'$ + R-COOH} \]

\begin{exemplebox}[Avec l'anhydride éthanoïque]
\[ \ce{(CH3-CO)2O + C2H5-NH2 -> CH3-CO-NH-C2H5 + CH3-COOH} \]
On obtient le même N-éthyléthanamide qu'avec le chlorure d'éthanoyle, accompagné d'acide éthanoïque. Industriellement, c'est cette voie (anhydride éthanoïque) qui est utilisée pour fabriquer le paracétamol à partir du para-aminophénol.
\end{exemplebox}

\courssub{La liaison amide : des protéines aux médicaments}

\begin{popfigure}
\begin{center}
\chemfig{[:-30]H_3C-@{C}C(=[2]O)-[6]@{N}N(-[2]H)-[6]**6((-HO)-=-=-=)}
\chemmove{
\draw[poporange, very thick, rounded corners] (C.north west) rectangle (N.south east);
\node[poporange, font=\small\bfseries, above=2pt] at ($(C)!0.5!(N)$) {groupe amide \ce{-CO-NH-}};
}
\end{center}
\legende{Molécule de paracétamol (N-(4-hydroxyphényl)éthanamide). Le groupe amide, encadré en orange, résulte de l'acylation du para-aminophénol par l'anhydride éthanoïque.}
\end{popfigure}

La liaison amide \ce{-CO-N-} est partout autour de toi. Dans les \textbf{protéines} (viande, poisson, œufs, niébé), elle relie les acides aminés les uns aux autres : on parle alors de \cle{liaison peptidique}. Dans les \textbf{médicaments}, on la retrouve dans le paracétamol, dans les pénicillines (leur cycle $\beta$-lactame est un amide cyclique) et dans de nombreux autres principes actifs. Sa stabilité explique qu'elle soit le ciment de la chimie du vivant.

\begin{savaistu}[Le nylon, un polyamide dans nos vêtements]
Le \textbf{nylon 6,6} est un \emph{polyamide} obtenu par polycondensation entre l'hexane-1,6-diamine \ce{H2N-(CH2)6-NH2} et l'acide hexane-1,6-dioïque (ou son dérivé acylé) : chaque molécule de diamine réagit aux deux bouts, formant de longues chaînes reliées par des liaisons amide. La remarquable solidité des fibres de nylon — présentes dans nos vêtements techniques, nos cordes et nos filets de pêche — vient des \textbf{liaisons hydrogène} qui s'établissent entre les groupes \ce{N-H} et \ce{C=O} des chaînes voisines, exactement comme entre les brins des protéines. La chimie des amides est donc aussi celle des textiles modernes vendus dans nos marchés.
\end{savaistu}

\courssub{Bilan : deux visages du caractère nucléophile des amines}

Nous avons rencontré deux réactions où l'amine joue le rôle de nucléophile. Il est essentiel de ne pas les confondre : l'une est une \emph{alkylation} (fixation d'un groupe alkyle), l'autre une \emph{acylation} (fixation d'un groupe acyle \ce{R-CO}).

\begin{center}
\begin{tabularx}{\linewidth}{|l|X|X|}
\hline
\rowcolor{popblueL} \textbf{Critère} & \textbf{Amine + dérivé halogéné \ce{R-X}} & \textbf{Amine + chlorure d'acyle \ce{R-COCl}} \\
\hline
Nom de la réaction & Alkylation (réactions d'Hofmann) & Acylation \\
\hline
Site électrophile attaqué & Carbone $\delta^+$ porteur de l'halogène & Carbone $\delta^+$ du groupe carbonyle \ce{C=O} \\
\hline
Groupe fixé sur l'azote & Groupe alkyle \ce{R} & Groupe acyle \ce{R-CO} \\
\hline
Produit obtenu & Amine de classe supérieure, puis ion ammonium quaternaire \ce{R4N+} & Amide \ce{R-CO-NR$'$R$''$} \\
\hline
Sous-produit & \ce{HX} (puis \ce{R-NH3+ X-}) & \ce{HCl} \\
\hline
Amine tertiaire & Réagit (donne l'ion ammonium quaternaire) & Ne réagit \textbf{pas} (pas de H sur N) \\
\hline
Caractéristiques & Réactions successives, mélange de produits & Rapide, totale, exothermique \\
\hline
\end{tabularx}
\end{center}

\begin{exoresolu}[Identifier le réactif et le produit]
On désire préparer le N-méthylpropanamide \ce{C2H5-CO-NH-CH3}.
\begin{enumerate}
\item Écrire la formule semi-développée du chlorure d'acyle et de l'amine nécessaires.
\item Écrire l'équation-bilan de la réaction et préciser ses caractéristiques.
\item La triméthylamine \ce{N(CH3)3} conviendrait-elle à la place de l'amine choisie ? Justifier.
\end{enumerate}
\tcblower
\textbf{1.} Le nom \emph{N-méthylpropanamide} indique un groupe acyle à 3 carbones (\ce{C2H5-CO-}) et un méthyle sur l'azote. Il faut donc le \textbf{chlorure de propanoyle} \ce{C2H5-COCl} et la \textbf{méthanamine} (méthylamine) \ce{CH3-NH2}, amine primaire.

\textbf{2.} Équation-bilan :
\[ \ce{C2H5-COCl + CH3-NH2 -> C2H5-CO-NH-CH3 + HCl} \]
La réaction est \textbf{rapide, totale et exothermique} ; on utilise en pratique un excès de méthanamine pour neutraliser le \ce{HCl} formé.

\textbf{3.} \textbf{Non.} La triméthylamine est une amine \emph{tertiaire} : son azote ne porte aucun atome d'hydrogène. L'élimination de \ce{HCl}, indispensable à la formation de la liaison amide, est donc impossible : aucun amide ne se forme.
\end{exoresolu}

\begin{exercice}[À toi de jouer]
\begin{enumerate}
\item Écrire l'équation-bilan de la réaction entre le chlorure d'éthanoyle et la diéthylamine \ce{(C2H5)2NH}. Nommer l'amide obtenu et préciser s'il est substitué.
\item Écrire l'équation-bilan de la réaction entre l'anhydride éthanoïque et l'ammoniac. Quel est le nom de l'amide formé ?
\item Expliquer pourquoi on ajoute parfois une solution de soude lors de la préparation d'un amide à partir d'un chlorure d'acyle.
\end{enumerate}
\end{exercice}

\begin{corrige}[Exercice]
\textbf{1.} \ce{CH3-COCl + (C2H5)2NH -> CH3-CO-N(C2H5)2 + HCl}. L'amide obtenu est le \textbf{N,N-diéthyléthanamide} : c'est un amide N,N-disubstitué (l'azote porte deux groupes éthyle).

\textbf{2.} \ce{(CH3-CO)2O + NH3 -> CH3-CO-NH2 + CH3-COOH}. L'amide formé est l'\textbf{éthanamide} (amide non substitué), accompagné d'acide éthanoïque.

\textbf{3.} La réaction dégage \ce{HCl}, un acide fort qui protonerait l'amine restante en ion alkylammonium \ce{R-NH3+}, laquelle n'a plus de doublet disponible et cesse d'être nucléophile. La soude neutralise \ce{HCl} (\ce{HCl + NaOH -> NaCl + H2O}) et maintient l'amine sous sa forme basique réactive, ce qui améliore le rendement en amide.
\end{corrige}

\begin{aretenir}[L'essentiel de la section]
\begin{itemize}
\item Le carbone du groupe \ce{C=O} d'un chlorure d'acyle \ce{R-COCl} est un \cle{électrophile} puissant (lié à O et Cl, tous deux attracteurs).
\item Amine primaire ou secondaire + chlorure d'acyle $\rightarrow$ \cle{amide} + \ce{HCl} : réaction \textbf{rapide, totale, exothermique} ; avec l'ammoniac, on obtient l'amide non substitué \ce{R-CO-NH2}.
\item Une \textbf{amine tertiaire ne réagit pas} : pas d'hydrogène sur l'azote, donc pas d'élimination de \ce{HCl}.
\item Avec un \textbf{anhydride d'acide}, la réaction est plus douce : \ce{(R-CO)2O + R$'$-NH2 -> R-CO-NH-R$'$ + R-COOH}.
\item La \cle{liaison amide} \ce{-CO-N-} (ou liaison peptidique) est le motif des protéines, du paracétamol, des pénicillines et du nylon.
\item Alkylation (\ce{R-X}) $\neq$ acylation (\ce{R-COCl}) : dans les deux cas l'amine est nucléophile, mais le produit est une amine de classe supérieure dans le premier cas, un amide dans le second.
\end{itemize}
\end{aretenir}

\newpage

\coursec{Travaux pratiques}

\courssub{Objectifs du TP}

À l'issue de cette séance de travaux pratiques, tu dois être capable de :

\begin{itemize}
\item mettre en évidence expérimentalement le \cle{caractère basique} d'une amine en solution aqueuse à l'aide d'\cle{indicateurs colorés} (BBT, phénolphtaléine) et du papier pH ;
\item observer l'action d'une amine sur les ions \ce{Cu^2+} (formation d'un précipité d'hydroxyde de cuivre(II)) ;
\item comparer le comportement d'une amine à celui de l'\cle{ammoniac} et à celui de l'eau distillée (témoin négatif) ;
\item écrire et interpréter l'équation-bilan de la réaction de l'eau sur une amine, en faisant intervenir le \cle{doublet non liant} de l'atome d'azote.
\end{itemize}

\courssub{Matériel et produits}

\begin{center}
\begin{tabularx}{\linewidth}{|l|X|}
\hline
\rowcolor{popblue!60} \textbf{Matériel} & \textbf{Produits} \\
\hline
8 tubes à essai + portoir & Solution aqueuse d'éthylamine \ce{C2H5NH2} à \SI{0,1}{mol.L^{-1}} \\
\hline
2 pipettes jaugées de \SI{2}{mL} (ou pipettes compte-gouttes graduées) & Solution d'ammoniac diluée \ce{NH3} à \SI{0,1}{mol.L^{-1}} (témoin positif) \\
\hline
Pipettes compte-gouttes & Eau distillée (témoin négatif) \\
\hline
Papier pH + baguettes en verre & Bleu de bromothymol (BBT) \\
\hline
Spatule, bouchons, étiquettes & Phénolphtaléine (solution alcoolique) \\
\hline
Lunettes de protection, gants, blouse & Solution de sulfate de cuivre(II) \ce{CuSO4} à \SI{0,1}{mol.L^{-1}} \\
\hline
\end{tabularx}
\end{center}

\begin{attention}[Si un produit manque au laboratoire]
L'éthylamine est parfois difficile à se procurer dans les lycées. On peut la remplacer par une solution de \cle{méthylamine} ou, à défaut, utiliser uniquement la solution d'\cle{ammoniac} (l'ammoniac possède le même doublet non liant sur l'azote et se comporte comme une base : les conclusions restent valables). Si le BBT manque, on peut utiliser du jus de \cle{hibiscus} (bissap) fraîchement préparé : il vire au vert-bleu en milieu basique. Le papier pH peut être remplacé par du papier tournesol.
\end{attention}

\courssub{Sécurité}

\begin{attention}[Consignes de sécurité — à lire avant toute manipulation]
\begin{itemize}
\item \textbf{Éthylamine} : liquide \emph{inflammable} (pictogramme : flamme) et \emph{corrosif} (pictogramme : corrosion). Travailler loin de toute flamme, sous la hotte ou près d'une fenêtre ouverte. Ne jamais respirer les vapeurs directement.
\item \textbf{Ammoniac} : vapeurs \emph{irritantes} pour les yeux et les voies respiratoires (pictogramme : danger). Manipuler le flacon bouché, ouvrir brièvement.
\item \textbf{Sulfate de cuivre(II)} : \emph{nocif} et \emph{dangereux pour l'environnement aquatique} (pictogrammes : point d'exclamation, environnement). Ne pas jeter à l'évier : récupérer les déchets dans un bidon prévu à cet effet.
\item \textbf{Gestes obligatoires} : porter blouse, lunettes et gants ; ne jamais pipeter à la bouche ; en cas de contact avec la peau, rincer abondamment à l'eau claire et prévenir l'enseignant.
\end{itemize}
\end{attention}

\courssub{Schéma du dispositif}

\begin{popfigure}
\begin{center}
\begin{tikzpicture}[scale=0.9, every node/.style={font=\small}]
% Tube 1 : BBT
\draw[thick, popdark] (0,0) -- (0,2.6);
\draw[thick, popdark] (1,0) -- (1,2.6);
\draw[thick, popdark] (0,0) arc (180:360:0.5);
\fill[popblue!50] (0.02,0.15) arc (180:360:0.48) -- (0.98,1.3) -- (0.02,1.3) -- cycle;
\draw[popdark] (0,1.3) -- (1,1.3);
\node[below] at (0.5,-0.35) {Tube 1 : BBT};
\node[popblue!80!black] at (0.5,0.7) {bleu};

% Tube 2 : phénolphtaléine
\begin{scope}[xshift=2.2cm]
\draw[thick, popdark] (0,0) -- (0,2.6);
\draw[thick, popdark] (1,0) -- (1,2.6);
\draw[thick, popdark] (0,0) arc (180:360:0.5);
\fill[poppink!60] (0.02,0.15) arc (180:360:0.48) -- (0.98,1.3) -- (0.02,1.3) -- cycle;
\draw[popdark] (0,1.3) -- (1,1.3);
\node[below] at (0.5,-0.35) {Tube 2 : phénolphtaléine};
\node[poppink!80!black] at (0.5,0.7) {rose};
\end{scope}

% Tube 3 : papier pH
\begin{scope}[xshift=4.4cm]
\draw[thick, popdark] (0,0) -- (0,2.6);
\draw[thick, popdark] (1,0) -- (1,2.6);
\draw[thick, popdark] (0,0) arc (180:360:0.5);
\fill[popblue!20] (0.02,0.15) arc (180:360:0.48) -- (0.98,1.3) -- (0.02,1.3) -- cycle;
\draw[popdark] (0,1.3) -- (1,1.3);
\draw[thick, popgold!80!black, fill=popgreen!50] (0.25,1.3) rectangle (0.45,2.4);
\node[below] at (0.5,-0.35) {Tube 3 : papier pH};
\node at (0.5,0.7) {pH $>7$};
\end{scope}

% Tube 4 : CuSO4
\begin{scope}[xshift=6.6cm]
\draw[thick, popdark] (0,0) -- (0,2.6);
\draw[thick, popdark] (1,0) -- (1,2.6);
\draw[thick, popdark] (0,0) arc (180:360:0.5);
\fill[popblue!70] (0.02,0.15) arc (180:360:0.48) -- (0.98,0.7) -- (0.02,0.7) -- cycle;
\fill[popblue!25] (0.02,0.7) -- (0.98,0.7) -- (0.98,1.3) -- (0.02,1.3) -- cycle;
\draw[popdark] (0,1.3) -- (1,1.3);
% pipette compte-gouttes
\draw[thick, popmuted] (0.5,3.4) -- (0.5,1.8);
\draw[thick, popmuted] (0.4,3.4) -- (0.6,3.4);
\fill[popblue!80] (0.5,1.55) circle (0.06);
\node[right, popmuted] at (0.65,3.0) {\ce{CuSO4} goutte à goutte};
\node[below] at (0.5,-0.35) {Tube 4 : \ce{CuSO4}};
\node[popblue!90!black] at (0.5,0.4) {précipité};
\end{scope}

% Légende solution
\node[align=center] at (3.8,3.2) {\textbf{Solution d'éthylamine \ce{C2H5NH2} (\SI{2}{mL} par tube)}};
\end{tikzpicture}
\end{center}
\legende{Dispositif expérimental : quatre tubes contenant chacun \SI{2}{mL} de solution d'éthylamine, testés respectivement au BBT, à la phénolphtaléine, au papier pH et au sulfate de cuivre(II).}
\end{popfigure}

\courssub{Protocole}

\begin{methode}[Mode opératoire]
\begin{enumerate}
\item \textbf{Préparation des tubes.} Numéroter 4 tubes à essai de 1 à 4. Verser dans chacun \SI{2}{mL} de solution aqueuse d'éthylamine à \SI{0,1}{mol.L^{-1}} à l'aide de la pipette jaugée.
\item \textbf{Tube 1 — BBT.} Ajouter 2 gouttes de solution de bleu de bromothymol. Agiter doucement. Noter la teinte obtenue.
\item \textbf{Tube 2 — Phénolphtaléine.} Ajouter 2 gouttes de solution de phénolphtaléine. Agiter. Noter la teinte.
\item \textbf{Tube 3 — Mesure du pH.} À l'aide d'une baguette en verre propre, déposer une goutte de solution sur un morceau de papier pH. Comparer la couleur à l'échelle étalon et relever la valeur du pH.
\item \textbf{Tube 4 — Ions cuivre(II).} Ajouter \emph{goutte à goutte}, en agitant, la solution de sulfate de cuivre(II). Observer attentivement dès les premières gouttes.
\item \textbf{Témoin positif.} Recommencer les quatre tests avec une solution d'ammoniac diluée à \SI{0,1}{mol.L^{-1}} (tubes 1' à 4').
\item \textbf{Témoin négatif.} Recommencer les quatre tests avec de l'eau distillée (tubes 1'' à 4'').
\item \textbf{Consignation.} Reporter toutes les observations dans le tableau de résultats.
\end{enumerate}
\end{methode}

\courssub{Observations et résultats attendus}

\begin{center}
\begin{tabularx}{\linewidth}{|l|X|X|X|}
\hline
\rowcolor{popgreen!60} \textbf{Test} & \textbf{Éthylamine} & \textbf{Ammoniac} & \textbf{Eau distillée} \\
\hline
BBT & Virage au bleu & Virage au bleu & Reste jaune-vert \\
\hline
Phénolphtaléine & Virage au rose & Virage au rose & Reste incolore \\
\hline
Papier pH & pH $\approx 11$ à $12$ & pH $\approx 11$ & pH $\approx 7$ \\
\hline
\ce{CuSO4} goutte à goutte & Précipité bleu de \ce{Cu(OH)2} & Précipité bleu, puis solution bleu foncé en excès d'ammoniac & Solution bleue homogène, sans précipité \\
\hline
\end{tabularx}
\end{center}

\begin{attention}[Piège fréquent au tube 4]
Avec l'ammoniac, si l'on ajoute un \emph{excès} de réactif, le précipité bleu de \ce{Cu(OH)2} se \emph{redissout} en donnant une solution bleu outremer très foncé (formation de l'ion complexe \ce{[Cu(NH3)4]^2+}). Pour rester dans le cadre du programme, ajouter le sulfate de cuivre(II) \emph{goutte à goutte} et s'arrêter dès l'apparition du précipité bleu.
\end{attention}

\courssub{Exploitation}

\begin{exoresolu}[Questions guidées d'exploitation]
\textbf{Énoncé.}
\begin{enumerate}
\item Que signifie le virage du BBT au bleu et celui de la phénolphtaléine au rose ?
\item Écrire l'équation-bilan de la réaction de l'eau sur l'éthylamine. Identifier le couple acide/base mis en jeu.
\item Expliquer, à l'aide du doublet non liant de l'azote, pourquoi l'éthylamine se comporte comme une base.
\item Écrire l'équation de la réaction de précipitation observée au tube 4.
\item Comparer les résultats obtenus avec l'éthylamine et avec l'ammoniac. Conclure.
\end{enumerate}
\tcblower
\textbf{Solution détaillée.}
\begin{enumerate}
\item Le BBT est bleu en milieu basique (pH $> 7{,}6$) et la phénolphtaléine est rose pour pH $> 8{,}2$. Le double virage prouve que la solution d'éthylamine est \cle{basique} : elle contient des ions hydroxyde \ce{HO-}.
\item L'éthylamine capte un proton \ce{H+} d'une molécule d'eau selon une réaction \emph{partielle} (d'où la double flèche) :
\[ \ce{C2H5NH2 + H2O <=> C2H5NH3+ + HO-} \]
Le couple mis en jeu est \ce{C2H5NH3+}/\ce{C2H5NH2} (ion éthylammonium / éthylamine). L'eau joue ici le rôle d'\emph{acide} (donneur de proton).
\item L'atome d'azote de l'amine porte un \cle{doublet non liant}. Ce doublet peut former une liaison de covalence dative avec un proton \ce{H+} : c'est exactement la définition d'une \emph{base} selon Brønsted (capteur de proton). La libération d'ions \ce{HO-} rend la solution basique, ce qui explique le pH $\approx 12$ et le virage des indicateurs.
\item Les ions \ce{HO-} formés réagissent avec les ions cuivre(II) pour donner un précipité bleu d'hydroxyde de cuivre(II) :
\[ \ce{Cu^2+ + 2 HO- -> Cu(OH)2\v{s}} \]
\item L'ammoniac donne des résultats quasi identiques (BBT bleu, phénolphtaléine rose, pH élevé, précipité bleu), car il possède lui aussi un doublet non liant sur l'azote : \ce{NH3 + H2O <=> NH4+ + HO-}. L'eau distillée, elle, ne fait virer aucun indicateur (pH $\approx 7$) : c'est le témoin qui prouve que les observations sont bien dues à l'amine. \textbf{Conclusion :} les amines, comme l'ammoniac, sont des \emph{bases} en solution aqueuse.
\end{enumerate}
\end{exoresolu}

\courssub{Conclusion}

Cette séance a permis de vérifier expérimentalement le \cle{caractère basique des amines}. En solution aqueuse, l'éthylamine réagit partiellement avec l'eau selon :
\[ \ce{C2H5NH2 + H2O <=> C2H5NH3+ + HO-} \]
La formation d'ions hydroxyde \ce{HO-} est confirmée par trois faits expérimentaux concordants : le virage du BBT au bleu, le virage de la phénolphtaléine au rose, et un pH supérieur à 7. Ces ions \ce{HO-} précipitent ensuite les ions \ce{Cu^2+} sous forme d'hydroxyde de cuivre(II) bleu. L'origine de ce comportement est le \cle{doublet non liant} porté par l'atome d'azote, capable de capter un proton : c'est la même raison qui rend l'ammoniac basique. Retiens l'idée directrice de toute cette leçon : \textbf{le doublet non liant de l'azote gouverne toute la chimie des amines} — leur basicité comme leur caractère nucléophile, que tu exploiteras dans les réactions d'Hofmann et la préparation des amides.

\begin{savaistu}[Les amines autour de nous]
Les amines ne sont pas seulement des réactifs de laboratoire ! On les retrouve dans la vie quotidienne au Cameroun :
\begin{itemize}
\item \textbf{Odeurs caractéristiques :} L'odeur forte du poisson pas tout à fait frais (triméthylamine), l'odeur de la putréfaction (putrescine, cadavérine) sont dues à des amines.
\item \textbf{Boissons traditionnelles :} Lors de la fermentation du \emph{bili-bili} ou du vin de palme, des levures et bactéries produisent des amines biogéniques (histamine, tyramine) qui peuvent donner des maux de tête en cas d'excès.
\item \textbf{Pharmacopée :} De nombreux alcaloïdes (caféine, quinine, morphine) sont des amines tertiaires ou hétérocycliques à propriétés médicinales.
\item \textbf{Industrie :} La fabrication de savons (savonneries artisanales ou industrielles comme \emph{Savonnerie du Cameroun}) utilise parfois des amines comme agents tensioactifs ou catalyseurs.
\end{itemize}
\end{savaistu}

\newpage

\coursec{Méthodes et exercices résolus}

\coursec{Les amines}

\courssub{De l'ammoniac aux amines : structure, définition et classes}

\begin{definition}[Structure de l'ammoniac et passage aux amines]
La molécule d'ammoniac \ce{NH3} adopte une géométrie \textbf{pyramidale trigonale} (angle H-N-H \approx \SI{107}{\degree}). L'atome d'azote est hybridé \cle{$sp^3$} : trois orbitaux forment les liaisons \ce{N-H}, le quatrième contient un \cle{doublet non liant}. Ce doublet est à l'origine du caractère basique et nucléophile de l'ammoniac.
Une \cle{amine} est un dérivé de l'ammoniac dans lequel un, deux ou trois atomes d'hydrogène ont été remplacés par des groupes alkyles (ou aryles). La formule générale des amines est \ce{R-NH2}, \ce{R2NH} ou \ce{R3N} (R = alkyle ou aryle, non tous H).
\end{definition}

\begin{definition}[Classes d'amines et ion ammonium quaternaire]
On distingue trois classes selon le nombre de groupes carbonés liés à l'atome d'azote :
\begin{itemize}
\item \textbf{Amine primaire (1\degre)} : un groupe alkyle/aryle \ce{R-NH2} (ex. : \ce{CH3-NH2} méthanamine).
\item \textbf{Amine secondaire (2\degre)} : deux groupes \ce{R-NH-R$'$} (ex. : \ce{CH3-NH-CH3} diméthylamine).
\item \textbf{Amine tertiaire (3\degre)} : trois groupes \ce{R-NR$'$R$''$} (ex. : \ce{(CH3)3N} triméthylamine).
\end{itemize}
Si l'azote porte \textbf{quatre} groupes alkyles, il porte une charge positive : c'est un \cle{ion ammonium quaternaire} \ce{[R4N]+} (ex. : \ce{[(CH3)4N]+} tétraéthylammonium). Il n'est plus une amine mais un sel.
\end{definition}

\begin{popfigure}
\begin{tikzpicture}[scale=0.9, every node/.style={transform shape}]
% Ammoniac
\begin{scope}[xshift=-5cm]
\chemfig{[:30]H-N(-[2]H)-[6]H}
\node[below=0.5cm] at (0,-1.5) {\textbf{Ammoniac \ce{NH3}}};
\node[below=1.2cm] at (0,-1.5) {Pyramidale, doublet non liant};
\end{scope}
% Amine primaire
\begin{scope}[xshift=-1.5cm]
\chemfig{[:30]CH_3-N(-[2]H)-[6]H}
\node[below=0.5cm] at (0,-1.5) {\textbf{Primaire } $\mathrm{R{-}NH_2}$};
\node[below=1.2cm] at (0,-1.5) {1 groupe R, 2 liaisons N-H};
\end{scope}
% Amine secondaire
\begin{scope}[xshift=2cm]
\chemfig{[:30]CH_3-N(-[2]CH_3)-[6]H}
\node[below=0.5cm] at (0,-1.5) {\textbf{Secondaire } $\mathrm{R_2NH}$};
\node[below=1.2cm] at (0,-1.5) {2 groupes R, 1 liaison N-H};
\end{scope}
% Amine tertiaire
\begin{scope}[xshift=5.5cm]
\chemfig{[:30]CH_3-N(-[2]CH_3)-[6]CH_3}
\node[below=0.5cm] at (0,-1.5) {\textbf{Tertiaire } $\mathrm{R_3N}$};
\node[below=1.2cm] at (0,-1.5) {3 groupes R, 0 liaison N-H};
\end{scope}
% Quaternaire
\begin{scope}[xshift=9cm]
\chemfig{[:30]CH_3-N^{+}(-[2]CH_3)(-[6]CH_3)-CH_3}
\node[below=0.5cm] at (0,-1.5) {\textbf{Quaternaire } $[\mathrm{R_4N}]^{+}$};
\node[below=1.2cm] at (0,-1.5) {4 groupes R, charge +};
\end{scope}
\end{tikzpicture}
\legende{De l'ammoniac aux amines : substitution progressive des H par des groupes alkyles. L'azote reste hybridé sp$^3$ (pyramidal).}
\end{popfigure}

\begin{methode}[Reconnaître la classe d'une amine]
\begin{enumerate}
\item Repérer l'atome d'\cle{azote} dans la formule : c'est lui qui porte la fonction amine.
\item \cle{Compter le nombre de groupes alkyles} (ou aryles) directement liés à l'azote :
\begin{itemize}
\item 1 groupe $\Rightarrow$ amine \cle{primaire} \ce{R-NH2} ;
\item 2 groupes $\Rightarrow$ amine \cle{secondaire} \ce{R-NH-R$'$} ;
\item 3 groupes $\Rightarrow$ amine \cle{tertiaire} \ce{R-NR$'$R$''$}.
\end{itemize}
\item \textbf{Réflexe de sécurité} : compter les \emph{carbones} liés à l'azote, jamais les hydrogènes. Un atome d'azote lié à 4 carbones porte une charge $+$ : c'est un \cle{ion ammonium quaternaire} \ce{[R4N]+}.
\item Pour écrire une formule semi-développée à partir d'un nom : placer l'azote, y accrocher les groupes cités dans le nom, puis compléter avec les hydrogènes pour que l'azote ait \textbf{3 liaisons} (et un doublet non liant).
\end{enumerate}
\end{methode}

\begin{exoresolu}[Classes d'amines (type Bac)]
On considère les composés suivants :
\ce{CH3-CH2-NH2} (A) ; \ce{CH3-NH-CH2-CH3} (B) ; \ce{(CH3)3N} (C) ; \ce{CH3-CH(NH2)-CH3} (D) ; \ce{[(CH3)4N]+} (E).
\begin{enumerate}
\item Écrire la formule semi-développée de chacun des composés C et E.
\item Indiquer la classe de chaque composé.
\item Parmi ces composés, lesquels sont des isomères ?
\end{enumerate}
\tcblower
\textbf{1. Formules semi-développées.}

Le composé C est la \emph{triméthylamine} : l'azote porte trois groupes méthyle :
\[ \chemfig{CH_3-N(-[2]CH_3)-CH_3} \qquad \text{soit} \qquad \chemfig{H_3C-N(-[2]CH_3)-CH_3} \]
Le composé E est l'ion \emph{tétraméthylammonium} : l'azote porte quatre groupes méthyle et une charge positive :
\[ \chemfig{CH_3-N^{+}(-[2]CH_3)(-[6]CH_3)-CH_3} \]

\textbf{2. Classe de chaque composé.} On compte les groupes alkyles liés à l'azote :
\begin{center}
\begin{tabularx}{\linewidth}{|l|X|l|}
\hline
\rowcolor{popblueL} \textbf{Composé} & \textbf{Groupes liés à N} & \textbf{Classe} \\
\hline
A : \ce{CH3-CH2-NH2} & 1 (éthyle) & primaire \\
\hline
B : \ce{CH3-NH-CH2-CH3} & 2 (méthyle et éthyle) & secondaire \\
\hline
C : \ce{(CH3)3N} & 3 (trois méthyle) & tertiaire \\
\hline
D : \ce{CH3-CH(NH2)-CH3} & 1 (isopropyle) & primaire \\
\hline
E : \ce{[(CH3)4N]+} & 4 & ion ammonium quaternaire \\
\hline
\end{tabularx}
\end{center}

\textbf{3. Isomères.} A a pour formule brute \ce{C2H7N} ; D a pour formule brute \ce{C3H9N} : ils ne sont pas isomères. En revanche B et C ont tous deux pour formule brute \ce{C3H9N} : ce sont des \cle{isomères de constitution} (B est secondaire, C est tertiaire).

\textbf{Conclusion :} A et D sont primaires, B est secondaire, C est tertiaire, E est un ion ammonium quaternaire ; B et C sont isomères.
\end{exoresolu}

\courssub{Nomenclature des amines}

\begin{methode}[Nomenclature des amines]
\begin{enumerate}
\item \textbf{Amines primaires} \ce{R-NH2} : deux méthodes possibles.
\begin{itemize}
\item Nom du groupe alkyle suivi du suffixe \emph{-amine} : \ce{CH3-NH2} = méthanamine (ou méthylamine).
\item Nom de l'alcane correspondant + suffixe \emph{-amine}, avec un indice de position si nécessaire : \ce{CH3-CH(NH2)-CH3} = propan-2-amine.
\end{itemize}
\item \textbf{Amines secondaires et tertiaires} : citer les groupes alkyles par ordre alphabétique devant le mot \emph{amine} ; les groupes identiques sont signalés par \emph{di-} ou \emph{tri-} : \ce{CH3-NH-C2H5} = éthylméthylamine ; \ce{(CH3)3N} = triméthylamine. En nomenclature systématique (IUPAC), le groupe principal est l'alcane le plus long, les autres sont des substituants \emph{N-} : \ce{CH3-NH-C2H5} = N-éthylméthanamine.
\item \textbf{Cas particulier à connaître par c\oe ur} : \ce{C6H5-NH2} = \cle{aniline} (nom systématique : phénylamine). C'est une amine aromatique primaire.
\item Vérifier que le nom rend compte de \emph{tous} les groupes portés par l'azote.
\end{enumerate}
\end{methode}

\begin{savaistu}[L'aniline et les alcaloïdes]
L'aniline (\ce{C6H5-NH2}), découverte dans l'indigo, est la base de l'industrie des colorants azoïques et des médicaments (paracétamol). Son doublet non liant est délocalisé dans le cycle benzénique : elle est \textbf{beaucoup moins basique} que les amines aliphatiques (pKa de l'anilinium \approx 4,6 contre \approx 10,6 pour le méthylammonium).
De nombreux alcaloïdes naturels (caféine, nicotine, quinine, morphine) contiennent des fonctions amines (souvent tertiaires ou hétérocycliques). Au Cameroun, le cacao et le café contiennent de la caféine (amine tertiaire purique) ; le tabac contient de la nicotine (amine secondaire et tertiaire).
\end{savaistu}

\begin{exoresolu}[Nommer et écrire des amines (type Bac)]
\begin{enumerate}
\item Nommer les amines suivantes : \ce{CH3-CH2-CH2-NH2} ; \ce{(C2H5)2NH} ; \ce{CH3-N(CH3)-C2H5}.
\item Écrire la formule semi-développée de la butan-2-amine et de l'aniline.
\end{enumerate}
\tcblower
\textbf{1. Noms.}
\begin{itemize}
\item \ce{CH3-CH2-CH2-NH2} : amine primaire à 3 carbones, groupe propyle $\Rightarrow$ \textbf{propan-1-amine} (ou propylamine).
\item \ce{(C2H5)2NH} : deux groupes éthyle sur l'azote $\Rightarrow$ \textbf{diéthylamine} (amine secondaire).
\item \ce{CH3-N(CH3)-C2H5} : trois groupes (deux méthyle, un éthyle), par ordre alphabétique $\Rightarrow$ \textbf{éthyldiméthylamine} (amine tertiaire).
\end{itemize}

\textbf{2. Formules.}

Butan-2-amine : chaîne de 4 carbones, \ce{-NH2} sur le carbone 2 :
\[ \chemfig{CH_3-CH(-[2]NH_2)-CH_2-CH_3} \]
Aniline (phénylamine) : groupe \ce{-NH2} fixé sur un noyau benzénique :
\[ \chemfig{*6(=-=-(-NH_2)=-)} \]

\textbf{Conclusion :} propan-1-amine, diéthylamine, éthyldiméthylamine ; la butan-2-amine est \ce{CH3-CH(NH2)-CH2-CH3} et l'aniline est \ce{C6H5-NH2}.
\end{exoresolu}

\courssub{Propriétés physiques des amines}

\begin{propriete}[État, solubilité, odeurs et liaisons hydrogène]
\begin{itemize}
\item \textbf{État} : Les amines de faible masse molaire (jusqu'à 3-4 carbones) sont des \textbf{gaz} à l'odeur forte d'ammoniac/poisson. Les suivantes sont des \textbf{liquides} (odeur persistante), puis des solides.
\item \textbf{Liaisons hydrogène} : L'atome d'azote (électronégativité 3,0) polarise les liaisons \ce{N-H}. Les amines primaires (2 \ce{N-H}) et secondaires (1 \ce{N-H}) forment des \cle{liaisons hydrogène intermoléculaires} \ce{N-H...N}. Les amines tertiaires (0 \ce{N-H}) \textbf{ne peuvent pas} en former \emph{entre elles}.
\item \textbf{Températures d'ébullition} : Pour une même masse molaire : Primaire $>$ Secondaire $>$ Tertiaire $>$ Alcanes. Les liaisons hydrogène augmentent la cohésion.
\item \textbf{Solubilité dans l'eau} : Les amines de faible masse molaire sont \textbf{miscibles} à l'eau. Les primaires et secondaires donnent \emph{et} acceptent des liaisons H avec l'eau. Les tertiaires n'\emph{acceptent} que les liaisons H de l'eau (via leur doublet) : elles restent solubles si la chaîne carbonée est courte. La solubilité diminue quand la chaîne carbonée s'allonge (partie hydrophobe croissante).
\item \textbf{Odeur} : Toutes les amines volatiles ont une odeur caractéristique, désagréable, rappelant l'ammoniac ou le poisson pourri (putrescine, cadavérine).
\end{itemize}
\end{propriete}

\begin{popfigure}
\begin{center}
\begin{tikzpicture}[scale=0.8, every node/.style={transform shape}]
% Primary amine H-bonds
\begin{scope}[xshift=-5cm]
\node[align=center] at (0,2.5) {\textbf{Amine primaire}\\\ce{R-NH2}};
\node at (0,1) {\chemfig{[:90]R-N(-[2]H)-[6]H}};
% H-bonds
\draw[popblue, thick, dashed] (0.5,0.5) -- (1.5,0);
\draw[popblue, thick, dashed] (-0.5,0.5) -- (-1.5,0);
\node[popblue, right] at (1.5,0) {$\cdots$H--NR$_2$};
\node[popblue, left] at (-1.5,0) {$\cdots$H--NR$_2$};
\end{scope}
% Secondary amine H-bonds
\begin{scope}[xshift=-1cm]
\node[align=center] at (0,2.5) {\textbf{Amine secondaire}\\\ce{R2NH}};
\node at (0,1) {\chemfig{[:90]R-N(-[2]R)-[6]H}};
\draw[popblue, thick, dashed] (0,0.3) -- (0,-0.8);
\node[popblue, below] at (0,-0.8) {$\cdots$H--NR$_2$};
\end{scope}
% Tertiary amine: NO H-bonds between them
\begin{scope}[xshift=3cm]
\node[align=center] at (0,2.5) {\textbf{Amine tertiaire}\\\ce{R3N}};
\node at (0,1) {\chemfig{[:90]R-N(-[2]R)-[6]R}};
\node[poporange, align=center] at (0,-0.5) {\textbf{Pas de liaison H}\\\textbf{intermoléculaire}};
\end{scope}
% Tertiary + Water
\begin{scope}[xshift=7cm]
\node[align=center] at (0,2.5) {\textbf{Tertiaire + Eau}};
\node at (0,1.2) {\chemfig{[:90]R-N(-[2]R)-[6]R}};
\node at (0,-0.2) {\chemfig{H-O-H}};
\draw[popgreen, thick, dashed] (0,0.5) -- (0,0);
\node[popgreen, right] at (0.5,0.2) {$N\cdots H-O$};
\node[popgreen, align=center] at (0,-1.5) {L'azote \textbf{accepte}\\\textbf{la liaison H de l'eau}};
\end{scope}
\end{tikzpicture}
\end{center}
\legende{Rôle des liaisons hydrogène : donneurs (prim./sec.) vs accepteurs seulement (tert.).}
\end{popfigure}

\begin{exoresolu}[Ébullition et solubilité (type Bac)]
On donne les températures d'ébullition : éthylamine \ce{C2H5-NH2} : $\theta_{eb} = \SI{17}{\celsius}$ ; diméthylamine \ce{(CH3)2NH} : $\theta_{eb} = \SI{7}{\celsius}$ ; triméthylamine \ce{(CH3)3N} : $\theta_{eb} = \SI{3}{\celsius}$ ; propane \ce{C3H8} : $\theta_{eb} = \SI{-42}{\celsius}$.
\begin{enumerate}
\item Pourquoi les trois amines bout-elles à des températures bien plus élevées que le propane ?
\item Expliquer l'ordre d'ébullition des trois amines.
\item La triméthylamine est pourtant soluble dans l'eau. Justifier.
\end{enumerate}
\tcblower
\textbf{1. Comparaison avec le propane.} Le propane est apolaire : seules de faibles interactions de Van der Waals existent entre ses molécules, d'où une ébullition très basse. Les amines possèdent un atome d'azote électronégatif porteur d'un doublet non liant ; celles qui ont des liaisons \ce{N-H} établissent des \cle{liaisons hydrogène} intermoléculaires, plus fortes, qui demandent davantage d'énergie pour être rompues.

\textbf{2. Ordre des amines.} L'éthylamine (primaire) possède \textbf{deux} liaisons \ce{N-H} : elle forme le plus de liaisons hydrogène $\Rightarrow$ ébullition la plus élevée (\SI{17}{\celsius}). La diméthylamine (secondaire) n'en possède qu'\textbf{une} $\Rightarrow$ \SI{7}{\celsius}. La triméthylamine (tertiaire) n'en possède \textbf{aucune} : pas de liaison hydrogène entre ses molécules, seulement des interactions dipôle-dipôle $\Rightarrow$ ébullition la plus basse (\SI{3}{\celsius}).

\textbf{3. Solubilité de la triméthylamine.} Même sans liaison \ce{N-H}, l'azote garde son \cle{doublet non liant} : il \emph{accepte} des liaisons hydrogène avec les molécules d'eau (\ce{N...H-O}). La triméthylamine, de petite taille, est donc soluble dans l'eau.

\textbf{Conclusion :} les liaisons hydrogène (données ou acceptées) expliquent à la fois les températures d'ébullition et la solubilité des amines.
\end{exoresolu}

\courssub{Caractère basique des amines}

\begin{propriete}[Réaction avec l'eau : basicité]
L'amine est une \cle{base de Brønsted} : le doublet non liant de l'azote capture un proton \ce{H+} fourni par l'eau (acide). C'est un \textbf{équilibre limité} (base faible).
\[ \ce{R-NH2 + H2O <=> R-NH3+ + HO-} \]
Le couple acide/base est \ce{R-NH3+}/\ce{R-NH2} (acide conjugué / base). La solution contient des ions \ce{HO-} : \textbf{pH > 7}.
\begin{itemize}
\item \textbf{Indicateurs} : Bleu de bromothymol (BBT) $\to$ \textbf{bleu} ; Phénolphtaléine $\to$ \textbf{rose}.
\item \textbf{Comparaison avec l'ammoniac} : Les groupes alkyles sont \emph{donneurs d'électrons} (+I), ils augmentent la densité électronique sur l'azote $\Rightarrow$ les amines aliphatiques sont \textbf{plus basiques} que l'ammoniac (pKa \ce{CH3NH3+} \approx 10,6 ; pKa \ce{NH4+} = 9,25). L'aniline est moins basique (pKa \approx 4,6) car le doublet est délocalisé dans le cycle.
\end{itemize}
\end{propriete}

\begin{methode}[Écrire l'équation-bilan amine + eau]
\begin{enumerate}
\item Identifier le couple : \ce{R-NH3+}/\ce{R-NH2} pour l'amine ; \ce{H2O}/\ce{HO-} pour l'eau.
\item Le doublet de l'azote capture \ce{H+} de l'eau.
\item Écrire l'équilibre (double flèche \ce{<=>}) :
\[ \ce{R-NH2 + H2O <=> R-NH3+ + HO-} \]
\item Conclure : solution basique (pH > 7), virage BBT (bleu), phénolphtaléine (rose).
\item \textbf{Piège} : ne jamais écrire de flèche simple \ce{->} ni oublier la charge $+$ de l'ion alkylammonium \ce{R-NH3+}.
\end{enumerate}
\end{methode}

\begin{popfigure}
\begin{center}
\begin{tikzpicture}
\begin{axis}[
    width=\linewidth, height=6cm,
    axis lines=middle,
    axis line style={popdark},
    tick style={popdark},
    tick label style={popdark},
    grid=major,
    grid style={popline},
    minor grid style={popline!50},
    xlabel={Volume d'acide fort ajout\'e (mL)},
    ylabel={pH},
    xlabel style={anchor=north, popdark},
    ylabel style={anchor=east, popdark},
    xmin=0, xmax=25, ymin=0, ymax=14,
    xtick={0,10,20},
    ytick={0,7,10.6,14},
    yticklabels={0,7,pK\_a,14},
    domain=0:25, samples=200,
]
% Courbe de titration schématique : base faible + acide fort
% Équivalence à 10 mL, pH ~ 5,5 à l'éq. pH initial ~ 11,5. Demi-éq à 5 mL, pH = pKa ~ 10,6
\addplot[popblue, thick, mark=none, smooth] coordinates {
    (0,11.5) (2,11.2) (5,10.6) (8,9.5) (9.5,7) (10,5.5) (10.5,4) (12,2.5) (15,1.5) (20,1)
};
\draw[dashed, popmuted] (0,10.6) -- (5,10.6) node[left, popmuted] {pK\_a = 10,6};
\draw[dashed, popmuted] (5,0) -- (5,10.6);
\draw[dashed, popmuted] (10,0) -- (10,5.5);
\node[popdark, anchor=north] at (axis cs:5,11) {Demi-\'equivalence};
\node[popdark, anchor=north] at (axis cs:10,6) {\'Equivalence};
\node[popdark, anchor=south west] at (axis cs:1,12) {Zone tampon};
\end{axis}
\end{tikzpicture}
\end{center}
\legende{Courbe de dosage d'une amine (base faible) par un acide fort. Au point demi-équivalence, pH = pKa du couple \ce{R-NH3+}/\ce{R-NH2}.}
\end{popfigure}

\begin{exoresolu}[Solution aqueuse de méthylamine (type Bac)]
On dissout $n = \SI{0,020}{mol}$ de méthylamine \ce{CH3-NH2} dans de l'eau distillée pour obtenir $V = \SI{200}{mL}$ de solution. Le pH mesuré vaut $11{,}8$.
\begin{enumerate}
\item Écrire l'équation-bilan de la réaction de la méthylamine avec l'eau.
\item Calculer la concentration apportée $C$ en méthylamine, puis la concentration en ions hydroxyde.
\item Montrer que la réaction est limitée et calculer son taux d'avancement final $\tau$.
\end{enumerate}
\tcblower
\textbf{1. Équation-bilan.} La méthylamine capte un proton de l'eau (couple \ce{CH3NH3+}/\ce{CH3NH2}) :
\[ \ce{CH3-NH2 + H2O <=> CH3-NH3+ + HO-} \]

\textbf{2. Concentrations.}
\[ C = \frac{n}{V} = \frac{\SI{0,020}{mol}}{\SI{0,200}{L}} = \SI{0,10}{mol.L^{-1}} \]
D'après le pH : $[\ce{H3O+}] = 10^{-11{,}8} = \SI{1,58e-12}{mol.L^{-1}}$, d'où, avec $K_e = \SI{1,0e-14}{}$ à \SI{25}{\celsius} :
\[ [\ce{HO-}] = \frac{K_e}{[\ce{H3O+}]} = \frac{10^{-14}}{10^{-11{,}8}} = 10^{-2{,}2} = \SI{6,31e-3}{mol.L^{-1}} \]

\textbf{3. Réaction limitée.} Si la réaction était totale, on aurait $[\ce{HO-}] = C = \SI{0,10}{mol.L^{-1}}$, soit pH $= 13$. Or $[\ce{HO-}] = \SI{6,31e-3}{mol.L^{-1}} \ll C$ : la réaction est \cle{limitée} (équilibre). Le taux d'avancement final est :
\[ \tau = \frac{[\ce{HO-}]}{C} = \frac{\num{6,31e-3}}{\num{0,10}} = \num{6,31e-2} \approx 6,3\,\% \]

\textbf{Conclusion :} la méthylamine est une base faible : seulement $6,3\,\%$ des molécules sont protonées ; la solution, de pH $11{,}8$, fait virer le BBT au bleu.
\end{exoresolu}

\courssub{Action d'une amine sur les ions métalliques}

\begin{methode}[Précipitation des hydroxydes métalliques]
\begin{enumerate}
\item Une solution d'amine contient des ions \ce{HO-} (réaction limitée avec l'eau).
\item Ajoutée à une solution d'ions métalliques (\ce{Cu^2+}, \ce{Fe^3+}, \ce{Al^3+}...), elle provoque la \cle{précipitation de l'hydroxyde métallique}, comme le ferait la soude ou l'ammoniac :
\[ \ce{Cu^2+ + 2 HO- -> Cu(OH)2 (s)} \]
\item Retenir les couleurs caractéristiques : \ce{Cu(OH)2} \textbf{bleu}, \ce{Fe(OH)3} \textbf{rouille}, \ce{Al(OH)3} \textbf{blanc gélatineux}.
\item Écrire l'équation globale avec l'amine si demandé :
\[ \ce{Cu^2+ + 2 R-NH2 + 2 H2O -> Cu(OH)2 (s) + 2 R-NH3+} \]
\end{enumerate}
\end{methode}

\begin{exoresolu}[Dépistage des ions cuivre (type Bac)]
Dans un laboratoire de lycée à Douala, on verse quelques gouttes d'une solution d'éthylamine dans un tube contenant une solution bleue de sulfate de cuivre(II). On observe l'apparition d'un précipité bleu.
\begin{enumerate}
\item Interpréter l'observation et écrire l'équation-bilan de la réaction.
\item On a utilisé $V = \SI{10,0}{mL}$ de solution de \ce{Cu^2+} à $C = \SI{0,050}{mol.L^{-1}}$. Calculer la masse maximale de précipité obtenue. On donne $M(\ce{Cu(OH)2}) = \SI{97,5}{g.mol^{-1}}$.
\end{enumerate}
\tcblower
\textbf{1. Interprétation.} L'éthylamine réagit avec l'eau : \ce{C2H5-NH2 + H2O <=> C2H5-NH3+ + HO-}. Les ions \ce{HO-} formés précipitent les ions cuivre(II) en hydroxyde de cuivre(II) bleu :
\[ \ce{Cu^2+ + 2 HO- -> Cu(OH)2 (s)} \]
soit globalement : \ce{Cu^2+ + 2 C2H5-NH2 + 2 H2O -> Cu(OH)2 (s) + 2 C2H5-NH3+}.

\textbf{2. Masse de précipité.} Quantité d'ions cuivre :
\[ n(\ce{Cu^2+}) = C \times V = \num{0,050} \times \num{10,0e-3} = \SI{5,0e-4}{mol} \]
D'après la st\oe chiométrie, $n(\ce{Cu(OH)2}) = n(\ce{Cu^2+}) = \SI{5,0e-4}{mol}$, d'où :
\[ m = n \times M = \num{5,0e-4} \times \num{97,5} = \SI{4,875e-2}{g} \approx \SI{49}{mg} \]

\textbf{Conclusion :} le précipité bleu de \ce{Cu(OH)2} (masse maximale $\approx \SI{49}{mg}$) met en évidence le caractère basique de l'éthylamine.
\end{exoresolu}

\courssub{Caractère nucléophile : alkylation et réactions d'Hofmann}

\begin{propriete}[Nucléophilicité de l'amine]
L'atome d'azote porte un doublet non liant disponible. L'amine est un \cle{nucléophile} (donneur de doublet). Elle attaque les centres \cle{électrophiles} (déficitaires en électrons) : carbone porteur d'un halogène (dérivé halogéné) ou carbone de carbonyle (chlorure d'acyle, anhydride). La réaction est une \textbf{substitution nucléophile}.
\end{propriete}

\begin{methode}[Amine + dérivé halogéné : réactions d'Hofmann successives]
\begin{enumerate}
\item \textbf{Rôles} : Amine = nucléophile (doublet sur N) ; Dérivé halogéné \ce{R$'$-X} = électrophile (C-\ce{X} polarisé, \ce{X} électronégatif).
\item \textbf{Mécanisme} : Attaque du doublet sur le carbone, départ de \ce{X-} (SN2), puis déprotonation par une seconde molécule d'amine.
\item \textbf{Équation-bilan globale} (2 mol d'amine pour 1 mol de \ce{R$'$-X}) :
\[ \ce{2 R-NH2 + R$'$-X -> R-NH-R$'$ + R-NH3+ + X-} \]
\item \textbf{Point clé (Hofmann)} : L'amine formée (\ce{R-NH-R$'$}) est \emph{encore nucléophile} (doublet libre). La réaction se poursuit successivement :
\[ \text{Amine 1\degre} \xrightarrow{\ce{R$'$-X}} \text{Amine 2\degre} \xrightarrow{\ce{R$'$-X}} \text{Amine 3\degre} \xrightarrow{\ce{R$'$-X}} \text{Ion ammonium quaternaire } \ce{[R4N]+} \]
La dernière étape ne consomme qu'\textbf{1 mole} d'amine tertiaire (pas de proton à capter) :
\[ \ce{R3N + R$'$-X -> [R4N]+ + X-} \]
\item \textbf{Conséquence} : L'alkylation n'est \textbf{jamais sélective} ; on obtient un mélange.
\end{enumerate}
\end{methode}

\begin{popfigure}
\begin{center}
\begin{tikzpicture}[node distance=1.5cm, >=Stealth]
\node (start) {\chemfig{CH_3-NH_2}};
\node[right=of start] (plus1) {+\ce{CH3I}};
\node[right=of plus1] (sec) {\chemfig{CH_3-NH-CH_3}};
\node[below=of sec] (label1) {\small Amine secondaire};
\node[right=of sec] (plus2) {+\ce{CH3I}};
\node[right=of plus2] (tert) {\chemfig{CH_3-N(-[2]CH_3)-CH_3}};
\node[below=of tert] (label2) {\small Amine tertiaire};
\node[right=of tert] (plus3) {+\ce{CH3I}};
\node[right=of plus3] (quat) {\chemfig{CH_3-N^{+}(-[2]CH_3)(-[6]CH_3)-CH_3} \ce{I-}};
\node[below=of quat] (label3) {\small Ion ammonium quaternaire};
\draw[->, popblue, thick] (start) -- (plus1);
\draw[->, popblue, thick] (plus1) -- (sec);
\draw[->, popblue, thick] (sec) -- (plus2);
\draw[->, popblue, thick] (plus2) -- (tert);
\draw[->, popblue, thick] (tert) -- (plus3);
\draw[->, popblue, thick] (plus3) -- (quat);
\node[poporange, align=center, above=0.5cm of sec] {2 \ce{CH3NH2}\\\ce{->} \ce{CH3NHCH3} + \ce{CH3NH3+I-}};
\node[poporange, align=center, above=0.5cm of tert] {2 \ce{CH3NHCH3}\\\ce{->} \ce{(CH3)3N} + \ce{CH3NHCH3+I-}};
\node[poporange, align=center, above=0.5cm of quat] {1 \ce{(CH3)3N}\\\ce{->} \ce{[(CH3)4N]+I-}};
\end{tikzpicture}
\end{center}
\legende{Suite des réactions d'Hofmann : alkylation successive de la méthylamine par l'iodométhane.}
\end{popfigure}

\begin{exoresolu}[Alkylation de l'ammoniac et de la méthylamine (type Bac)]
Pour synthétiser un intermédiaire de médicament, on fait réagir la méthylamine \ce{CH3-NH2} avec l'iodoéthane \ce{C2H5-I}.
\begin{enumerate}
\item Justifier le caractère nucléophile de la méthylamine et électrophile de l'iodoéthane.
\item Écrire l'équation-bilan de la réaction conduisant à l'amine secondaire et nommer le produit.
\item Quels produits obtient-on si l'on poursuit l'alkylation ? Écrire l'équation de la dernière étape.
\end{enumerate}
\tcblower
\textbf{1. Rôles des réactifs.} L'azote de \ce{CH3-NH2} porte un \cle{doublet non liant} : site riche en électrons, donc nucléophile. Dans \ce{C2H5-I}, l'iode, très électronégatif, attire les électrons de la liaison \ce{C-I} : le carbone lié à l'iode est déficitaire en électrons, donc électrophile.

\textbf{2. Formation de l'amine secondaire.} Substitution nucléophile :
\[ \ce{2 CH3-NH2 + C2H5-I -> CH3-NH-C2H5 + CH3-NH3+ + I-} \]
Le produit organique est l'\textbf{éthylméthylamine} (amine secondaire). La seconde molécule de méthylamine neutralise \ce{HI} formé.

\textbf{3. Poursuite de l'alkylation (réactions d'Hofmann).} L'amine secondaire formée possède encore un doublet : elle réagit à son tour pour donner l'amine tertiaire \textbf{éthyldiméthylamine} \ce{CH3-N(CH3)-C2H5}, puis celle-ci donne l'\cle{ion ammonium quaternaire} :
\[ \ce{CH3-N(CH3)-C2H5 + C2H5-I -> [CH3-N(CH3)(C2H5)-C2H5]+ + I-} \]
soit l'ion \textbf{diéthyldiméthylammonium} associé à l'ion iodure.

\textbf{Conclusion :} l'alkylation d'une amine n'est jamais sélective : elle conduit à un mélange d'amine secondaire, tertiaire et de sel d'ammonium quaternaire.
\end{exoresolu}

\courssub{Action sur les chlorures d'acyle : préparation des amides}

\begin{propriete}[Électrophilicité du chlorure d'acyle]
Le chlorure d'acyle \ce{R-COCl} possède un carbone de carbonyle très électrophile : il est lié à deux atomes électronégatifs (O et Cl) qui attirent la densité électronique. C'est un réactif très réactif en substitution nucléophile acylque.
\end{propriete}

\begin{methode}[Amine + chlorure d'acyle : synthèse d'un amide]
\begin{enumerate}
\item \textbf{Rôles} : Amine = nucléophile ; Chlorure d'acyle \ce{R-COCl} = électrophile (C=O).
\item \textbf{Mécanisme} : Attaque du doublet de N sur le C=O, formation d'un tétraédrique, élimination de \ce{Cl-}, puis déprotonation par une seconde amine.
\item \textbf{Équation-bilan} (réaction \textbf{rapide et totale}, flèche simple) :
\[ \ce{R-COCl + 2 R$'$-NH2 -> R-CO-NH-R$'$ + R$'$-NH3+ + Cl-} \]
La deuxième mole d'amine neutralise le \ce{HCl} formé (acide fort).
\item \textbf{Nomenclature de l'amide} : Nom de l'acide (sans « acide ...-oïque ») + « -amide », précédé de \emph{N-} + nom du groupe sur l'azote. Ex. : \ce{CH3-CO-NH-C2H5} = \textbf{N-éthyléthanamide}.
\item \textbf{Variante} : L'anhydride d'acide \ce{(R-CO)2O} peut remplacer le chlorure d'acyle (réaction plus lente, souvent chauffage).
\end{enumerate}
\end{methode}

\begin{exoresolu}[Synthèse d'un amide (type Bac)]
On fait réagir le chlorure d'éthanoyle \ce{CH3-COCl} avec la méthylamine \ce{CH3-NH2} en solution dans l'éther.
\begin{enumerate}
\item Écrire l'équation-bilan de la réaction et nommer le produit organique obtenu.
\item Pourquoi utilise-t-on au moins deux moles de méthylamine par mole de chlorure d'acyle ?
\item On fait réagir $n_1 = \SI{0,050}{mol}$ de chlorure d'éthanoyle avec un excès de méthylamine. La réaction étant totale, calculer la masse d'amide obtenue. On donne $M(\ce{C3H7NO}) = \SI{73}{g.mol^{-1}}$.
\end{enumerate}
\tcblower
\textbf{1. Équation-bilan.} La méthylamine, nucléophile, attaque le carbone électrophile du chlorure d'acyle :
\[ \ce{CH3-COCl + 2 CH3-NH2 -> CH3-CO-NH-CH3 + CH3-NH3+ + Cl-} \]
Le produit organique est le \textbf{N-méthyléthanamide}, de formule semi-développée :
\[ \chemfig{CH_3-C(=[1]O)-[7]NH-CH_3} \]

\textbf{2. Rôle de la deuxième mole d'amine.} La réaction libère du chlorure d'hydrogène \ce{HCl}, acide fort, qui protonerait l'amine nucléophile et la rendrait inerte. La deuxième mole de méthylamine sert à \cle{neutraliser} \ce{HCl} sous forme de chlorure de méthylammonium \ce{[CH3-NH3]+ Cl-}.

\textbf{3. Masse d'amide.} La réaction est totale et la st\oe chiométrie donne $n(\text{amide}) = n_1 = \SI{0,050}{mol}$ (la méthylamine est en excès, le chlorure d'acyle est limitant) :
\[ m = n \times M = \num{0,050} \times 73 = \SI{3,65}{g} \approx \SI{3,7}{g} \]

\textbf{Conclusion :} on obtient environ \SI{3,7}{g} de N-méthyléthanamide ; c'est ce type de réaction (acylation d'amine) qui permet de fabriquer de nombreux médicaments (paracétamol, pénicillines, anesthésiques locaux) contenant la fonction amide.
\end{exoresolu}

\begin{savaistu}[Applications industrielles et locales]
\begin{itemize}
\item \textbf{Paracétamol} : synthétisé par acylation de la p-aminophénol (amine aromatique) par l'anhydride éthanoïque.
\item \textbf{Pénicillines} : contiennent un cycle $\beta$-lactame (amide cyclique) issu d'une condensation amine/acide.
\item \textbf{Industrie camerounaise} : Les savonneries artisanales (huile de palme + soude) produisent de la glycérine ; les amines grasses (dérivées des acides gras du palmier) servent à fabriquer des tensioactifs cationiques (assouplissants, désinfectants) via alkylation (Hofmann) jusqu'au quaternaire.
\item \textbf{SONARA / CIMENCAM} : Utilisation d'amines comme inhibiteurs de corrosion (amines alicycliques) ou additifs pour ciments.
\end{itemize}
\end{savaistu}

\begin{attention}[Les 5 erreurs qui coûtent des points au Bac]
\begin{enumerate}
\item \textbf{Confondre la classe de l'amine avec le nombre d'hydrogènes.} On compte les \emph{groupes alkyles liés à l'azote}, pas les atomes H : \ce{(CH3)2NH} est secondaire, \ce{(CH3)3N} est tertiaire.
\item \textbf{Écrire la réaction amine + eau avec une flèche simple.} C'est un équilibre (base faible) : \ce{R-NH2 + H2O <=> R-NH3+ + HO-}. Oublier la charge $+$ de \ce{R-NH3+} est également sanctionné.
\item \textbf{Oublier la deuxième mole d'amine} dans les équations d'alkylation ou d'acylation : il faut \textbf{2 mol} d'amine pour 1 mol de \ce{R-X} ou de \ce{R-COCl}, car une mole neutralise \ce{HX} ou \ce{HCl}.
\item \textbf{Croire que l'alkylation s'arrête à l'amine voulue.} L'amine produite reste nucléophile : les réactions d'Hofmann se poursuivent jusqu'à l'ion ammonium quaternaire \ce{[R4N]+}. Le jury attend cette remarque.
\item \textbf{Attribuer des liaisons hydrogène intermoléculaires aux amines tertiaires.} Sans liaison \ce{N-H}, elles ne peuvent en former \emph{entre elles} (ébullition basse) ; elles ne font qu'\emph{accepter} celles de l'eau (d'où leur solubilité).
\end{enumerate}
\end{attention}

\newpage

\coursec{Exercices}

\courssub{Je vérifie mes connaissances}

\begin{exercice}[QCM : classes et propriétés]
Pour chaque question, choisir la (ou les) bonne(s) réponse(s).
\begin{enumerate}
\item La molécule d'ammoniac \ce{NH3} est :
\begin{enumerate}
\item plane ; \item pyramidale ; \item tétraédrique régulière ; \item linéaire.
\end{enumerate}
\item L'amine de formule \ce{CH3-NH-CH2-CH3} est :
\begin{enumerate}
\item primaire ; \item secondaire ; \item tertiaire ; \item un ion ammonium quaternaire.
\end{enumerate}
\item L'ion \ce{(CH3)4N+} est :
\begin{enumerate}
\item une amine tertiaire ; \item un ion ammonium quaternaire ; \item une base ; \item obtenu par alkylation totale de l'ammoniac.
\end{enumerate}
\item À masse molaire comparable, une amine primaire a une température d'ébullition :
\begin{enumerate}
\item inférieure à celle de l'alcane correspondant ; \item supérieure à celle de l'alcane correspondant ; \item inférieure à celle de l'alcool correspondant ; \item supérieure à celle de l'alcool correspondant.
\end{enumerate}
\item Le caractère nucléophile des amines est dû :
\begin{enumerate}
\item à la présence du doublet non liant de l'azote ; \item à la polarité de la liaison \ce{N-H} ; \item à la charge positive de l'azote ; \item à l'effet donneur des groupes alkyles.
\end{enumerate}
\end{enumerate}
\end{exercice}

\begin{exercice}[Vrai ou faux justifié]
Répondre par \emph{vrai} ou \emph{faux}, puis justifier brièvement.
\begin{enumerate}
\item Une amine tertiaire ne peut pas établir de liaisons hydrogène entre ses propres molécules.
\item L'aniline \ce{C6H5-NH2} est une amine secondaire.
\item Une solution aqueuse d'éthylamine fait virer le bleu de bromothymol (BBT) au jaune.
\item Dans la réaction d'une amine avec un chlorure d'acyle, l'atome de carbone du groupe carbonyle est le site électrophile.
\item Une amine tertiaire réagit avec un dérivé halogéné pour donner un ion ammonium quaternaire.
\end{enumerate}
\end{exercice}

\begin{exercice}[Texte à trous]
Compléter le texte suivant avec les mots ou expressions : \emph{doublet non liant ; basique ; nucléophile ; hydrogène ; pyramidale ; ammonium quaternaire ; alkylation ; électrophile}.

La molécule d'ammoniac a une géométrie \ldots\ldots\ldots\ldots et possède sur l'atome d'azote un \ldots\ldots\ldots\ldots\ldots\ldots\ldots, à l'origine de ses propriétés chimiques. Les amines sont des bases : elles captent un proton, ce qui traduit leur caractère \ldots\ldots\ldots\ldots. Elles sont aussi des réactifs \ldots\ldots\ldots\ldots : elles attaquent les sites \ldots\ldots\ldots\ldots des molécules, comme l'atome de carbone porteur d'un halogène. La réaction d'une amine avec un dérivé halogéné est une réaction d'\ldots\ldots\ldots\ldots ; répétée, elle conduit finalement à un ion \ldots\ldots\ldots\ldots\ldots\ldots\ldots. Les amines primaires et secondaires établissent entre elles des liaisons \ldots\ldots\ldots\ldots, ce qui explique leurs températures d'ébullition relativement élevées.
\end{exercice}

\begin{exercice}[Définitions]
Définir en une ou deux phrases chacun des termes suivants et donner un exemple de formule semi-développée :
\begin{enumerate}
\item amine primaire ; amine secondaire ; amine tertiaire ;
\item ion ammonium quaternaire ;
\item réaction d'Hofmann (alkylation d'une amine) ;
\item caractère nucléophile.
\end{enumerate}
\end{exercice}

\courssub{Je m'entraîne}

\begin{exercice}[Classes et formules]
On considère les amines suivantes :
\ce{CH3-CH2-NH2} ; \quad \ce{(CH3)2NH} ; \quad \ce{CH3-N(CH3)-CH2-CH3} ; \quad \ce{CH3-CH(NH2)-CH3} ; \quad \ce{C6H5-NH2}.
\begin{enumerate}
\item Écrire la formule semi-développée de chaque amine.
\item Préciser la classe de chacune d'elles en justifiant.
\item Identifier deux amines isomères de formule brute \ce{C3H9N} et écrire tous les autres isomères possibles de cette formule brute, en précisant leur classe.
\end{enumerate}
\end{exercice}

\begin{exercice}[Nomenclature]
\begin{enumerate}
\item Nommer, en nomenclature systématique, les amines suivantes : \ce{CH3-CH2-CH2-NH2} ; \ce{CH3-NH-CH2-CH3} ; \ce{(CH3)3N} ; \ce{CH3-CH2-CH(NH2)-CH3} ; \ce{C6H5-NH2}.
\item Écrire la formule semi-développée des amines dont les noms suivent : éthanamine ; N-méthylpropan-1-amine ; N,N-diméthylméthanamine ; butan-2-amine ; phénylamine (aniline).
\end{enumerate}
\end{exercice}

\begin{exercice}[Basicité et équations-bilans]
On dissout de l'éthylamine \ce{C2H5-NH2} dans l'eau distillée.
\begin{enumerate}
\item Écrire l'équation-bilan de la réaction de l'éthylamine avec l'eau.
\item Identifier les deux couples acide/base mis en jeu.
\item Le pH de la solution obtenue est supérieur à 7. Justifier. Quelle est la couleur prise par l'hélianthine dans cette solution ?
\item On ajoute goutte à goutte une solution de chlorure d'hydrogène. Écrire l'équation-bilan de la réaction qui se produit et nommer le produit ionique formé.
\end{enumerate}
\end{exercice}

\begin{exercice}[Équations de réactions caractéristiques]
Écrire les équations-bilans des réactions suivantes, en précisant le nom et la classe (ou la nature) du produit organique obtenu :
\begin{enumerate}
\item action de la méthylamine \ce{CH3-NH2} sur l'iodoéthane \ce{C2H5-I} ;
\item action de la diéthylamine \ce{(C2H5)2NH} sur le bromométhane \ce{CH3-Br} ;
\item action de la triméthylamine \ce{(CH3)3N} sur l'iodométhane ;
\item action de l'éthylamine \ce{C2H5-NH2} sur le chlorure d'éthanoyle \ce{CH3-COCl} ;
\item action de l'ammoniac sur le chlorure d'éthanoyle.
\end{enumerate}
\end{exercice}

\courssub{Je me prépare au Bac}

\begin{exercice}[Type Bac 1 — Structure, classes et solubilité des amines]
\emph{(D'après Baccalauréat C, session normale)}

On donne les formules de cinq composés azotés notés A, B, C, D et E :
\begin{center}
A : \ce{(CH3)3C-NH2} \qquad B : \ce{CH3-NH-C2H5} \qquad C : \ce{C6H5-NH2} \qquad D : \ce{(C2H5)3N} \qquad E : \ce{CH3-CH2-CH2-CH2-NH2}
\end{center}
Données : masses molaires atomiques en \si{g.mol^{-1}} : $M(\ce{H})=1$ ; $M(\ce{C})=12$ ; $M(\ce{N})=14$.
\begin{enumerate}
\item Écrire la formule semi-développée de chacun des composés A, B, C, D et E, puis préciser la classe de chaque amine en justifiant.
\item Nommer chaque amine en nomenclature systématique ; donner également le nom courant du composé C.
\item Vérifier que les composés A et E sont isomères. Quel type d'isomérie les relie ?
\item On constate que l'amine E est nettement plus soluble dans l'eau que l'amine D, alors que leurs masses molaires sont voisines. Expliquer cette différence de solubilité en faisant intervenir les liaisons hydrogène.
\item Comparer, en les justifiant, les températures d'ébullition du butane ($M = \SI{58}{g.mol^{-1}}$), de l'amine E et du butan-1-ol ($M = \SI{74}{g.mol^{-1}}$).
\end{enumerate}
\end{exercice}

\begin{exercice}[Type Bac 2 — Caractère basique de la méthylamine]
\emph{(D'après Baccalauréat D, session normale)}

On prépare une solution aqueuse $(S)$ de méthylamine \ce{CH3-NH2} de concentration molaire $C = \SI{0,10}{mol.L^{-1}}$. La mesure du pH à \SI{25}{\celsius} donne $\text{pH} = 11{,}3$.

Données : zones de virage : hélianthine : 3,1 -- 4,4 (rouge/jaune-orange) ; BBT : 6,0 -- 7,6 (jaune/bleu) ; phénolphtaléine : 8,2 -- 10,0 (incolore/rose). $\text{p}K_a$ du couple \ce{CH3NH3+}/\ce{CH3NH2} : 10,7 ; $\text{p}K_a$ du couple \ce{NH4+}/\ce{NH3} : 9,2 ; $\text{p}K_e = 14$ à \SI{25}{\celsius}.
\begin{enumerate}
\item Écrire l'équation-bilan de la réaction de la méthylamine avec l'eau et identifier les deux couples acide/base mis en jeu.
\item Calculer la concentration en ions hydroxyde \ce{HO-} dans la solution $(S)$. La méthylamine est-elle une base forte ou une base faible ? Justifier à l'aide de la valeur du $\text{p}K_a$.
\item Prévoir la teinte prise par le BBT et par la phénolphtaléine dans la solution $(S)$. Justifier.
\item On ajoute à un échantillon de $(S)$ quelques gouttes d'une solution de sulfate de cuivre(II) \ce{CuSO4}. Décrire ce qu'on observe et écrire l'équation-bilan de la réaction responsable du précipité formé.
\item Comparer la basicité de la méthylamine à celle de l'ammoniac. Justifier cette différence par un effet électronique lié au groupe méthyle.
\item On désire préparer \SI{500}{mL} de la solution $(S)$ à partir de méthylamine commerciale de masse molaire $M = \SI{31}{g.mol^{-1}}$. Calculer la masse de méthylamine à dissoudre.
\end{enumerate}
\end{exercice}

\begin{exercice}[Type Bac 3 — Réactions d'Hofmann et synthèse d'un médicament]
\emph{(D'après Baccalauréat C, session de rattrapage)}

\textbf{Partie A : alkylation de l'ammoniac.}

On fait réagir de l'ammoniac \ce{NH3} avec un excès d'iodométhane \ce{CH3-I}, en chauffant dans un solvant approprié.
\begin{enumerate}
\item Écrire les équations-bilans des quatre réactions successives qui conduisent de l'ammoniac à l'ion tétraméthylammonium.
\item Nommer chaque produit organique obtenu et préciser sa classe (amine primaire, secondaire, tertiaire ou ion ammonium quaternaire).
\item Expliquer, en termes de nucléophilie et d'électrophilie, le mécanisme de la première étape : quel site de l'iodométhane est attaqué et pourquoi ?
\item Quel réactif faut-il mettre en grand excès pour obtenir majoritairement la méthylamine \ce{CH3-NH2} ? Justifier.
\end{enumerate}

\textbf{Partie B : synthèse du paracétamol.}

Le paracétamol, principe actif de nombreux médicaments vendus dans les pharmacies du Cameroun, peut être obtenu par action du chlorure d'éthanoyle \ce{CH3-COCl} sur le 4-aminophénol \ce{HO-C6H4-NH2}.
\begin{enumerate}
\setcounter{enumi}{4}
\item Écrire l'équation-bilan de la réaction de formation du paracétamol (on notera le groupe phényle \ce{C6H4}).
\item Justifier le caractère électrophile du chlorure d'éthanoyle et le caractère nucléophile de l'amine.
\item Expliquer pourquoi on utilise en pratique un excès d'amine.
\item Une amine tertiaire, par exemple \ce{(CH3)3N}, pourrait-elle conduire à un amide stable par action sur un chlorure d'acyle ? Justifier la réponse.
\end{enumerate}
\end{exercice}

\newpage

\coursec{Corrigés des exercices}

\courssub{Je vérifie mes connaissances}

\begin{corrige}[Exercice 1 — QCM : classes et propriétés]
\begin{enumerate}
\item \textbf{Réponse : b (pyramidale).} L'atome d'azote est entouré de trois atomes d'hydrogène et d'un doublet non liant. Selon la méthode VSEPR (type \ce{AX3E}), ce doublet occupe le sommet d'un tétraèdre et repousse les trois liaisons \ce{N-H} : la géométrie observée des atomes est donc \cle{pyramidale} (angle $\widehat{H-N-H} \approx 107^{\circ}$). Elle n'est ni plane, ni tétraédrique régulière (ce serait le cas de \ce{CH4}, sans doublet non liant), ni linéaire.
\item \textbf{Réponse : b (secondaire).} L'atome d'azote est lié à \textbf{deux} groupes alkyles (un méthyle et un éthyle) : c'est une amine \cle{secondaire}, de formule générale $R\text{-}NH\text{-}R'$. Ce n'est pas un ion (la molécule est neutre).
\item \textbf{Réponses : b et d.} \ce{(CH3)4N+} est un ion \cle{ammonium quaternaire} : l'azote porte quatre groupes alkyles et une charge positive. Il est obtenu par alkylation totale de l'ammoniac (quatre réactions d'Hofmann successives). Ce n'est pas une amine tertiaire (qui est neutre, $R_3N$) ni une base : il ne possède plus de doublet non liant disponible pour capter un proton.
\item \textbf{Réponses : b et c.} Une amine primaire établit des liaisons hydrogène intermoléculaires $N-H\cdots N$ : sa température d'ébullition est donc \textbf{supérieure} à celle de l'alcane de masse voisine (pas de liaisons H dans l'alcane). Mais la liaison $O-H\cdots O$ des alcools est plus forte que $N-H\cdots N$ (l'oxygène est plus électronégatif) : l'ébullition de l'amine reste \textbf{inférieure} à celle de l'alcool correspondant. Ordre : alcane $<$ amine $<$ alcool.
\item \textbf{Réponse : a.} Le caractère \cle{nucléophile} (« qui aime les noyaux », c'est-à-dire les sites pauvres en électrons) provient du \textbf{doublet non liant} de l'azote, qui peut attaquer un site électrophile. L'effet donneur des groupes alkyles renforce la basicité mais n'est pas la cause première de la nucléophilie.
\end{enumerate}
\end{corrige}

\begin{corrige}[Exercice 2 — Vrai ou faux justifié]
\begin{enumerate}
\item \textbf{Vrai.} Une liaison hydrogène intermoléculaire exige un atome d'hydrogène porté par l'azote (liaison \ce{N-H}). Une amine tertiaire $R_3N$ ne possède \textbf{aucune} liaison \ce{N-H} : ses molécules ne peuvent pas s'associer entre elles par liaisons hydrogène. (En revanche, elle peut former des liaisons hydrogène avec l'eau, grâce à son doublet non liant accepteur.)
\item \textbf{Faux.} L'aniline \ce{C6H5-NH2} possède un azote lié à \textbf{un seul} groupe carboné (le phényle) et à deux hydrogènes : c'est une amine \cle{primaire}. La nature aromatique du groupe ne change pas le critère de classement, qui porte uniquement sur le nombre de groupes carbonés liés à l'azote.
\item \textbf{Faux.} L'éthylamine est une \textbf{base} : sa solution aqueuse a un pH $> 7$. Le BBT, jaune en milieu acide et bleu en milieu basique, vire donc au \textbf{bleu}, pas au jaune.
\item \textbf{Vrai.} Dans le groupe carbonyle \ce{C=O}, l'oxygène, plus électronégatif, attire les électrons de la double liaison : le carbone porte une charge partielle $\delta^+$. De plus, le chlore du chlorure d'acyle renforce ce déficit électronique. Ce carbone est donc le site \cle{électrophile} attaqué par le doublet de l'azote de l'amine.
\item \textbf{Vrai.} L'amine tertiaire $R_3N$ possède encore son doublet non liant : elle attaque le carbone du dérivé halogéné $R'\!-\!X$ et forme un quatrième lien $N\!-\!C$. L'azote porte alors quatre groupes alkyles et une charge positive : c'est l'ion \cle{ammonium quaternaire} $[R_3N-R']^+$, associé à l'ion halogénure $X^-$.
\end{enumerate}
\end{corrige}

\begin{corrige}[Exercice 3 — Texte à trous]
La molécule d'ammoniac a une géométrie \textbf{pyramidale} et possède sur l'atome d'azote un \textbf{doublet non liant}, à l'origine de ses propriétés chimiques. Les amines sont des bases : elles captent un proton, ce qui traduit leur caractère \textbf{basique}. Elles sont aussi des réactifs \textbf{nucléophiles} : elles attaquent les sites \textbf{électrophiles} des molécules, comme l'atome de carbone porteur d'un halogène. La réaction d'une amine avec un dérivé halogéné est une réaction d'\textbf{alkylation} ; répétée, elle conduit finalement à un ion \textbf{ammonium quaternaire}. Les amines primaires et secondaires établissent entre elles des liaisons \textbf{hydrogène}, ce qui explique leurs températures d'ébullition relativement élevées.
\end{corrige}

\begin{corrige}[Exercice 4 — Définitions]
\begin{enumerate}
\item \textbf{Amine primaire} : composé organique dérivé de l'ammoniac dans lequel l'atome d'azote est lié à \textbf{un seul} groupe carboné ; formule générale $R-\mathrm{NH}_2$. Exemple : \ce{CH3-CH2-NH2} (éthylamine).

\textbf{Amine secondaire} : l'azote est lié à \textbf{deux} groupes carbonés ; formule générale $R-\mathrm{NH}-R'$. Exemple : \ce{CH3-NH-CH3} (diméthylamine).

\textbf{Amine tertiaire} : l'azote est lié à \textbf{trois} groupes carbonés ; formule générale $R-\mathrm{N}(R')-R''$. Exemple : \ce{(CH3)3N} (triméthylamine).
\item \textbf{Ion ammonium quaternaire} : ion positif dans lequel l'atome d'azote est lié à \textbf{quatre} groupes carbonés ; formule générale $[R_4N]^+$. Exemple : \ce{[(CH3)4N]+} (ion tétraméthylammonium). Il ne possède plus de doublet non liant : il n'est ni basique ni nucléophile.
\item \textbf{Réaction d'Hofmann} : réaction d'\cle{alkylation} d'une amine (ou de l'ammoniac) par un dérivé halogéné \ce{R-X}, au cours de laquelle le doublet non liant de l'azote attaque le carbone porteur de l'halogène et forme une nouvelle liaison \ce{C-N}. Exemple :
\[ \ce{CH3-NH2 + C2H5-I -> CH3-NH-C2H5 + HI} \]
\item \textbf{Caractère nucléophile} : propriété d'une espèce chimique riche en électrons (ici grâce au doublet non liant de l'azote) qui attaque les sites déficitaires en électrons (sites électrophiles) d'autres molécules pour former une nouvelle liaison covalente.
\end{enumerate}
\end{corrige}

\courssub{Je m'entraîne}

\begin{corrige}[Exercice 5 — Classes et formules]
\begin{enumerate}
\item Formules semi-développées :
\begin{center}
\chemfig{CH_3-CH_2-NH_2} \qquad \chemfig{CH_3-NH-CH_3} \qquad \chemfig{CH_3-N(-[2]CH_3)-CH_2-CH_3} \qquad \chemfig{CH_3-CH(-[2]NH_2)-CH_3} \qquad \chemfig{C_6H_5-NH_2}
\end{center}
\item Classes (on compte le nombre de groupes carbonés liés à l'azote) :
\begin{center}
\begin{tabularx}{\linewidth}{|l|X|l|}
\hline
\rowcolor{popblueL} \textbf{Amine} & \textbf{Justification} & \textbf{Classe} \\
\hline
\ce{CH3-CH2-NH2} & 1 groupe carboné (éthyle) sur N & primaire \\
\hline
\ce{(CH3)2NH} & 2 groupes carbonés (2 méthyles) sur N & secondaire \\
\hline
\ce{CH3-N(CH3)-CH2-CH3} & 3 groupes carbonés sur N & tertiaire \\
\hline
\ce{CH3-CH(NH2)-CH3} & 1 groupe carboné (isopropyle) sur N & primaire \\
\hline
\ce{C6H5-NH2} & 1 groupe carboné (phényle) sur N & primaire \\
\hline
\end{tabularx}
\end{center}
\item \ce{CH3-CH(NH2)-CH3} et \ce{CH3-CH2-CH2-NH2} (propan-1-amine) sont toutes deux de formule brute \ce{C3H9N} : ce sont des isomères de \textbf{position} (le groupe \ce{-NH2} n'est pas porté par le même carbone).

Tous les isomères de \ce{C3H9N} :
\begin{itemize}
\item \ce{CH3-CH2-CH2-NH2} : propan-1-amine, \textbf{primaire} ;
\item \ce{CH3-CH(NH2)-CH3} : propan-2-amine, \textbf{primaire} ;
\item \ce{CH3-CH2-NH-CH3} : N-méthyléthanamine, \textbf{secondaire} ;
\item \ce{N(CH3)3} : N,N-diméthylméthanamine (triméthylamine), \textbf{tertiaire}.
\end{itemize}
Soit \textbf{quatre} isomères au total.
\end{enumerate}
\end{corrige}

\begin{corrige}[Exercice 6 — Nomenclature]
\begin{enumerate}
\item Noms systématiques :
\begin{itemize}
\item \ce{CH3-CH2-CH2-NH2} : \textbf{propan-1-amine} ;
\item \ce{CH3-NH-CH2-CH3} : \textbf{N-méthyléthanamine} (amine secondaire : le plus long groupe carboné donne le nom de l'amine principale, l'autre est cité en préfixe N-alkyle) ;
\item \ce{(CH3)3N} : \textbf{N,N-diméthylméthanamine} ;
\item \ce{CH3-CH2-CH(NH2)-CH3} : \textbf{butan-2-amine} ;
\item \ce{C6H5-NH2} : \textbf{phénylamine} (nom courant : aniline).
\end{itemize}
\item Formules semi-développées :
\begin{itemize}
\item éthanamine : \ce{CH3-CH2-NH2} ;
\item N-méthylpropan-1-amine : \ce{CH3-CH2-CH2-NH-CH3} ;
\item N,N-diméthylméthanamine : \ce{(CH3)3N}, soit \chemfig{CH_3-N(-[2]CH_3)-CH_3} ;
\item butan-2-amine : \ce{CH3-CH2-CH(NH2)-CH3} ;
\item phénylamine (aniline) : \ce{C6H5-NH2}.
\end{itemize}
\end{enumerate}
\end{corrige}

\begin{corrige}[Exercice 7 — Basicité et équations-bilans]
\begin{enumerate}
\item L'éthylamine capte un proton de l'eau grâce au doublet non liant de l'azote :
\[ \ce{C2H5-NH2 + H2O <=> C2H5-NH3+ + HO-} \]
\item Les deux couples acide/base mis en jeu :
\begin{itemize}
\item \ce{C2H5NH3+}/\ce{C2H5NH2} (ion éthylammonium / éthylamine) ;
\item \ce{H2O}/\ce{HO-} (l'eau joue ici le rôle d'acide).
\end{itemize}
\item La réaction produit des ions hydroxyde \ce{HO-} : la solution est \textbf{basique}, d'où pH $> 7$. L'hélianthine, rouge en milieu acide et jaune-orange en milieu basique (zone de virage 3,1 -- 4,4), prend donc la teinte \textbf{jaune-orange} dans cette solution.
\item L'acide chlorhydrique apporte des ions \ce{H3O+} qui réagissent totalement avec l'amine :
\[ \ce{C2H5-NH2 + H3O+ -> C2H5-NH3+ + H2O} \]
Le produit ionique formé est le \textbf{chlorure d'éthylammonium} \ce{(C2H5NH3+, Cl-)}, sel soluble dans l'eau (c'est le principe des sels d'ammonium utilisés en pharmacie pour rendre les amines médicamenteuses solubles).
\end{enumerate}
\end{corrige}

\begin{corrige}[Exercice 8 — Équations de réactions caractéristiques]
\begin{enumerate}
\item Méthylamine + iodoéthane (réaction d'Hofmann) :
\[ \ce{CH3-NH2 + C2H5-I -> CH3-NH-C2H5 + HI} \]
Produit : \textbf{N-méthyléthanamine}, amine \textbf{secondaire}.
\item Diéthylamine + bromométhane :
\[ \ce{(C2H5)2NH + CH3-Br -> (C2H5)2N-CH3 + HBr} \]
Produit : \textbf{N,N-diéthylméthanamine}, amine \textbf{tertiaire}.
\item Triméthylamine + iodométhane : l'amine tertiaire n'a plus d'hydrogène sur l'azote ; elle donne directement l'ion ammonium quaternaire :
\[ \ce{(CH3)3N + CH3-I -> [(CH3)4N]+ + I-} \]
Produit : \textbf{iodure de tétraméthylammonium}, \textbf{ion ammonium quaternaire}.
\item Éthylamine + chlorure d'éthanoyle (préparation d'un amide) :
\[ \ce{CH3-COCl + C2H5-NH2 -> CH3-CO-NH-C2H5 + HCl} \]
Produit : \textbf{N-éthyléthanamide} (N-éthylacétamide), \textbf{amide substitué}.
\item Ammoniac + chlorure d'éthanoyle :
\[ \ce{CH3-COCl + NH3 -> CH3-CO-NH2 + HCl} \]
Produit : \textbf{éthanamide} (acétamide), \textbf{amide non substitué}.
\end{enumerate}
\textbf{Remarque :} en pratique, l'acide \ce{HX} formé est neutralisé par un excès d'amine ou d'ammoniac (formation du sel d'ammonium correspondant).
\end{corrige}

\courssub{Je me prépare au Bac}

\begin{corrige}[Exercice 9 — Type Bac 1 : structure, classes et solubilité des amines]
\textbf{Barème indicatif (sur 10) :} question 1 : 2,5 pts ; question 2 : 2 pts ; question 3 : 1,5 pt ; question 4 : 2 pts ; question 5 : 2 pts.

\begin{enumerate}
\item Formules semi-développées et classes :
\begin{center}
\chemfig{CH_3-C(-[2]CH_3)(-[6]CH_3)-NH_2} \qquad \chemfig{CH_3-NH-CH_2-CH_3} \qquad \chemfig{C_6H_5-NH_2}
\end{center}
\begin{center}
\chemfig{CH_3-CH_2-N(-[2]CH_2-CH_3)-CH_2-CH_3} \qquad \chemfig{CH_3-CH_2-CH_2-CH_2-NH_2}
\end{center}
\begin{center}
\begin{tabularx}{\linewidth}{|l|X|l|}
\hline
\rowcolor{popblueL} \textbf{Composé} & \textbf{Groupes carbonés sur N} & \textbf{Classe} \\
\hline
A \ce{(CH3)3C-NH2} & 1 (tert-butyle) & primaire \\
\hline
B \ce{CH3-NH-C2H5} & 2 (méthyle + éthyle) & secondaire \\
\hline
C \ce{C6H5-NH2} & 1 (phényle) & primaire \\
\hline
D \ce{(C2H5)3N} & 3 (3 éthyles) & tertiaire \\
\hline
E \ce{CH3-CH2-CH2-CH2-NH2} & 1 (butyle) & primaire \\
\hline
\end{tabularx}
\end{center}
\item Noms systématiques :
\begin{itemize}
\item A : \textbf{2-méthylpropan-2-amine} ;
\item B : \textbf{N-méthyléthanamine} ;
\item C : \textbf{phénylamine} (nom courant : \textbf{aniline}) ;
\item D : \textbf{N,N-diéthyléthanamine} (triéthylamine) ;
\item E : \textbf{butan-1-amine}.
\end{itemize}
\item Formules brutes :
\begin{itemize}
\item A : \ce{C4H11N} ($4 \times 12 + 11 \times 1 + 14 = \SI{73}{g.mol^{-1}}$) ;
\item E : \ce{C4H11N} ($M = \SI{73}{g.mol^{-1}}$).
\end{itemize}
A et E ont la \textbf{même formule brute} \ce{C4H11N} mais des formules semi-développées différentes : ce sont des \textbf{isomères}. Comme leurs chaînes carbonées sont différentes (ramifiée pour A, linéaire pour E), il s'agit d'une \textbf{isomérie de chaîne}.
\item E (butan-1-amine, primaire) possède des liaisons \ce{N-H} : elle peut à la fois \textbf{donner} des liaisons hydrogène aux molécules d'eau (via ses H) et en \textbf{accepter} (via son doublet non liant). D (triéthylamine, tertiaire) n'a \textbf{aucune} liaison \ce{N-H} : elle ne peut qu'accepter des liaisons hydrogène de l'eau. De plus, ses trois groupes éthyle volumineux créent un encombrement hydrophobe autour de l'azote. E est donc nettement plus soluble dans l'eau que D.
\item Masses molaires : butane $M = \SI{58}{g.mol^{-1}}$ ; E : $M = \SI{73}{g.mol^{-1}}$ ; butan-1-ol : $M = \SI{74}{g.mol^{-1}}$.

Le \textbf{butane} n'établit aucune liaison hydrogène : seules de faibles interactions de Van der Waals existent entre ses molécules $\Rightarrow$ ébullition la plus \textbf{faible} ($\approx \SI{-0,5}{\celsius}$).

L'amine \textbf{E} établit des liaisons hydrogène \ce{N-H...N}, d'intensité moyenne $\Rightarrow$ ébullition intermédiaire ($\approx \SI{78}{\celsius}$).

Le \textbf{butan-1-ol} établit des liaisons hydrogène \ce{O-H...O}, plus fortes car l'oxygène est plus électronégatif que l'azote $\Rightarrow$ ébullition la plus \textbf{élevée} ($\approx \SI{118}{\celsius}$).

Conclusion : $\theta_{\text{éb}}(\text{butane}) < \theta_{\text{éb}}(\text{E}) < \theta_{\text{éb}}(\text{butan-1-ol})$.
\end{enumerate}
\textbf{Points de vigilance :} ne pas confondre la classe de l'amine avec le nombre de carbones de la molécule ; pour A, l'amine est primaire même si le carbone porteur est tertiaire ; justifier toute comparaison d'ébullition par la nature des liaisons intermoléculaires, pas seulement par la masse molaire.
\end{corrige}

\begin{corrige}[Exercice 10 — Type Bac 2 : caractère basique de la méthylamine]
\textbf{Barème indicatif (sur 12) :} question 1 : 2 pts ; question 2 : 3 pts ; question 3 : 2 pts ; question 4 : 2 pts ; question 5 : 1,5 pt ; question 6 : 1,5 pt.

\begin{enumerate}
\item Équation-bilan :
\[ \ce{CH3-NH2 + H2O <=> CH3-NH3+ + HO-} \]
Couples mis en jeu : \ce{CH3NH3+}/\ce{CH3NH2} et \ce{H2O}/\ce{HO-}.
\item Concentration en ions hydroxyde :
\[ [\ce{H3O+}] = 10^{-\text{pH}} = 10^{-11{,}3} = 5{,}0 \times 10^{-12}\ \si{mol.L^{-1}} \]
\[ [\ce{HO-}] = \frac{K_e}{[\ce{H3O+}]} = \frac{10^{-14}}{5{,}0 \times 10^{-12}} = 2{,}0 \times 10^{-3}\ \si{mol.L^{-1}} \]
Vérification : $\text{pOH} = 14 - 11{,}3 = 2{,}7$, d'où $[\ce{HO-}] = 10^{-2{,}7} \approx 2{,}0 \times 10^{-3}\ \si{mol.L^{-1}}$. Cohérent.

Pour une base \textbf{forte} à $C = \SI{0,10}{mol.L^{-1}}$, on aurait $[\ce{HO-}] = C = \SI{0,10}{mol.L^{-1}}$, soit $\text{pH} = 13$. Or $[\ce{HO-}] = 2{,}0 \times 10^{-3}\ \si{mol.L^{-1}} \ll C$ : la réaction avec l'eau est \textbf{partielle} (taux d'avancement $\tau \approx \frac{2{,}0 \times 10^{-3}}{0{,}10} = 2\ \%$). De plus, le $\text{p}K_a$ du couple \ce{CH3NH3+}/\ce{CH3NH2} vaut 10,7 (fini, non nul). La méthylamine est donc une \textbf{base faible}.
\item $\text{pH} = 11{,}3$ :
\begin{itemize}
\item \textbf{BBT} (zone de virage 6,0 -- 7,6) : pH $> 7{,}6$, le BBT prend sa teinte basique : \textbf{bleu} ;
\item \textbf{phénolphtaléine} (zone de virage 8,2 -- 10,0) : pH $> 10{,}0$, elle prend sa teinte basique : \textbf{rose} (fuchsia).
\end{itemize}
\item Les ions \ce{Cu^2+} réagissent avec les ions hydroxyde présents en solution : on observe la formation d'un \textbf{précipité bleu} d'hydroxyde de cuivre(II) :
\[ \ce{Cu^2+ + 2HO- -> Cu(OH)2 (s)} \]
(C'est un test caractéristique du caractère basique des solutions d'amines.)
\item $\text{p}K_a(\ce{CH3NH3+}/\ce{CH3NH2}) = 10{,}7 > \text{p}K_a(\ce{NH4+}/\ce{NH3}) = 9{,}2$ : plus le $\text{p}K_a$ du couple est élevé, plus la base est forte. La méthylamine est donc \textbf{plus basique que l'ammoniac}.

Justification : le groupe méthyle \ce{CH3} est \textbf{donneur d'électrons} (effet inductif donneur $+I$) : il enrichit l'atome d'azote en électrons, ce qui rend le doublet non liant \textbf{plus disponible} pour capter un proton et stabilise le cation \ce{CH3NH3+} formé.
\item Quantité de méthylamine nécessaire :
\[ n = C \times V = 0{,}10 \times 0{,}500 = 5{,}0 \times 10^{-2}\ \si{mol} \]
Masse à dissoudre :
\[ m = n \times M = 5{,}0 \times 10^{-2} \times 31 = \SI{1,55}{g} \]
Il faut donc peser environ $\boxed{\SI{1,6}{g}}$ de méthylamine et compléter à \SI{500}{mL} avec de l'eau distillée.
\end{enumerate}
\textbf{Points de vigilance :} ne pas conclure « base forte » à partir du seul pH élevé : comparer $[\ce{HO-}]$ à $C$ ; bien écrire la réaction avec une double flèche \ce{<=>} (équilibre) ; dans le calcul de masse, convertir le volume en litres.
\end{corrige}

\begin{corrige}[Exercice 11 — Type Bac 3 : réactions d'Hofmann et synthèse d'un médicament]
\textbf{Barème indicatif (sur 14) :} Partie A : question 1 : 4 pts ; question 2 : 2 pts ; question 3 : 1,5 pt ; question 4 : 1 pt. Partie B : question 5 : 2 pts ; question 6 : 1,5 pt ; question 7 : 1 pt ; question 8 : 1 pt.

\textbf{Partie A : alkylation de l'ammoniac.}
\begin{enumerate}
\item Les quatre réactions successives (réactions d'Hofmann) :
\begin{align*}
\ce{NH3 + CH3-I &-> CH3-NH2 + HI} \\
\ce{CH3-NH2 + CH3-I &-> (CH3)2NH + HI} \\
\ce{(CH3)2NH + CH3-I &-> (CH3)3N + HI} \\
\ce{(CH3)3N + CH3-I &-> [(CH3)4N]+ + I-}
\end{align*}
(En pratique, l'acide \ce{HI} formé est neutralisé par l'ammoniac ou l'amine en excès, donnant des sels d'ammonium.)
\item Produits obtenus :
\begin{center}
\begin{tabularx}{\linewidth}{|l|X|l|}
\hline
\rowcolor{popblueL} \textbf{Produit} & \textbf{Nom} & \textbf{Classe / nature} \\
\hline
\ce{CH3-NH2} & méthanamine (méthylamine) & amine primaire \\
\hline
\ce{(CH3)2NH} & N-méthylméthanamine (diméthylamine) & amine secondaire \\
\hline
\ce{(CH3)3N} & N,N-diméthylméthanamine (triméthylamine) & amine tertiaire \\
\hline
\ce{[(CH3)4N]+} & ion tétraméthylammonium & ion ammonium quaternaire \\
\hline
\end{tabularx}
\end{center}
\item Dans l'iodométhane \ce{CH3-I}, l'iode est plus électronégatif que le carbone : la liaison \ce{C-I} est polarisée, le carbone porte une charge partielle $\delta^+$ et constitue un site \textbf{électrophile}. Le doublet non liant de l'azote de \ce{NH3} (site \textbf{nucléophile}) attaque ce carbone, tandis que l'iode part avec le doublet de la liaison (groupe partant \ce{I-}) : c'est une \textbf{substitution nucléophile}. Une nouvelle liaison \ce{C-N} se forme.
\item Pour obtenir majoritairement la méthylamine (produit de la première étape), il faut mettre \textbf{l'ammoniac en grand excès} : statistiquement, une molécule d'iodométhane rencontre alors beaucoup plus souvent une molécule de \ce{NH3} qu'une molécule de méthylamine déjà formée, ce qui limite les alkylations ultérieures.

\setcounter{enumi}{4}
\item \textbf{Partie B.} Équation-bilan de la formation du paracétamol :
\[ \ce{HO-C6H4-NH2 + CH3-COCl -> HO-C6H4-NH-CO-CH3 + HCl} \]
Le groupe amine du 4-aminophénol est acylé ; le groupe hydroxyle \ce{-OH} n'est pas touché dans ces conditions.
\item \textbf{Chlorure d'éthanoyle électrophile} : le carbone du groupe carbonyle est lié à deux atomes très électronégatifs (O et Cl) qui attirent les électrons ; il porte une charge partielle $\delta^+$ importante : c'est un site électrophile.

\textbf{Amine nucléophile} : l'azote du groupe \ce{-NH2} possède un doublet non liant disponible qui attaque ce carbone déficitaire en électrons.
\item On utilise un excès d'amine pour \textbf{neutraliser l'acide chlorhydrique} \ce{HCl} formé au cours de la réaction (l'amine donne alors un chlorure d'alkylammonium). Sans excès de base, \ce{HCl} protonerait l'amine de départ, la transformant en ion ammonium \ce{R-NH3+} privé de doublet non liant, donc \textbf{non nucléophile} : la réaction s'arrêterait.
\item \textbf{Non.} Une amine tertiaire \ce{(CH3)3N} peut bien attaquer le chlorure d'acyle et former un intermédiaire ionique acylammonium \ce{[CH3-CO-N(CH3)3]+ Cl-}, mais elle ne possède \textbf{aucun atome d'hydrogène sur l'azote} : aucun proton ne peut être éliminé pour stabiliser le produit par formation d'une liaison \ce{C-N} neutre d'amide. L'intermédiaire se décompose et aucun \textbf{amide stable} ne se forme. Seules les amines primaires et secondaires (et l'ammoniac) conduisent à des amides.
\end{enumerate}
\textbf{Points de vigilance :} bien distinguer la quatrième étape (pas de \ce{HI} formé, mais un sel \ce{[(CH3)4N]+ I-}) ; ne pas oublier \ce{HCl} comme produit de l'acylation ; pour la question 8, l'argument décisif est l'absence de proton sur l'azote de l'amine tertiaire.
\end{corrige}

\newpage

\coursec{Fiche bilan}

\begin{popfigure}
\begin{center}
\begin{tikzpicture}[
  mindmap, concept/.append style={minimum size=1.6cm},
  level 1/.append style={level distance=3.6cm, sibling angle=60},
  every node/.append style={font=\small},
  root concept/.append style={font=\bfseries\small, fill=popblue!85, text=white},
  level 1 concept/.append style={font=\footnotesize, text=popdark}
]
\node[concept, root concept]{Les amines}
  child[concept color=popgreen!55]{ node[concept]{Structure et classes\\ \footnotesize R--N pyramidale\\ \footnotesize primaire, secondaire,\\ \footnotesize tertiaire, ammonium\\ \footnotesize quaternaire} }
  child[concept color=poporange!55]{ node[concept]{Nomenclature\\ \footnotesize alc\`ene $\to$ alcanamine\\ \footnotesize aniline (ph\'enylamine)} }
  child[concept color=poppurple!55]{ node[concept]{Propri\'et\'es physiques\\ \footnotesize liaisons hydrog\`ene\\ \footnotesize solubilit\'e, odeur} }
  child[concept color=popgold!65]{ node[concept]{Caract\`ere basique\\ \footnotesize $R\text{-}NH_2 + \ce{H2O} \rightleftharpoons R\text{-}NH_3^+ + \ce{HO-}$\\ \footnotesize action sur \ce{Cu^2+}} }
  child[concept color=poppink!55]{ node[concept]{Caract\`ere nucl\'eophile\\ \footnotesize alkylation (Hofmann)\\ \footnotesize jusqu'\`a $R_4N^+$} }
  child[concept color=popblue!40]{ node[concept]{Avec chlorure d'acyle\\ \footnotesize pr\'eparation des amides\\ \footnotesize $R\text{-}COCl + 2 R'\text{-}NH_2$} };
\end{tikzpicture}
\end{center}
\legende{Carte mentale de la le\c{c}on : les six comp\'etences cl\'es autour des amines.}
\end{popfigure}

\begin{bilan}[Fiche bilan --- Les amines]
\textbf{D\'efinitions cl\'es}
\begin{itemize}
  \item \cle{Amine} : compos\'e organique d\'eriv\'e de l'ammoniac \ce{NH3} par remplacement de 1, 2 ou 3 atomes d'hydrog\`ene par des groupes alkyles \ce{R} ou aryles. L'atome d'azote garde sa g\'eom\'etrie \cle{pyramidale} et son \cle{doublet non liant} : c'est lui qui explique toute la chimie des amines (basicit\'e + nucl\'eophilie).
  \item \cle{Classes} : amine primaire \ce{R-NH2}, secondaire \ce{R2NH}, tertiaire \ce{R3N} ; \cle{ion ammonium quaternaire} \ce{R4N+} (associ\'e \`a un anion \ce{X-}).
  \item \cle{Aniline} (ph\'enylamine) : \ce{C6H5-NH2}, amine aromatique de r\'ef\'erence.
  \item \cle{Couple acide/base} : \ce{R-NH3+}/\ce{R-NH2} ; les amines sont des bases plus fortes que \ce{NH3} (effet donneur des groupes alkyles), sauf l'aniline, base plus faible (doublet d\'elocalis\'e dans le cycle).
\end{itemize}

\textbf{\'Equations-types \`a conna\^itre par c\oe ur}
\begin{itemize}
  \item Action de l'eau (\'equilibre, solution basique) :
  \[ \ce{\text{R-NH2} + H2O <=> \text{R-NH3}+ + HO-} \]
  \item Pr\'ecipitation puis redissolution des ions cuivre(II) :
  \[ \ce{Cu^2+ + 2 HO- -> Cu(OH)2} \qquad \ce{Cu(OH)2 + 4 \text{R-NH2} -> [Cu(\text{R-NH2})4]^2+ + 2 HO-} \]
  \item R\'eactions d'Hofmann successives (alkylation, avec \ce{NaOH} pour lib\'erer l'amine) :
  \[ \ce{\text{R-NH2} + \text{R$'$-X} -> \text{R-NH-R$'$} + \text{HX}} \quad\text{puis}\quad \ce{\text{R2N-R$'$} + \text{R$'$-X} -> \text{R2N(R$'$)2} + \text{HX}} \quad\text{puis}\quad \ce{\text{R3N} + \text{R-X} -> [\text{R4N}]+ \text{X}-} \]
  \item Formation d'un amide (rapide et totale, contrairement \`a l'est\'erification) :
  \[ \ce{\text{R-COCl} + \text{R$'$-NH2} -> \text{R-CO-NH-R$'$} + HCl} \qquad\text{(2 mol d'amine : 1 pour l'amide, 1 pour pi\'eger \ce{HCl})} \]
\end{itemize}

\textbf{Formules et classes (rep\`eres)}
\begin{center}
\begin{tabularx}{\linewidth}{|l|X|X|}
\hline
\rowcolor{popblueL} \textbf{Classe} & \textbf{Formule g\'en\'erale} & \textbf{Exemple} \\
\hline
Primaire & \ce{\text{R}-NH2} & \'ethanamine \ce{CH3-CH2-NH2} \\
\hline
Secondaire & \ce{\text{R}-NH-\text{R$'$}} & N-m\'ethyl\'ethanamine \ce{CH3-CH2-NH-CH3} \\
\hline
Tertiaire & \ce{N(\text{R})3} & N,N-dim\'ethylm\'ethanamine \ce{N(CH3)3} \\
\hline
Ammonium quaternaire & \ce{[\text{R}4N]+ \text{X}-} & bromure de t\'etram\'ethylammonium \ce{[N(CH3)4]+ Br-} \\
\hline
\end{tabularx}
\end{center}

\textbf{M\'ethodes express}
\begin{itemize}
  \item \cle{Classer une amine} : compter le nombre de groupes carbon\'es li\'es \`a l'azote (1 $\to$ primaire, 2 $\to$ secondaire, 3 $\to$ tertiaire) --- et non le nombre de carbones de la cha\^ine !
  \item \cle{Nommer} : cha\^ine principale + suffixe \emph{-amine} avec indice de position (\ce{CH3-CH(NH2)-CH3} : propan-2-amine) ; substituants sur N pr\'ec\'ed\'es de \emph{N-}.
  \item \cle{Propri\'et\'es physiques} : les amines primaires et secondaires (liaison \ce{N-H}) font des liaisons hydrog\`ene $\Rightarrow$ points d'\'ebullition plus \'elev\'es et bonne solubilit\'e dans l'eau des petites mol\'ecules ; les tertiaires ne s'auto-associent pas. Odeur caract\'eristique de poisson.
  \item \cle{Pr\'eparer un amide} : chlorure d'acyle (ou anhydride) + amine, r\'eaction exothermique et totale ; l'amine joue le r\^ole de nucl\'eophile, le carbone du groupe \ce{C=O} est le site \'electrophile.
\end{itemize}
\end{bilan}

\begin{aretenir}[Checklist avant le Bac]
\begin{itemize}
  \item $\square$ Je sais dessiner la structure pyramidale de \ce{NH3} et justifier le passage aux amines.
  \item $\square$ Je sais \'ecrire la formule semi-d\'evelopp\'ee d'une amine et identifier sa classe (primaire, secondaire, tertiaire).
  \item $\square$ Je sais nommer une amine aliphatique et reconna\^itre l'aniline \ce{C6H5-NH2}.
  \item $\square$ Je sais expliquer solubilit\'e et points d'\'ebullition par les liaisons hydrog\`ene.
  \item $\square$ Je sais \'ecrire l'\'equation de l'action de l'eau sur une amine et citer le couple \ce{R-NH3+}/\ce{R-NH2}.
  \item $\square$ Je sais comparer la basicit\'e des amines \`a celle de l'ammoniac et pr\'evoir l'action sur un indicateur color\'e (B.B.T. au bleu).
  \item $\square$ Je sais d\'ecrire le test aux ions \ce{Cu^2+} : pr\'ecipit\'e bleu de \ce{Cu(OH)2} puis solution bleu fonc\'e du complexe.
  \item $\square$ Je sais \'ecrire les r\'eactions d'Hofmann successives jusqu'\`a l'ion ammonium quaternaire \ce{[R4N]+}.
  \item $\square$ Je sais \'ecrire l'\'equation amine + chlorure d'acyle et justifier l'emploi de 2 moles d'amine.
  \item $\square$ Je sais d\'efinir le caract\`ere nucl\'eophile (doublet non liant de N) et le site \'electrophile (carbone de \ce{C=O}).
  \item $\square$ Je sais construire un tableau d'avancement pour calculer un rendement de pr\'eparation d'amide.
  \item $\square$ Je v\'erifie toujours mes \'equations : conservation des atomes, des charges, et fl\`eches adapt\'ees (\ce{<=>} pour l'\'equilibre avec l'eau, \ce{->} pour Hofmann et les amides).
\end{itemize}
\end{aretenir}

\findecours

\end{document}
